% Applying project management fundamentals to run sample projects — ICT Upper Sixth — Cours signé OnBuch+
% Official MINESEC syllabus. Compiler avec : tectonic ICT20-applying-project-management-fundamentals-to-run-sample-proje.tex
\newcommand{\DOCMATIERE}{ICT}
\newcommand{\DOCNIVEAU}{Upper Sixth}
\newcommand{\DOCMODULE}{Upper Sixth — ICT}
\newcommand{\DOCLECON}{20}
\newcommand{\DOCTITRE}{Applying project management fundamentals to run sample projects}
% OnBuch+ — préambule « cours signés » ENRICHI (compilé avec tectonic / XeLaTeX)
% Système visuel « Pop » d'OnBuch+ : fond crème (jamais blanc), bordures encre
% épaisses, ombres « dures » décalées, titres Archivo Black en capitales.
% Version MATHÉMATIQUES : graphes (pgfplots/tikz), boîtes pédagogiques
% (définition, propriété, méthode, attention, le savais-tu, exercice résolu,
% exercice, TP, bilan), page de garde et signature OnBuch+ ancrée en bas.
%
% Macros attendues dans main.tex AVANT \input{preamble} :
%   \DOCMATIERE  \DOCNIVEAU  \DOCMODULE  \DOCLECON (numéro)  \DOCTITRE
\documentclass[11pt]{article}

\usepackage{fontspec}
\usepackage[english]{babel}
\usepackage[a4paper,margin=2cm,top=3.4cm,bottom=2.8cm,headheight=1.4cm,footskip=1.3cm]{geometry}
\usepackage{amsmath,amssymb,amsfonts}
\usepackage{mathtools} % \xmapsto, \coloneqq... souvent utilisés par les modèles
\usepackage{cancel} % simplifications barrées (bilans de forces)
% \abs{z} (module, valeur absolue) et \norm{u} (norme), utilisés par les modèles
\providecommand{\abs}{}\let\abs\relax\DeclarePairedDelimiter{\abs}{\lvert}{\rvert}
\providecommand{\norm}{}\let\norm\relax\DeclarePairedDelimiter{\norm}{\lVert}{\rVert}
\usepackage[table]{xcolor} % [table] : fournit aussi \rowcolors (alternance de lignes)
\usepackage{pagecolor}
\usepackage{tikz}
\usetikzlibrary{shapes.symbols,circuits.logic.US,circuits.logic.IEC,arrows.meta,positioning,calc,shapes.geometric,shapes.misc,shapes.arrows,decorations.pathmorphing,decorations.pathreplacing,decorations.markings,patterns,shadows,mindmap,trees,fit,backgrounds,matrix,intersections,angles,quotes}
\usepackage{pgfplots}
\usepackage[version=4]{mhchem} % équations nucléaires \ce{^{235}_{92}U}
\usepackage[european,siunitx]{circuitikz} % schémas électriques
\pgfplotsset{compat=1.18}
\usepgfplotslibrary{fillbetween,groupplots,polar,statistics,dateplot} % aires entre courbes (calcul intégral)
\usepackage{enumitem}
\usepackage{fancyhdr}
% [most] charge toutes les bibliothèques tcolorbox (listings, minted, raster,
% documentation...) et gonfle inutilement la mémoire TeX sur un document long
% (~300 boîtes) : on ne charge que ce qui est réellement utilisé.
\usepackage[skins,breakable]{tcolorbox}
\usepackage{graphicx}
\usepackage{colortbl}
\usepackage{booktabs}
\usepackage{tabularx}
\usepackage{diagbox} % case d'en-tête diagonale des tableaux à double entrée
% Colonnes extensibles centrée / gauche / droite (C, L, R), souvent utilisées par les modèles
\newcolumntype{C}{>{\centering\arraybackslash}X}
\newcolumntype{L}{>{\raggedright\arraybackslash}X}
\newcolumntype{R}{>{\raggedleft\arraybackslash}X}
\usepackage{array}
\usepackage{multicol}
\usepackage{siunitx}
% parse-numbers=false : accepte aussi \SI{2,5 \times 10^{-3}}{...}
\sisetup{locale=UK,output-decimal-marker={.},group-minimum-digits=4,parse-numbers=false}
\DeclareSIUnit\ohm{\ensuremath{\symup{\Omega}}} % U+2126 absent de la police
\DeclareSIUnit\division{div} % réglages d'oscilloscope (ms/div, V/div)
\DeclareSIUnit\year{yr}\DeclareSIUnit\annum{yr}\DeclareSIUnit\Ma{Ma}\DeclareSIUnit\tr{tr} % tours (vitesse de rotation, ex. \SI{1500}{\tr\per\minute})
\usepackage{qrcode}
% \degree : commande fréquemment utilisée par les modèles pour un angle ou une
% température en dehors de \SI{}{} (non fournie par les paquets chargés).
\providecommand{\degree}{^{\circ}}
% \qed (fin de démonstration), utilisé par les modèles sans amsthm
\providecommand{\qed}{\hfill$\square$}
% \barre : double barre des valeurs interdites dans un tableau de variations (tkz-tab)
\providecommand{\barre}{\ensuremath{\|}}
% environnement proof (amsthm non chargé) utilisé par les modèles
\ifdefined\proof\else\newenvironment{proof}[1][Proof]{\par\noindent\textit{#1.}\ }{\qed\par}\fi
\usepackage{lastpage}
\usepackage{hyperref}
\hypersetup{hidelinks}

% ============================================================
% Palette « Pop » OnBuch+ — mêmes hex que la classe P (plus_ui.dart)
% ============================================================
\definecolor{poporange}{HTML}{F59321}
\definecolor{popdark}{HTML}{E07A0C}
\definecolor{popcream}{HTML}{FFF6E8}
\definecolor{popink}{HTML}{1C1714}
\definecolor{poppurple}{HTML}{7A5AE0}
\definecolor{popgreen}{HTML}{2E9E6A}
\definecolor{poppink}{HTML}{E0457B}
\definecolor{popblue}{HTML}{2F7DE1}
\definecolor{popgold}{HTML}{C98A12}
\definecolor{popmuted}{HTML}{8A7F76}
\definecolor{popline}{HTML}{EADFD2}
\definecolor{popgray}{HTML}{8A7F76}
\colorlet{popgrey}{popgray}
% Alias tolérants (le modèle nomme parfois les couleurs différemment)
\colorlet{popwhite}{white}
\colorlet{popblack}{popink}
\colorlet{popred}{poppink}
\colorlet{popyellow}{popgold}
% Teintes claires pour fonds de tableaux / zones
\colorlet{popblueL}{popblue!10!white}
\colorlet{poporangeL}{poporange!14!white}
\colorlet{popgreenL}{popgreen!12!white}
\colorlet{poppurpleL}{poppurple!12!white}
\colorlet{poppinkL}{poppink!12!white}
\colorlet{popgoldL}{popgold!14!white}

\pagecolor{popcream}

% ============================================================
% Typographie
% ============================================================
\defaultfontfeatures{Ligatures=TeX}
\setmainfont{PlusJakartaSans-Medium.ttf}[Path=../../../fonts/,BoldFont=PlusJakartaSans-Bold.ttf]
% Maths en sans-serif (Fira Math) avec chiffres et lettres droites de Plus
% Jakarta, pour que les formules se fondent dans le texte.
\usepackage[math-style=ISO,bold-style=ISO]{unicode-math}
\setmathfont{FiraMath-Regular.otf}[Path=../../../fonts/]
\setmathfont{PlusJakartaSans-Medium.ttf}[Path=../../../fonts/,range={up/num,up/{Latin,latin},\mathup}]
% (icomma retiré : décimales avec point en anglais)
% Caractères mathématiques Unicode absents des polices de texte (Plus Jakarta,
% Archivo Black) : rendus par leur équivalent mathématique, en texte comme en
% formule — sinon ils s'affichent en carré vide (ex. « Arithmétique dans ℤ »).
\usepackage{newunicodechar}
\newunicodechar{ℝ}{\ensuremath{\mathbb{R}}}
\newunicodechar{ℤ}{\ensuremath{\mathbb{Z}}}
\newunicodechar{ℂ}{\ensuremath{\mathbb{C}}}
\newunicodechar{ℕ}{\ensuremath{\mathbb{N}}}
\newunicodechar{ℚ}{\ensuremath{\mathbb{Q}}}
\newunicodechar{𝔻}{\ensuremath{\mathbb{D}}}
\newunicodechar{α}{\ensuremath{\alpha}}
\newunicodechar{β}{\ensuremath{\beta}}
\newunicodechar{γ}{\ensuremath{\gamma}}
\newunicodechar{δ}{\ensuremath{\delta}}
\newunicodechar{ε}{\ensuremath{\varepsilon}}
\newunicodechar{θ}{\ensuremath{\theta}}
\newunicodechar{λ}{\ensuremath{\lambda}}
\newunicodechar{μ}{\ensuremath{\mu}}
\newunicodechar{σ}{\ensuremath{\sigma}}
\newunicodechar{τ}{\ensuremath{\tau}}
\newunicodechar{φ}{\ensuremath{\varphi}}
\newunicodechar{ψ}{\ensuremath{\psi}}
\newunicodechar{ω}{\ensuremath{\omega}}
\newunicodechar{Σ}{\ensuremath{\Sigma}}
% majuscules produites par \MakeUppercase dans les titres (α-aminés -> Α)
\newunicodechar{Α}{\ensuremath{\alpha}}
\newunicodechar{Β}{\ensuremath{\beta}}
\newunicodechar{Γ}{\ensuremath{\Gamma}}
\newunicodechar{Δ}{\ensuremath{\Delta}}
\newunicodechar{Ε}{\ensuremath{\varepsilon}}
\newunicodechar{Θ}{\ensuremath{\Theta}}
\newunicodechar{Λ}{\ensuremath{\Lambda}}
\newunicodechar{Μ}{\ensuremath{\mu}}
\newunicodechar{Τ}{\ensuremath{\tau}}
\newunicodechar{Φ}{\ensuremath{\Phi}}
\newunicodechar{Ψ}{\ensuremath{\Psi}}
\newunicodechar{Ω}{\ensuremath{\Omega}}
\newunicodechar{↦}{\ensuremath{\mapsto}}
\newunicodechar{∈}{\ensuremath{\in}}
\newunicodechar{∉}{\ensuremath{\notin}}
\newunicodechar{∧}{\ensuremath{\wedge}}
\newunicodechar{⇔}{\ensuremath{\Leftrightarrow}}
\newunicodechar{⟺}{\ensuremath{\Longleftrightarrow}}
\newunicodechar{⇒}{\ensuremath{\Rightarrow}}
\newunicodechar{⟹}{\ensuremath{\Longrightarrow}}
\newunicodechar{∘}{\ensuremath{\circ}}
\newunicodechar{∪}{\ensuremath{\cup}}
\newunicodechar{∩}{\ensuremath{\cap}}
\newunicodechar{≡}{\ensuremath{\equiv}}
\newunicodechar{⊂}{\ensuremath{\subset}}
\newunicodechar{⊥}{\ensuremath{\perp}}
\newunicodechar{∃}{\ensuremath{\exists}}
\newunicodechar{∀}{\ensuremath{\forall}}
\newunicodechar{∛}{\ensuremath{\sqrt[3]{\,}}}
\newunicodechar{∜}{\ensuremath{\sqrt[4]{\,}}}
\newunicodechar{⃗}{\ensuremath{{}^{\to}}}
\newunicodechar{ⁿ}{\ifmmode^{n}\else\textsuperscript{\ensuremath{n}}\fi}
\newunicodechar{ˣ}{\ifmmode^{x}\else\textsuperscript{\ensuremath{x}}\fi}
\newunicodechar{ᵇ}{\ifmmode^{b}\else\textsuperscript{\ensuremath{b}}\fi}
\newunicodechar{ʳ}{\ifmmode^{r}\else\textsuperscript{\ensuremath{r}}\fi}
\newunicodechar{ᵘ}{\ifmmode^{u}\else\textsuperscript{\ensuremath{u}}\fi}
\newunicodechar{ʸ}{\ifmmode^{y}\else\textsuperscript{\ensuremath{y}}\fi}
\newunicodechar{ᵃ}{\ifmmode^{a}\else\textsuperscript{\ensuremath{a}}\fi}
\newunicodechar{ᵉ}{\ifmmode^{e}\else\textsuperscript{\ensuremath{e}}\fi}
\newunicodechar{ᵏ}{\ifmmode^{k}\else\textsuperscript{\ensuremath{k}}\fi}
\newunicodechar{ᵖ}{\ifmmode^{p}\else\textsuperscript{\ensuremath{p}}\fi}
\newunicodechar{ᴺ}{\ifmmode^{N}\else\textsuperscript{\ensuremath{N}}\fi}
\newunicodechar{ᶜ}{\ifmmode^{c}\else\textsuperscript{\ensuremath{c}}\fi}
\newunicodechar{ʰ}{\ifmmode^{h}\else\textsuperscript{\ensuremath{h}}\fi}
\newunicodechar{ᵠ}{\ifmmode^{q}\else\textsuperscript{\ensuremath{q}}\fi}
\newunicodechar{ₙ}{\ifmmode_{n}\else\textsubscript{\ensuremath{n}}\fi}
\newunicodechar{ₐ}{\ifmmode_{a}\else\textsubscript{\ensuremath{a}}\fi}
\newunicodechar{ᵢ}{\ifmmode_{i}\else\textsubscript{\ensuremath{i}}\fi}
\newunicodechar{ᵣ}{\ifmmode_{r}\else\textsubscript{\ensuremath{r}}\fi}
\newunicodechar{ₖ}{\ifmmode_{k}\else\textsubscript{\ensuremath{k}}\fi}
\newunicodechar{ᵦ}{\ifmmode_{\beta}\else\textsubscript{\ensuremath{\beta}}\fi}
\newunicodechar{ₓ}{\ifmmode_{x}\else\textsubscript{\ensuremath{x}}\fi}
\newfontfamily\popdisplay{ArchivoBlack-Regular.ttf}[Path=../../../fonts/]
\newfontfamily\poplogo{SpaceGrotesk-Bold.ttf}[Path=../../../fonts/]
\newfontfamily\popsemi{PlusJakartaSans-SemiBold.ttf}[Path=../../../fonts/]
\newfontfamily\popxbold{PlusJakartaSans-ExtraBold.ttf}[Path=../../../fonts/]
\newfontfamily\popxboldls{PlusJakartaSans-ExtraBold.ttf}[Path=../../../fonts/,LetterSpace=3.0]

\setlist[enumerate]{leftmargin=1.7em,itemsep=5pt,topsep=4pt}
\setlist[itemize]{leftmargin=1.7em,itemsep=4pt,topsep=4pt,label={\color{poporange}\textbullet}}
\setlength{\parindent}{0pt}
\setlength{\parskip}{5pt}
\renewcommand{\arraystretch}{1.25}

% pgfplots : style Pop par défaut
\pgfplotsset{
  every axis/.append style={axis line style={popink,line width=1pt},tick style={popink},
    grid style={popline,line width=0.5pt},label style={font=\small},tick label style={font=\footnotesize}},
  popcurve/.style={poporange,line width=1.8pt,no markers,smooth},
}
% Même style disponible dans un \draw TikZ simple (hors pgfplots)
\tikzset{popcurve/.style={poporange,line width=1.8pt}}
\tikzset{popgrid/.style={popline,very thin}}
\colorlet{popcurve}{poporange} % le nom du style est parfois utilisé comme couleur

% ============================================================
% Pill / squiggle
% ============================================================
\newcommand{\poppill}[3]{%
  \tikz[baseline=-0.55ex]\node[draw=popink,line width=1.2pt,rounded corners=2.6pt,%
    fill=#2,inner xsep=6.5pt,inner ysep=3pt]{%
    \popxboldls\fontsize{7}{7}\selectfont\color{#3}\MakeUppercase{#1}};%
}
\newcommand{\popsquiggle}[1]{%
  \tikz[baseline=-0.4ex]{
    \draw[#1,line width=2.6pt,line cap=round]
      plot [smooth] coordinates {(0,0.09) (0.34,0.01) (0.68,0.21) (1.02,0.01) (1.36,0.10)};
  }%
}
% Mot-clé surligné (marqueur orange sous le texte)
\newcommand{\cle}[1]{{\popxbold\color{popdark}#1}}

% ============================================================
% En-tête / pied
% ============================================================
\pagestyle{fancy}
\fancyhf{}
\renewcommand{\headrulewidth}{0pt}
\renewcommand{\footrulewidth}{0pt}
\lhead{\raisebox{-0.15ex}{\poplogo\fontsize{13}{13}\selectfont\color{popink}OnBuch\color{poporange}+}}
\rhead{\poppill{\DOCMATIERE\ \textbullet\ \DOCNIVEAU\ \textbullet\ Chapter \DOCLECON}{white}{popink}}
\lfoot{\popxbold\fontsize{7.6}{7.6}\selectfont\color{popmuted}\MakeUppercase{Signed course OnBuch+}\ \textbullet\ {\popsemi\DOCTITRE}}
\rfoot{\tikz[baseline=-0.6ex]\node[rounded corners=8.5pt,fill=popink,minimum height=17pt,minimum width=17pt,inner xsep=5pt,inner sep=0]{\popxbold\fontsize{8}{8}\selectfont\color{popcream}\thepage\,/\,\pageref*{LastPage}};}

% ============================================================
% Titres
% ============================================================
\newcommand{\chaptitre}[2]{%
  \begin{tcolorbox}[enhanced,colback=poporange,colframe=popink,boxrule=2.6pt,arc=11pt,
      drop shadow=popink,left=15pt,right=15pt,top=12pt,bottom=11pt,before skip=3pt,after skip=18pt]
    \poppill{#2}{popink}{popcream}\par\vspace{9pt}
    {\raggedright\hyphenpenalty=10000\exhyphenpenalty=10000\popdisplay\fontsize{22}{24}\selectfont\color{popcream}\MakeUppercase{#1}\par}
  \end{tcolorbox}
}
% Section numérotée : pastille orange avec le numéro + titre Archivo
\newcounter{popsec}
\newcommand{\coursec}[1]{%
  \stepcounter{popsec}\setcounter{popsub}{0}%
  \par\vspace{16pt}\noindent
  \begin{minipage}{\linewidth}
  \tikz[baseline=-0.7ex]\node[draw=popink,line width=1.6pt,fill=poporange,rounded corners=4pt,minimum size=22pt,inner sep=0]{\popdisplay\fontsize{12}{12}\selectfont\color{popcream}\thepopsec};%
  \hspace{8pt}{\popdisplay\fontsize{15}{17}\selectfont\color{popink}\MakeUppercase{#1}}\par
  \vspace{4pt}\noindent{\color{poporange}\rule{36pt}{3.6pt}}
  \end{minipage}\par\nopagebreak\vspace{8pt}
}
\newcounter{popsub}
\newcommand{\courssub}[1]{%
  \stepcounter{popsub}%
  \par\vspace{9pt}\noindent{\popxbold\fontsize{11.5}{13}\selectfont\color{popink}{\color{poporange}\thepopsec.\thepopsub}\hspace{6pt}#1}\par\nopagebreak\vspace{4pt}
}

% ============================================================
% Boîtes pédagogiques — même recette : fond blanc, bordure épaisse colorée,
% ombre dure encre, badge plein. Titre optionnel [#1].
% ============================================================
\tcbset{popbase/.style={breakable,enhanced,colback=white,boxrule=2.3pt,arc=9pt,
  shadow={3pt}{-3pt}{0pt}{popink},left=14pt,right=14pt,top=11pt,bottom=11pt,before skip=11pt,after skip=15pt,
  fontupper=\popsemi\fontsize{10.6}{14.5}\selectfont\color{popink},
  fontlower=\popsemi\fontsize{10.6}{14.5}\selectfont\color{popink}}}
% Coche dessinée (le glyphe ✓ n'existe pas dans les polices du document)
\newcommand{\popcheck}[1]{\tikz[baseline=-0.5ex]\draw[#1,line width=1.6pt,line cap=round,line join=round](-0.13,0.0)--(-0.04,-0.09)--(0.13,0.1);}
\providecommand{\coche}{\popcheck{popgreen}}
\providecommand{\resultat}[1]{\par\smallskip\noindent\fcolorbox{popgreen}{popgreenL}{\parbox{\dimexpr\linewidth-2\fboxsep-2\fboxrule\relax}{\textbf{Result.} #1}}\par\smallskip}
\newcommand{\popboxhead}[3]{\poppill{#1}{#2}{white}\ifx\relax#3\relax\else\hspace{6pt}{\popxbold\fontsize{10.6}{12}\selectfont\color{popink}#3}\fi\par\vspace{7pt}}

\newtcolorbox{aretenir}[1][]{popbase,colframe=popgreen,before upper={\popboxhead{Key points}{popgreen}{#1}}}
\newtcolorbox{exemplebox}[1][]{popbase,colframe=poporange,before upper={\popboxhead{Exemple}{poporange}{#1}}}
\newtcolorbox{definition}[1][]{popbase,colframe=popblue,before upper={\popboxhead{Definition}{popblue}{#1}}}
\newtcolorbox{propriete}[1][]{popbase,colframe=poppurple,before upper={\popboxhead{Property}{poppurple}{#1}}}
\newtcolorbox{methode}[1][]{popbase,colframe=popgold,before upper={\popboxhead{Method}{popgold}{#1}}}
\newtcolorbox{attention}[1][]{popbase,colframe=poppink,before upper={\popboxhead{Warning}{poppink}{#1}}}
\newtcolorbox{experience}[1][]{popbase,colframe=popblue,colback=popblueL,before upper={\popboxhead{Experiment}{popblue}{#1}}}
% « Le savais-tu ? » : carte violette pleine (contexte, culture, Cameroun)
\newtcolorbox{savaistu}[1][]{popbase,colback=poppurple,colframe=popink,
  fontupper=\popsemi\fontsize{10.4}{14.2}\selectfont\color{white},
  before upper={\poppill{Did you know?}{white}{poppurple}\ifx\relax#1\relax\else\hspace{6pt}{\popxbold\color{white}#1}\fi\par\vspace{7pt}}}
% Exercice résolu : énoncé en haut, \tcblower puis la solution sur fond crème
\newcounter{popexo}\newcounter{popexr}
\newtcolorbox{exoresolu}[1][]{popbase,colframe=poporange,colbacklower=popcream,
  segmentation style={popink,line width=1.2pt,dashed},
  before upper={\stepcounter{popexr}\popboxhead{Worked example \thepopexr}{poporange}{#1}},
  before lower={\poppill{Solution}{popink}{popcream}\par\vspace{6pt}}}
% Exercice (non résolu en place) — la correction est dans la partie Corrigés
\newtcolorbox{exercice}[1][]{popbase,colframe=popink,
  before upper={\stepcounter{popexo}\popboxhead{Exercise \thepopexo}{popink}{#1}}}
% Corrigé
\newtcolorbox{corrige}[1][]{popbase,colframe=popgreen,colback=popgreenL,
  before upper={\popboxhead{Solution}{popgreen}{#1}}}
% Objectifs / prérequis / bilan
\newtcolorbox{objectifs}{popbase,colframe=popink,colback=poporangeL,
  before upper={\popboxhead{Chapter objectives}{popink}{}}}
\newtcolorbox{prerequis}{popbase,colframe=popink,colback=popblueL,
  before upper={\popboxhead{Prerequisites}{popblue}{}}}
\newtcolorbox{bilan}{popbase,colframe=popink,colback=white,boxrule=2.8pt,
  before upper={\popboxhead{Fiche bilan}{poporange}{L'essentiel à connaître pour l'examen}}}
% Figure encadrée avec légende
\newtcolorbox{popfigure}[1][]{enhanced,breakable=false,colback=white,colframe=popline,boxrule=1.4pt,arc=8pt,
  left=10pt,right=10pt,top=10pt,bottom=8pt,before skip=10pt,after skip=12pt,halign=center,
  lower separated=false,halign lower=center,
  fontlower=\popsemi\fontsize{9}{11.5}\selectfont\color{popmuted},
  #1}
\newcommand{\legende}[1]{\par\vspace{6pt}{\popsemi\fontsize{9}{11.5}\selectfont\color{popmuted}#1}}

% ============================================================
% Signature OnBuch+
% ============================================================
\newcommand{\poplogoinline}[1]{{\poplogo\fontsize{#1}{#1}\selectfont\color{popink}OnBuch\color{poporange}+}}
\newcommand{\signatureonbuch}{%
  \begin{tcolorbox}[enhanced,colback=popink,colframe=popink,boxrule=2.2pt,arc=10pt,
      drop shadow=poporange,left=14pt,right=14pt,top=10pt,bottom=10pt,before skip=0pt,after skip=0pt]
    \tikz[baseline=-0.7ex]\node[circle,fill=popgreen,draw=popcream,line width=1.3pt,minimum size=22pt,inner sep=0]{\popcheck{white}};%
    \hspace{9pt}\begin{minipage}[c]{0.62\linewidth}
      {\popxbold\fontsize{12}{13}\selectfont\color{popcream}Signed\ \textbullet\ {\poplogo OnBuch\color{poporange}+}}\par\vspace{2pt}
      {\raggedright\popsemi\fontsize{8}{10.5}\selectfont\color{popline}\MakeUppercase{OnBuch+ teaching team}\\\MakeUppercase{Follows the official MINESEC syllabus}\par}
    \end{minipage}\hfill
    {\poplogo\fontsize{22}{22}\selectfont\color{popcream}OnBuch\color{poporange}+}
  \end{tcolorbox}%
}

% ============================================================
% Page de garde
% #1 titre · #2 matière · #3 niveau · #4 module · #5 n° leçon · #6 durée ·
% #7 description / objectifs (texte ou itemize)
% ============================================================
\newcommand{\pagegarde}[7]{%
  \thispagestyle{empty}
  \newgeometry{a4paper,margin=1.7cm,top=1.6cm,bottom=1.4cm}%
  \begin{tikzpicture}[remember picture,overlay]
    \fill[poporange!18] ([xshift=-3.2cm,yshift=-3.6cm]current page.north east) circle (4.6cm);
    \fill[poppurple!14] ([xshift=2.4cm,yshift=7.6cm]current page.south west) circle (3.8cm);
  \end{tikzpicture}%
  {\poplogo\fontsize{30}{30}\selectfont\color{popink}OnBuch\color{poporange}+}\hfill
  \raisebox{0.6ex}{\poppill{Signed course \textbullet\ Premium}{popink}{popcream}}\par
  \vspace{1pt}\popsquiggle{poppurple}\par
  \vspace{8pt}
  \begin{tcolorbox}[enhanced,colback=poporange,colframe=popink,boxrule=2.8pt,arc=13pt,
      drop shadow=popink,left=18pt,right=18pt,top=12pt,bottom=13pt,after skip=10pt]
    \poppill{#2 \textbullet\ #3}{popink}{popcream}\hspace{5pt}\poppill{#4}{white}{popink}\par\vspace{8pt}
    {\popdisplay\fontsize{14}{14}\selectfont\color{popink}CHAPTER #5}\par\vspace{4pt}
    {\raggedright\hyphenpenalty=10000\exhyphenpenalty=10000\popdisplay\fontsize{25}{27}\selectfont\color{popcream}\MakeUppercase{#1}\par}
    \vspace{8pt}
    {\popxbold\fontsize{9.5}{9.5}\selectfont\color{popink}\MakeUppercase{Suggested time: #6}}
  \end{tcolorbox}
  \begin{tcolorbox}[enhanced,colback=white,colframe=popink,boxrule=2.2pt,arc=10pt,
      drop shadow=popink,left=16pt,right=16pt,top=10pt,bottom=10pt,after skip=8pt]
    \poppill{In this chapter}{poppurple}{white}\par\vspace{5pt}
    \popsemi\fontsize{9.3}{12.6}\selectfont\color{popink}#7
  \end{tcolorbox}
  \noindent\begin{minipage}[t]{0.487\linewidth}
    \begin{tcolorbox}[enhanced,colback=popgreen,colframe=popink,boxrule=2.2pt,arc=10pt,
        drop shadow=popink,left=11pt,right=11pt,top=10pt,bottom=10pt,height=3.1cm,valign=center]
      \tcbox[colback=white,colframe=popink,boxrule=1.3pt,arc=5pt,boxsep=0pt,nobeforeafter,
        left=3pt,right=3pt,top=3pt,bottom=3pt,valign=center]{\qrcode[height=1.9cm]{https://play.google.com/store/apps/details?id=com.luvvix.onbuch&hl=en}}%
      \hspace{8pt}\begin{minipage}[c]{0.5\linewidth}
        {\popdisplay\fontsize{11}{12.5}\selectfont\color{white}\MakeUppercase{Download OnBuch}}\par\vspace{4pt}
        {\raggedright\popsemi\fontsize{8.4}{11}\selectfont\color{white}Free \textbullet\ Google Play\\AI tutor, past papers, results\par}
      \end{minipage}
    \end{tcolorbox}
  \end{minipage}\hfill
  \begin{minipage}[t]{0.487\linewidth}
    \begin{tcolorbox}[enhanced,colback=poppurple,colframe=popink,boxrule=2.2pt,arc=10pt,
        drop shadow=popink,left=11pt,right=11pt,top=10pt,bottom=10pt,height=3.1cm,valign=center]
      \tikz[baseline=-0.7ex]\node[circle,fill=white,draw=popink,line width=1.3pt,minimum size=1.6cm,inner sep=0]{\popdisplay\fontsize{24}{24}\selectfont\color{poppurple}+};%
      \hspace{8pt}\begin{minipage}[c]{0.6\linewidth}
        {\popdisplay\fontsize{11}{12.5}\selectfont\color{white}\MakeUppercase{Go OnBuch+}}\par\vspace{4pt}
        {\raggedright\popsemi\fontsize{8.4}{11}\selectfont\color{white}All signed courses, unlimited AI tutor, PDF sheets — OnBuch+ tab in the app\par}
      \end{minipage}
    \end{tcolorbox}
  \end{minipage}
  \vspace{8pt}
  \popcontenu
  \vfill
  \signatureonbuch
  \restoregeometry
  \newpage
}

% Rangée « Ce cours contient » (page de garde)
\newcommand{\poptile}[3]{% #1 couleur #2 gros texte #3 légende
  \begin{tcolorbox}[enhanced,colback=#1,colframe=popink,boxrule=2pt,arc=9pt,shadow={2.5pt}{-2.5pt}{0pt}{popink},
      left=6pt,right=6pt,top=9pt,bottom=9pt,halign=center,valign=center,height=1.75cm,width=\linewidth,nobeforeafter]
    {\popdisplay\fontsize{12.5}{13}\selectfont\color{white}\MakeUppercase{#2}}\par\vspace{3pt}
    {\centering\popsemi\fontsize{7.8}{9.5}\selectfont\color{white}#3\par}
  \end{tcolorbox}}
\newcommand{\popcontenu}{%
  \par\noindent{\popxboldls\fontsize{7.5}{7.5}\selectfont\color{popmuted}THIS SIGNED COURSE CONTAINS}\par\vspace{7pt}
  \noindent\begin{minipage}[t]{0.235\linewidth}\poptile{popblue}{Course}{complete, illustrated, step by step}\end{minipage}\hfill
  \begin{minipage}[t]{0.235\linewidth}\poptile{poporange}{Methods}{and worked examples}\end{minipage}\hfill
  \begin{minipage}[t]{0.235\linewidth}\poptile{poppink}{Exercises}{exam style + solutions}\end{minipage}\hfill
  \begin{minipage}[t]{0.235\linewidth}\poptile{popgreen}{Summary sheet}{the essentials to remember}\end{minipage}\par}

% Sommaire
\newtcolorbox{sommaire}{enhanced,colback=white,colframe=popink,boxrule=2.2pt,arc=10pt,
  drop shadow=popink,left=18pt,right=18pt,top=13pt,bottom=13pt,before skip=6pt,after skip=20pt,
  before upper={\poppill{Contents}{poppurple}{white}\par\vspace{9pt}\popsemi\fontsize{10.6}{14.5}\selectfont\color{popink}}}

% Signature de fin de cours — poussée tout en bas de la dernière page
\newcommand{\findecours}{%
  \par\vfill\nopagebreak
  \begin{center}\popsquiggle{poporange}\end{center}\vspace{4pt}
  \signatureonbuch
}
\AtBeginDocument{\def\llbracket{\mathopen{[\mkern-3.5mu[}}\def\rrbracket{\mathclose{]\mkern-3.5mu]}}} % intervalles d'entiers (glyphe ⟦ absent de la police)
% Glyphes absents de Fira Math : redéfinis à partir de glyphes disponibles
\AtBeginDocument{%
  \def\setminus{\mathbin{\backslash}}\let\smallsetminus\setminus
  \def\vdots{\mathord{\vbox{\baselineskip3.2pt\lineskiplimit0pt\kern4pt\hbox{.}\hbox{.}\hbox{.}}}}%
  \def\ddots{\mathinner{\mkern1mu\raise6.5pt\hbox{.}\mkern2mu\raise3.6pt\hbox{.}\mkern2mu\raise0.7pt\hbox{.}\mkern1mu}}%
  \def\triangle{\mathord{\scalebox{0.9}{$\symup{\Delta}$}}}%
  \def\nabla{\mathord{\raisebox{\depth}{\scalebox{1}[-1]{$\symup{\Delta}$}}}}%
  \def\checkmark{\popcheck{popdark}}%
  \def\star{\mathbin{\ast}}\def\bigstar{\mathord{\ast}}%
  \def\diamond{\mathbin{\scalebox{0.6}{$\Diamond$}}}%
  \def\bigcup{\mathop{\mathchoice{\raisebox{-0.3em}{\scalebox{1.7}{$\cup$}}}{\scalebox{1.25}{$\cup$}}{\cup}{\cup}}\displaylimits}%
  \def\bigcap{\mathop{\mathchoice{\raisebox{-0.3em}{\scalebox{1.7}{$\cap$}}}{\scalebox{1.25}{$\cap$}}{\cap}{\cap}}\displaylimits}%
  \def\lesseqqgtr{\mathrel{\lessgtr}}\def\bigoplus{\mathop{\oplus}\displaylimits}%
}
\providecommand{\ssi}{if and only if} % abréviation fréquente
\AtBeginDocument{\providecommand{\wideparen}{\overparen}} % arc de cercle

% Fonctions usuelles que les modèles utilisent sans que amsmath les fournisse
\providecommand{\arsinh}{\operatorname{arsinh}}\providecommand{\arcosh}{\operatorname{arcosh}}
\providecommand{\artanh}{\operatorname{artanh}}\providecommand{\arsech}{\operatorname{arsech}}
\providecommand{\arcsch}{\operatorname{arcsch}}\providecommand{\arcoth}{\operatorname{arcoth}}
\providecommand{\arcsinh}{\operatorname{arsinh}}\providecommand{\arccosh}{\operatorname{arcosh}}
\providecommand{\arctanh}{\operatorname{artanh}}\providecommand{\sech}{\operatorname{sech}}
\providecommand{\csch}{\operatorname{csch}}\providecommand{\arcsec}{\operatorname{arcsec}}
\providecommand{\arccsc}{\operatorname{arccsc}}\providecommand{\arccot}{\operatorname{arccot}}
\providecommand{\sgn}{\operatorname{sgn}}\providecommand{\lcm}{\operatorname{lcm}}
\providecommand{\Var}{\operatorname{Var}}\providecommand{\Cov}{\operatorname{Cov}}
\providecommand{\permil}{\ensuremath{{}^{0}\!/_{00}}}\providecommand{\textperthousand}{\ensuremath{{}^{0}\!/_{00}}}

\begin{document}

\pagegarde{Applying project management fundamentals to run sample projects}{ICT}{Upper Sixth}{Upper Sixth — ICT}{20}{6 periods}{This chapter introduces the fundamentals of project management and shows how to plan, schedule and control a real ICT project. Learners master the project life cycle, the Work Breakdown Structure, Gantt and PERT charts, critical path and slack analysis, and apply them to a complete sample project.
\vspace{4pt}
\begin{multicols}{2}\small
\begin{itemize}
\item By the end of this chapter I can define a project, project management and distinguish a project from routine operations.
\item By the end of this chapter I can describe the phases of the project life cycle and the role of the project manager and stakeholders.
\item By the end of this chapter I can explain the triple constraint (scope, time, cost) and the key knowledge areas of project management.
\item By the end of this chapter I can decompose a project into a Work Breakdown Structure (WBS) with work packages.
\item By the end of this chapter I can build and read a Gantt chart to schedule project activities.
\item By the end of this chapter I can construct a PERT/CPM network diagram from a task list with dependencies and durations.
\end{itemize}\end{multicols}}

\begin{sommaire}
\begin{multicols}{2}
\begin{enumerate}
\item Fundamentals of Project Management
\item Project Planning and the Work Breakdown Structure (WBS)
\item Project Scheduling with Gantt and PERT Charts
\item Project Control, Critical Path and Slack Analysis
\item Integration: Running a Complete Sample Project
\item Integration activity
\item Methods and worked examples
\item Exercises
\item Solutions to the exercises
\item Summary sheet
\end{enumerate}
\end{multicols}
\end{sommaire}

\begin{objectifs}
\begin{itemize}
\item By the end of this chapter I can define a project, project management and distinguish a project from routine operations.
\item By the end of this chapter I can describe the phases of the project life cycle and the role of the project manager and stakeholders.
\item By the end of this chapter I can explain the triple constraint (scope, time, cost) and the key knowledge areas of project management.
\item By the end of this chapter I can decompose a project into a Work Breakdown Structure (WBS) with work packages.
\item By the end of this chapter I can build and read a Gantt chart to schedule project activities.
\item By the end of this chapter I can construct a PERT/CPM network diagram from a task list with dependencies and durations.
\item By the end of this chapter I can compute earliest/latest dates, identify the critical path and calculate slack (float) of activities.
\item By the end of this chapter I can monitor and control a project using milestones, progress reports and corrective actions.
\item By the end of this chapter I can plan and manage a small sample ICT project from initiation to closure.
\end{itemize}
\end{objectifs}

\begin{prerequis}
\begin{itemize}
\item Basic ICT project vocabulary seen in Lower Sixth (needs analysis, specifications, testing).
\item Reading and drawing simple diagrams and tables.
\item Elementary arithmetic: sums, differences, maximum and minimum of numbers.
\item Notion of algorithms: sequencing steps and dependencies between tasks.
\item Spreadsheet basics (rows, columns, simple formulas) useful for drawing Gantt charts.
\end{itemize}
\end{prerequis}

\newpage

\coursec{Fundamentals of Project Management}

In Yaoundé, a small cybercafé owner, Mr.~Tchoupo, wants to open a second branch in Bafoussam in \emph{exactly three months}. Computers must be bought, a room rented, an Internet line installed by an ISP, and staff trained. Without a plan, the money runs out before the opening day: the room is paid for but the computers arrive late, the ISP technician comes when nobody is on site, and the opening is postponed again and again. With a plan, every franc and every day is accounted for, and the branch opens on schedule. Big projects like the Lom Pangar dam or the rollout of the national fibre-optic backbone succeed or fail for exactly the same reason: \textbf{how well they are managed}.

\textbf{Problem:} how can we organise a temporary, one-off effort --- with limited money, limited time and many people involved --- so that it delivers exactly what was promised? The answer is the discipline called \cle{project management}, and this section lays its foundations.

\courssub{What is a Project?}

Before managing a project, we must be able to \emph{recognise} one. Consider two activities of a school: (a) building a new science block, and (b) running the school canteen. Building the science block starts one day, ends another day, and produces something that never existed before. Running the canteen repeats the same tasks every day, indefinitely. Only the first is a project.

\begin{definition}[Project]
A \cle{project} is a \textbf{temporary} endeavour, with a definite \textbf{start} and a definite \textbf{end}, undertaken to create a \textbf{unique} product, service or result.
\end{definition}

Three characteristics flow directly from this definition:

\begin{itemize}
\item \textbf{Uniqueness:} the result has never been produced in exactly this way before. Even if Mr.~Tchoupo opens a third branch later, the location, room and staff will differ.
\item \textbf{Temporariness:} the project ends when its objective is reached (or abandoned). ``Temporary'' does not mean short: a dam can take years.
\item \textbf{Progressive elaboration:} the details are developed step by step as the project advances. At the start, Mr.~Tchoupo knows only ``a cybercafé in Bafoussam''; later he knows the exact room, the exact number of computers and the exact tariff plan.
\end{itemize}

\begin{attention}[Project or routine operation?]
A very common GCE mistake is to confuse a \cle{project} (temporary, unique) with \cle{routine operations} (repetitive, ongoing, no defined end). Building the school is a project; teaching in it every morning is an operation. Installing a VSAT link for a bank branch is a project; operating that link daily is an operation. In the exam, always justify your classification using \emph{temporariness} and \emph{uniqueness}.
\end{attention}

\begin{center}
\begin{tabularx}{\linewidth}{|l|X|X|}
\hline
\rowcolor{popblueL} \textbf{Criterion} & \textbf{Project} & \textbf{Routine operation} \\
\hline
Duration & Temporary, defined start and end & Ongoing, no planned end \\
\hline
Result & Unique product, service or result & Repetitive, identical output \\
\hline
Team & Assembled for the project, dissolved after & Permanent staff \\
\hline
Example & Opening a new cybercafé branch & Running the existing cybercafé daily \\
\hline
\end{tabularx}
\end{center}

\courssub{What is Project Management?}

Knowing that a project exists is not enough: someone must make it happen. That is the job of project management.

\begin{definition}[Project management]
\cle{Project management} is the application of \textbf{knowledge}, \textbf{skills}, \textbf{tools} and \textbf{techniques} to project activities in order to \textbf{meet the project requirements}.
\end{definition}

In practice this means answering, at every moment, four questions: \emph{What must be done?} (scope), \emph{By when?} (time), \emph{With what resources and budget?} (cost), and \emph{How well?} (quality). Tools such as the WBS, the Gantt chart and the PERT chart, studied in the next sections, exist precisely to answer these questions rigorously instead of by guesswork.

\courssub{The Triple Constraint: Scope, Time, Cost}

Every project is squeezed between three demands that pull in opposite directions. Mr.~Tchoupo wants \emph{many services} (scope), \emph{opening in three months} (time) and \emph{spending little} (cost). He cannot maximise all three at once.

\begin{definition}[The triple constraint]
The \cle{triple constraint} states that every project is bounded by three interdependent constraints: \textbf{scope} (the work and features to deliver), \textbf{time} (the schedule) and \textbf{cost} (the budget). \textbf{Quality} sits at the centre, affected by all three.
\end{definition}

\begin{propriete}[Interdependence of the constraints]
Any change to one of the three constraints (scope, time, cost) impacts at least one of the other two, and usually quality as well. (Statement to know and to be able to illustrate; no formal proof is examinable, but a justified example is.)
\end{propriete}

\begin{popfigure}
\begin{tikzpicture}[scale=1.05]
\coordinate (S) at (90:3.2);
\coordinate (T) at (210:3.2);
\coordinate (C) at (330:3.2);
\draw[very thick, popblue] (S) -- (T) -- (C) -- cycle;
\node[above, font=\bfseries, popblue] at (S) {SCOPE};
\node[below left, font=\bfseries, popblue] at (T) {TIME};
\node[below right, font=\bfseries, popblue] at (C) {COST};
\node[circle, draw=poporange, very thick, fill=popblueL, minimum size=1.6cm, font=\bfseries, align=center] at (0,0) {QUALITY};
\draw[->, thick, poporange] (S) -- (0,1.0);
\draw[->, thick, poporange] (T) -- (-0.9,-0.5);
\draw[->, thick, poporange] (C) -- (0.9,-0.5);
\end{tikzpicture}
\legende{The triple constraint: scope, time and cost bound the project; quality sits at the centre and suffers when the triangle is unbalanced.}
\end{popfigure}

\begin{exemplebox}[Playing with the constraints]
Mr.~Tchoupo's client (his own business plan) suddenly demands \textbf{20 computers instead of 12} (scope increases). Consequences: either the \textbf{cost} rises (more machines to buy), or the \textbf{time} stretches (longer to raise the money and install), or \textbf{quality} falls (cheaper, slower machines). Conversely, if the opening must be \textbf{one month earlier} (time decreases), he must either pay for express delivery (cost up) or open with fewer services (scope down). There is no free lunch: the triangle always rebalances.
\end{exemplebox}

\courssub{The Project Life Cycle}

Every project, from a cybercafé to a dam, passes through the same five phases. Knowing them tells you \emph{where you are} and \emph{what document you should be producing}.

\begin{definition}[Project life cycle]
The \cle{project life cycle} is the sequence of phases a project goes through from its conception to its closure: \textbf{initiation}, \textbf{planning}, \textbf{execution}, \textbf{monitoring and control}, and \textbf{closure}.
\end{definition}

\begin{popfigure}
\begin{tikzpicture}[node distance=0.55cm,
phase/.style={rectangle, draw=popblue, very thick, fill=popblueL, rounded corners, minimum width=2.5cm, minimum height=1.0cm, align=center, font=\small\bfseries}]
\node[phase] (ini) {1. Initiation};
\node[phase, right=of ini] (pla) {2. Planning};
\node[phase, right=of pla] (exe) {3. Execution};
\node[phase, right=of exe] (mon) {4. Monitoring \& Control};
\node[phase, right=of mon] (clo) {5. Closure};
\draw[->, very thick, poporange] (ini) -- (pla);
\draw[->, very thick, poporange] (pla) -- (exe);
\draw[->, very thick, poporange] (exe) -- (mon);
\draw[->, very thick, poporange] (mon) -- (clo);
\draw[->, very thick, poppurple, dashed] (mon.north) to[bend left=35] node[above, font=\scriptsize]{corrective action} (exe.north);
\end{tikzpicture}
\legende{The five phases of the project life cycle. Monitoring and control runs alongside execution and can send the project back for corrective action.}
\end{popfigure}

\begin{aretenir}[Phases and their main deliverables]
\begin{itemize}
\item \textbf{Initiation:} the idea is validated and authorised. Deliverable: the \emph{project charter} or terms of reference (objectives, sponsor, rough budget).
\item \textbf{Planning:} the work is decomposed (WBS), scheduled (Gantt, PERT) and budgeted. Deliverable: the \emph{project plan}.
\item \textbf{Execution:} the planned work is actually carried out. Deliverable: the \emph{product, service or result} itself (the fitted-out cybercafé).
\item \textbf{Monitoring and control:} progress, cost and quality are measured against the plan and deviations corrected. Deliverable: \emph{progress reports} and corrective decisions.
\item \textbf{Closure:} the result is accepted, contracts closed, lessons recorded. Deliverable: \emph{acceptance certificate} and \emph{lessons-learned report}.
\end{itemize}
\end{aretenir}

\begin{attention}[Skipping the planning phase]
Under pressure, teams jump straight from initiation to execution: ``we have no time to plan.'' This is the single most frequent cause of failure. Time ``saved'' in planning is repaid many times over in execution, in rework and wasted materials. In the GCE, if a scenario describes a failing project, check first whether planning was done properly.
\end{attention}

\courssub{The Project Manager and the Project Team}

\begin{definition}[Project manager]
The \cle{project manager} is the person assigned by the sponsoring organisation to lead the team and to be responsible for achieving the project objectives within the constraints of scope, time, cost and quality.
\end{definition}

The project manager's main \textbf{responsibilities} are:

\begin{itemize}
\item define and clarify objectives with the sponsor;
\item plan the work (WBS, schedule, budget) and keep the plan up to date;
\item assemble, coordinate and motivate the \cle{project team};
\item communicate with all stakeholders and manage their expectations;
\item monitor progress, identify risks and take corrective action;
\item deliver the result and close the project formally.
\end{itemize}

The essential \textbf{qualities} are leadership, communication, organisation, negotiation, and enough technical understanding of the field (here, ICT) to make sound judgements. The \textbf{project team} is the group of people --- technicians, trainers, buyers, an accountant --- assembled temporarily to carry out the work; it is normally dissolved at closure.

\courssub{Stakeholders and Their Expectations}

A project never lives in a vacuum. Many people can affect it or are affected by it, and each one expects something different.

\begin{definition}[Stakeholder]
A \cle{stakeholder} is any person, group or organisation that can \textbf{affect}, is \textbf{affected by}, or believes itself affected by the project: the sponsor, the users, the suppliers, the project team, the administration, the local community, and so on.
\end{definition}

\begin{popfigure}
\begin{tikzpicture}
\node[circle, draw=popblue, very thick, fill=popblueL, minimum size=2.4cm, align=center, font=\bfseries] (P) at (0,0) {Cybercafé\\project};
\foreach \a/\t in {90/Sponsor (owner), 30/Customers (users), -30/ISP \& suppliers, -90/Landlord, -150/Project team \& staff, 150/Local community \& council}{
  \node[rectangle, draw=popgreen, thick, fill=popgreenL, rounded corners, align=center, font=\small, text width=2.6cm] (s) at (\a:4.3) {\t};
  \draw[<->, thick, popmuted] (P) -- (s);
}
\end{tikzpicture}
\legende{Stakeholder map: the project sits at the centre; each stakeholder exchanges influence and expectations with it.}
\end{popfigure}

\begin{methode}[Identifying the stakeholders of a project in 4 steps]
\begin{enumerate}
\item \textbf{List} everyone who can influence the project or who is touched by its result: sponsor, users, team, suppliers, authorities, community.
\item \textbf{Analyse} each stakeholder: what do they expect? What power do they have? Are they favourable or hostile to the project?
\item \textbf{Classify} them by importance (power and interest) to decide who must be managed closely and who simply kept informed.
\item \textbf{Plan communication}: decide what to tell each stakeholder, how often, and through which channel (meetings, reports, phone calls).
\end{enumerate}
\end{methode}

Managing expectations means preventing unpleasant surprises: if the landlord expects the rent before the 5th of each month and the sponsor expects an opening ceremony, both must appear in the plan.

\courssub{Success and Failure Factors}

\begin{exemplebox}[Two Cameroonian stories]
\textbf{Failure.} All over the country one sees abandoned construction sites: a market or a hall whose walls stop at one metre. Typical causes: no clear scope (the building grows as money comes), no realistic budget, no monitoring, and stakeholders (contractor, council) who stopped communicating. \textbf{Success.} A bank rolling out VSAT links to its rural branches succeeds when each installation is treated as a small project: a precise scope per branch, a schedule agreed with the satellite provider, a trained team, and a checklist signed at each closure. Same country, same constraints --- different management.
\end{exemplebox}

\begin{aretenir}[What makes projects succeed or fail]
\textbf{Success factors:} clear and agreed objectives; realistic scope, budget and schedule; competent and available project manager; strong sponsor support; early identification and management of stakeholders; continuous monitoring and honest reporting.\\
\textbf{Failure factors:} vague or changing objectives; underestimated costs and durations; planning skipped; poor communication; risks ignored; scope allowed to grow uncontrolled (\emph{scope creep}).
\end{aretenir}

\begin{exoresolu}[Is it a project? Classify and justify]
For each activity, say whether it is a project or a routine operation, and justify: (a) developing a mobile-money awareness campaign for a microfinance in Douala; (b) processing daily mobile-money transactions at the same microfinance; (c) connecting a school in Bamenda to the fibre-optic network.
\tcblower
\textbf{Solution.}
\begin{itemize}
\item (a) \textbf{Project:} it is temporary (the campaign has a launch and an end date) and unique (this specific campaign, for this audience, has never been run).
\item (b) \textbf{Routine operation:} the same transactions are processed every day, repetitively, with no planned end; the output is not unique.
\item (c) \textbf{Project:} temporary (a start and a hand-over date) and unique (this school, this route, this configuration). Operating the link afterwards would be a routine operation.
\end{itemize}
The justification must always invoke \emph{temporariness} and \emph{uniqueness} --- that is what the examiner looks for.
\end{exoresolu}

\begin{savaistu}[Waterfall and agile methodologies]
Beyond the GCE syllabus, professionals organise the life cycle in two broad styles. The \textbf{waterfall} method runs the phases strictly one after another: suitable when requirements are stable (a building, a network installation). The \textbf{agile} methods (e.g. Scrum) deliver the product in short repeated cycles called iterations, welcoming change: suitable for software, where users discover their needs progressively. Many real projects mix both. You only need to recognise the names and the idea.
\end{savaistu}

\begin{exercice}[Applying the fundamentals]
Mr.~Tchoupo finally launches his Bafoussam branch project.
\begin{enumerate}
\item Identify four stakeholders of this project and state one expectation of each.
\item He decides, two weeks before opening, to add a printing and photocopying service. Using the triple constraint, explain two possible consequences.
\item For each of the five phases of the life cycle, give one concrete activity of this project.
\end{enumerate}
\end{exercice}

\begin{corrige}[Exercise 1]
\begin{enumerate}
\item Possible stakeholders: the sponsor/owner (expects profitability), the customers (expect reliable, affordable Internet), the ISP (expects a signed contract and payment), the landlord (expects rent on time), the staff (expect training and salaries), the Bafoussam council (expects taxes and a legal business).
\item Scope increases, so the triangle must rebalance: either \textbf{cost} rises (buying a printer/photocopier) or \textbf{time} stretches (opening delayed to install and test the new service); if neither budget nor date moves, \textbf{quality} of the rest will suffer (less money and attention for the computers).
\item Initiation: writing the project charter (objective, budget ceiling). Planning: drawing the WBS, the budget and a three-month schedule. Execution: renting the room, buying and installing computers, training staff. Monitoring and control: weekly comparison of spending and progress against the plan, corrective decisions. Closure: testing all services, signing off the installation, opening ceremony and lessons-learned note.
\end{enumerate}
\end{corrige}

\coursec{Project Planning and the Work Breakdown Structure (WBS)}

In the previous section we discovered the fundamentals of project management: what a project is, the triple constraint (scope, time, cost), the project life cycle and the role of the project manager. We saw that a project is a \emph{temporary} undertaking with a \emph{unique} objective. But between the decision to launch a project and the actual execution of the work, there is an essential step that distinguishes successful projects from failed ones: \cle{planning}. This section answers a very practical question: \emph{how do we move from a vague idea --- ``set up a computer laboratory'' --- to a precise, organised list of tasks that a team can execute?} The central tool for this is the \cle{Work Breakdown Structure} (WBS).

\courssub{Why Plan? Objectives, Scope and Deliverables}

\textbf{An observation to start with.} Imagine a school principal in Bamenda who tells the PTA: ``This year, we shall computerise the school.'' Six months later, money has been spent, three computers have been bought, a room has been painted, but no laboratory exists and nobody knows what remains to be done. What went wrong? The project started \emph{without a plan}: nobody had defined precisely what ``computerise the school'' meant, what results were expected, or by when.

\begin{definition}[Project planning]
\cle{Project planning} is the process of defining, \emph{before any work starts}, what the project must produce (objectives, scope, deliverables), how the work will be organised (tasks, schedule, resources) and how progress will be measured. Planning answers the questions: \emph{What? Why? How? By whom? When? At what cost?}
\end{definition}

The first output of planning is a set of well-formulated \cle{objectives}. A vague objective (``improve the network'') cannot be managed. Good practice requires objectives to be \cle{SMART}.

\begin{aretenir}[SMART objectives]
A project objective is \cle{SMART} when it is:
\begin{itemize}
\item \textbf{S}pecific: it states clearly and precisely what is to be achieved;
\item \textbf{M}easurable: its achievement can be verified with numbers or observable facts;
\item \textbf{A}chievable: it can be reached with the available skills and resources;
\item \textbf{R}ealistic: it is relevant and reasonable given the constraints (budget, context);
\item \textbf{T}ime-bound: it has a clear deadline.
\end{itemize}
\textbf{Example.} Not SMART: ``Get computers for the school.'' SMART: ``Install and commission a computer laboratory of 20 networked PCs with Internet access in GBHS Nkwen before 30 April, within a budget of 15 million FCFA.''
\end{aretenir}

Once objectives are fixed, the project manager writes the \cle{scope statement}: a short document that defines the \emph{boundaries} of the project --- what is included and, just as importantly, what is \emph{excluded}. From the scope statement we derive the \cle{deliverables}: the tangible, verifiable results the project must hand over (for example: a cabled room, 20 installed PCs, a configured LAN, a training report, an acceptance certificate).

\begin{attention}[Scope creep]
A very common cause of project failure is \cle{scope creep}: new requests are added during execution (``while you are at it, also connect the staff room and the library!'') without adjusting time or budget. A clear scope statement, agreed and signed before work starts, is the project manager's protection. Everything in the plan must serve the scope --- and nothing outside the scope should consume project resources.
\end{attention}

\courssub{From Objectives to Tasks: the Logic of Decomposition}

Even with a SMART objective, the project remains too big to execute in one step. Nobody can ``do'' a computer laboratory directly. The natural human strategy, used in mathematics (dividing a difficult problem into sub-problems) as in everyday life (organising a wedding hall by hall, meal by meal), is \emph{divide and conquer}.

\begin{definition}[Decomposition]
\cle{Decomposition} is the technique of breaking down the project into smaller and smaller components, level by level, until each component is simple enough to be estimated, assigned to one responsible person, scheduled and controlled.
\end{definition}

Decomposition has three major benefits:
\begin{itemize}
\item \textbf{Nothing is forgotten}: by covering the project systematically, hidden work (delivery, testing, documentation) is revealed;
\item \textbf{Estimates become reliable}: it is impossible to estimate ``build a laboratory'', but easy to estimate ``install one operating system on one PC'';
\item \textbf{Responsibilities become clear}: each small component can be assigned to a specific person or team.
\end{itemize}

The structured result of this decomposition is the \cle{Work Breakdown Structure}.

\courssub{The Work Breakdown Structure (WBS)}

\begin{definition}[Work Breakdown Structure (WBS)]
The \cle{Work Breakdown Structure} (WBS) is a \emph{deliverable-oriented hierarchical decomposition} of \textbf{all} the work to be carried out by the project team. It organises and defines the total scope of the project as a tree: the project at the top, then phases or major deliverables, then work packages, and finally the activities needed to produce them.
\end{definition}

The key phrase is \emph{deliverable-oriented}: the WBS decomposes the project into \textbf{results} (nouns: ``a cabled room'', ``20 installed PCs''), not into actions (verbs). The actions come later, at the activity level.

\textbf{The levels of a WBS.} A typical WBS has the following levels:
\begin{enumerate}
\item \textbf{Level 1 --- the project}: the complete project itself (e.g.\ ``Computer laboratory of 20 PCs'');
\item \textbf{Level 2 --- phases or major deliverables}: the main results that together constitute the project (e.g.\ ``Room preparation'', ``Equipment'', ``Network'');
\item \textbf{Level 3 --- work packages}: the lowest level of the WBS; each one is a chunk of work small enough to be managed;
\item \textbf{Below the WBS --- activities}: the concrete tasks (with durations) needed to complete each work package. Activities belong to the schedule, strictly speaking, but they are derived directly from the work packages.
\end{enumerate}

\begin{definition}[Work package]
A \cle{work package} is the lowest-level component of the WBS. It must be:
\begin{itemize}
\item \textbf{measurable}: one can tell objectively when it is complete;
\item \textbf{assignable}: it can be given to one person or one team who is responsible for it;
\item \textbf{estimable}: its duration and its cost can be estimated with reasonable confidence.
\end{itemize}
If a component is still too big or too vague to satisfy these three properties, it must be decomposed further.
\end{definition}

\begin{propriete}[The 100\% rule]
The WBS must capture \textbf{100\% of the work} of the project --- \emph{nothing more, nothing less}. The sum of all the components at any level must equal exactly the work of their parent component: all deliverables of the scope (including project management work itself) must appear, and no element outside the scope may be added. (Statement to know; not a formal proof at GCE level, but its application is examinable.)
\end{propriete}

\begin{attention}[Common mistake: listing actions instead of deliverables]
The most frequent error when building a WBS is to write \emph{actions} (verbs) at the upper levels: ``Buy computers'', ``Paint the room'', ``Configure the network''. A WBS is \textbf{deliverable-oriented}: write \emph{results} (nouns): ``Equipment purchased and delivered'', ``Room prepared'', ``Network operational''. Verbs are reserved for the \emph{activities} derived from the work packages. A quick test: a WBS element should answer ``\emph{what} must exist at the end?'', not ``\emph{what} do we do?''.
\end{attention}

\begin{methode}[Five steps to build a WBS from the scope statement]
\begin{enumerate}
\item \textbf{Read the scope statement} and list the major deliverables of the project (including project management).
\item \textbf{Place the project at Level 1} and the major deliverables (or phases) at Level 2.
\item \textbf{Decompose each Level 2 element} into work packages until each package is measurable, assignable and estimable.
\item \textbf{Verify the 100\% rule}: check that the children of each element cover all of its work, and that nothing outside the scope has slipped in.
\item \textbf{Number the elements} (1, 1.1, 1.1.1, \dots) and write the WBS dictionary describing each work package.
\end{enumerate}
\end{methode}

\courssub{Worked Example: WBS of a Computer Laboratory Project}

Let us apply the method to our running project: \emph{``Setting up a computer laboratory of 20 PCs in a college''}. The scope statement includes: preparing the room, supplying and installing 20 PCs, networking them with Internet access, and managing the project. Training of teachers and maintenance after handover are \emph{excluded} from the scope.

Applying steps 1 to 3, we obtain four Level 2 deliverables, each decomposed into work packages:

\begin{popfigure}
\begin{center}
\begin{tikzpicture}[
  grow=down,
  level distance=1.5cm,
  level 1/.style={sibling distance=4.6cm},
  level 2/.style={sibling distance=2.2cm, level distance=1.7cm},
  every node/.style={draw=popblue, fill=popblueL, rounded corners=2pt,
    font=\scriptsize, align=center, text width=1.9cm, inner sep=2pt},
  edge from parent/.style={draw=popmuted, thick}
]
\node[fill=popblue, text=white, text width=3.4cm]{\textbf{1. Computer laboratory of 20 PCs}}
  child { node {\textbf{1.1} Room prepared}
    child { node {1.1.1 Electrical wiring}}
    child { node {1.1.2 Furniture installed}}
  }
  child { node {\textbf{1.2} Equipment delivered}
    child { node {1.2.1 20 PCs supplied}}
    child { node {1.2.2 Peripherals supplied}}
  }
  child { node {\textbf{1.3} Network operational}
    child { node {1.3.1 LAN cabled}}
    child { node {1.3.2 Internet configured}}
  }
  child { node {\textbf{1.4} Project managed}
    child { node {1.4.1 Planning done}}
    child { node {1.4.2 Acceptance report}}
  };
\end{tikzpicture}
\end{center}
\legende{A 3-level WBS for the computer laboratory project: project (Level 1), major deliverables (Level 2), work packages (Level 3).}
\end{popfigure}

\textbf{Numbering and coding.} Notice the hierarchical code: the project is \textbf{1}, each deliverable is \textbf{1.1}, \textbf{1.2}, \dots, and each work package inherits its parent's code plus one more level: \textbf{1.3.1} is the first work package of deliverable \textbf{1.3}. This coding has two advantages: every element has a unique identifier (used in budgets, schedules and reports), and the code itself shows the position of the element in the tree.

\textbf{The WBS dictionary.} A diagram alone is not enough: two people can read ``Network operational'' and understand different things. The \cle{WBS dictionary} is a companion document that describes each work package precisely: its code, its content, its deliverable, the person responsible, and its estimated duration and cost.

\begin{popfigure}
\begin{center}
\begin{tabularx}{\linewidth}{|l|X|l|l|}
\hline
\rowcolor{popblueL} \textbf{Code} & \textbf{Description of the work package} & \textbf{Responsible} & \textbf{Est.\ cost} \\
\hline
1.1.1 & Rewire the room: 25 power sockets, circuit breaker, earthing & Electrician & 350\,000 F \\
\hline
1.1.2 & Supply and install 20 computer tables and 20 chairs & PTA member & 600\,000 F \\
\hline
1.2.1 & Purchase and deliver 20 desktop PCs (with OS licences) & Supplier & 6\,000\,000 F \\
\hline
1.3.1 & Cable the LAN: switch, 20 Ethernet points, testing & Network tech. & 450\,000 F \\
\hline
\end{tabularx}
\end{center}
\legende{Extract of the WBS dictionary for the computer laboratory project.}
\end{popfigure}

\courssub{From Work Packages to Activities: Estimation and Dependencies}

The WBS tells us \emph{what} must exist; it does not yet tell us \emph{in what order} to work or \emph{how long} it will take. The next planning step is to derive, from each work package, the \cle{activities} (concrete tasks expressed with verbs) needed to complete it, and to estimate for each activity its \cle{duration} and its \cle{resources} (people, equipment, money).

\begin{exemplebox}[From work package to activities]
Work package \textbf{1.3.1 ``LAN cabled''} decomposes into activities such as: \emph{draw the cabling plan} (1 day), \emph{install the switch and patch panel} (1 day), \emph{run and terminate 20 Ethernet cables} (2 days), \emph{test all network points} (1 day). Each activity is short enough to be estimated reliably and assigned to one technician.
\end{exemplebox}

Activities are rarely independent: most of the time, one activity cannot start before another finishes. The most common relationship is \cle{finish-to-start} (FS): activity $B$ can start only when activity $A$ is finished. For example, ``install PCs on tables'' (B) requires ``deliver PCs'' (A) to be finished: we write $A \rightarrow B$. Other relationships exist (start-to-start, finish-to-finish), but at GCE level we use finish-to-start. Recording all these links produces the \cle{precedence table}, which is the starting point for the Gantt and PERT charts of the next section.

\begin{popfigure}
\begin{center}
\begin{tabularx}{\linewidth}{|c|X|c|c|}
\hline
\rowcolor{popblueL} \textbf{Code} & \textbf{Activity} & \textbf{Duration (days)} & \textbf{Predecessors} \\
\hline
A & Prepare the room (wiring, painting) & 5 & --- \\
\hline
B & Purchase and deliver the 20 PCs & 10 & --- \\
\hline
C & Install furniture (tables, chairs) & 2 & A \\
\hline
D & Cable the LAN and install the switch & 3 & A \\
\hline
E & Install and configure the PCs & 4 & B, C \\
\hline
F & Configure Internet access & 1 & D \\
\hline
G & Test the whole laboratory & 2 & E, F \\
\hline
H & Write the acceptance report & 1 & G \\
\hline
\end{tabularx}
\end{center}
\legende{Precedence table for the computer laboratory project: each activity has a code, a duration and the list of its immediate predecessors (finish-to-start).}
\end{popfigure}

Reading the table: activity E (install and configure the PCs) has \emph{two} predecessors, B and C, because the technicians need both the delivered PCs \emph{and} the installed furniture before they can work. Activities A and B have no predecessor: they can start on day 1, possibly in parallel. This simple observation --- that some activities can run in parallel while others must wait --- is the seed of the critical path analysis we shall develop later.

\begin{exoresolu}[Building a mini-WBS and a precedence table]
A school wants to organise a \emph{one-day ICT awareness seminar} for Form 5 students. The scope includes: a prepared venue, invited resource persons, and information material. (1) Propose a 2-level WBS (project + 3 deliverables). (2) List one activity per deliverable and give a plausible predecessor relationship. \tcblower
\textbf{Solution.} (1) Applying the five-step method, with deliverables as nouns:
\begin{itemize}
\item \textbf{1. ICT awareness seminar} (Level 1)
  \begin{itemize}
  \item \textbf{1.1} Venue prepared (hall booked, chairs and projector installed);
  \item \textbf{1.2} Resource persons confirmed (speakers invited and confirmed);
  \item \textbf{1.3} Information material produced (flyers and handouts printed).
  \end{itemize}
\end{itemize}
Check of the 100\% rule: the three deliverables cover the whole scope and add nothing outside it. (2) Activities (verbs) and dependencies: \emph{Print handouts} (from 1.3) depends on \emph{Confirm speakers} (from 1.2), because the handouts mention the speakers' names --- a finish-to-start link. \emph{Install the projector} (from 1.1) is independent and can be done in parallel. Hence a precedence table with, for example: A ``Confirm speakers'' (no predecessor), B ``Print handouts'' (predecessor A), C ``Install projector'' (no predecessor).
\end{exoresolu}

\begin{exercice}[WBS for a school website project]
Your college decides to create its first website presenting the school, its results and its announcements. (1) Write one SMART objective for this project. (2) Build a WBS with at least three Level 2 deliverables and two work packages under each, using correct numbering. (3) Choose one work package, list two activities, estimate their durations and state their dependency. (4) Explain why ``Write articles'' would be a poor Level 2 element and propose a better wording.
\end{exercice}

\begin{corrige}[Exercise 1]
(1) Example of a SMART objective: ``Publish a functional school website with the pages Home, Results and Announcements, online before the end of the second term, within a budget of 500\,000 FCFA.'' (2) Possible WBS: \textbf{1. School website}; \textbf{1.1} Content ready (1.1.1 Texts written, 1.1.2 Photos collected); \textbf{1.2} Website developed (1.2.1 Pages designed, 1.2.2 Site tested); \textbf{1.3} Site hosted online (1.3.1 Hosting purchased, 1.3.2 Site deployed). (3) For 1.3.2: ``Upload the site files'' (1 day) then ``Verify the site online'' (1 day); the second depends on the first (finish-to-start). (4) ``Write articles'' is an \emph{action} (verb), whereas the WBS must be deliverable-oriented; the correct Level 2 wording is the result: ``Content ready'' or ``Articles written''.
\end{corrige}

\begin{aretenir}[Key points of this section]
\begin{itemize}
\item Planning comes \emph{before} execution: SMART objectives, a scope statement and a list of deliverables are its first outputs.
\item The \cle{WBS} is a deliverable-oriented hierarchical decomposition of all the project work: project $\rightarrow$ deliverables $\rightarrow$ work packages $\rightarrow$ activities.
\item The \cle{100\% rule}: the WBS covers all the work, nothing more, nothing less.
\item A \cle{work package} is measurable, assignable and estimable; the \cle{WBS dictionary} describes each one.
\item From work packages we derive \cle{activities}, estimate their durations and resources, and record their \cle{dependencies} (mostly finish-to-start) in a \cle{precedence table} --- the bridge to Gantt and PERT scheduling, our next section.
\end{itemize}
\end{aretenir}

\coursec{Project Scheduling with Gantt and PERT Charts}

In the previous section we broke our project down with the Work Breakdown Structure (WBS): the big goal ``set up a computer laboratory'' became a tree of small, manageable work packages, each with an owner, a duration and a deliverable. But a WBS tells us \emph{what} must be done, not \emph{when}. A list of tasks without dates is like a timetable with no hours written on it: nobody knows when to come to class. In this section we learn to place our tasks in \textbf{time} using the two most famous scheduling tools in project management: the \cle{Gantt chart} and the \cle{PERT chart}.

\courssub{What is Scheduling?}

Imagine the WBS of our computer laboratory project has produced these tasks: prepare the site, install electrical wiring, paint the room, buy the computers, install the network, set up the hardware, install the software, train the teachers. Some questions immediately arise. Can we set up hardware before the computers arrive? Certainly not. Can we paint while the electrician works? Probably not safely. Can we buy computers while the room is being prepared? Yes, with different people.

\begin{definition}[Project scheduling]
\cle{Scheduling} is the process of placing the activities of a project in time: deciding when each activity \textbf{starts} and \textbf{ends}, which activities can run \textbf{in parallel}, which must \textbf{follow} others, and which \textbf{resources} (people, money, equipment) are assigned to each period.
\end{definition}

A good schedule answers three questions at any moment: \emph{What should be happening now? Who is doing it? Are we late?} The two tools we study answer these questions in complementary ways.

\courssub{The Gantt Chart}

\textbf{Intuition.} Take a sheet of squared paper. Write the task names down the left side, and the weeks of the project across the top. For each task, colour the cells of the weeks during which the task runs. You have just built a Gantt chart --- the most widely used planning tool in the world, invented by the engineer Henry Gantt around 1910.

\begin{definition}[Gantt chart]
A \cle{Gantt chart} is a bar chart that represents a project schedule. It has two axes:
\begin{itemize}
\item the \textbf{vertical axis} lists the \textbf{tasks} (usually in chronological order of start);
\item the \textbf{horizontal axis} is a \textbf{timeline} (days, weeks or months).
\end{itemize}
Each task is drawn as a horizontal \textbf{bar} whose position and length show its start date, end date and duration.
\end{definition}

\begin{definition}[Milestone]
A \cle{milestone} is a significant event in a project with \textbf{zero duration}: it marks the completion of an important stage or a key decision point (for example ``all computers delivered'' or ``laboratory officially opened''). On a Gantt chart it is drawn as a \textbf{diamond} ($\blacklozenge$) on the timeline.
\end{definition}

\textbf{Reading a Gantt chart.} To read a Gantt chart, pick a week on the timeline and look vertically: every bar crossing that week is a task in progress during that week. Overlapping bars reveal tasks done \textbf{in parallel}; a bar starting exactly where another ends reveals a \textbf{sequence}. Dependencies can be drawn as small arrows from the end of one bar to the start of the next.

\begin{popfigure}
\begin{tikzpicture}[x=0.72cm,y=0.62cm]
% timeline grid
\foreach \w in {1,...,12}{
  \draw[popline, very thin] (\w,0) -- (\w,-8.6);
  \node[above, font=\small\bfseries, text=popdark] at (\w-0.5,0) {W\w};
}
\draw[popline, very thin] (0,0) -- (12,0);
\draw[popline, very thin] (0,-8.6) -- (12,-8.6);
% task labels and bars
\node[right, font=\small, text=popink] at (-5.4,-0.5) {A. Site preparation};
\fill[popblue] (0,-0.75) rectangle (2,-0.25);
\node[right, font=\small, text=popink] at (-5.4,-1.5) {B. Electrical wiring};
\fill[popblue] (2,-1.75) rectangle (4,-1.25);
\node[right, font=\small, text=popink] at (-5.4,-2.5) {C. Painting};
\fill[popblue] (2,-2.75) rectangle (4,-2.25);
\node[right, font=\small, text=popink] at (-5.4,-3.5) {D. Purchase computers};
\fill[poporange] (0,-3.75) rectangle (5,-3.25);
\node[right, font=\small, text=popink] at (-5.4,-4.5) {E. Network installation};
\fill[popgreen] (4,-4.75) rectangle (6,-4.25);
\node[right, font=\small, text=popink] at (-5.4,-5.5) {F. Hardware setup};
\fill[popgreen] (6,-5.75) rectangle (8,-5.25);
\node[right, font=\small, text=popink] at (-5.4,-6.5) {G. Software installation};
\fill[popgreen] (8,-6.75) rectangle (10,-6.25);
\node[right, font=\small, text=popink] at (-5.4,-7.5) {H. Teacher training};
\fill[poppurple] (10,-7.75) rectangle (12,-7.25);
% milestone
\node[rotate=45, fill=popgold, draw=popdark, minimum size=0.28cm, inner sep=0pt] at (12,-8.25) {};
\node[right, font=\small\itshape, text=popdark] at (-5.4,-8.25) {Milestone: laboratory opens (end of W12)};
% dependency arrows
\draw[-latex, popmuted, thin] (2,-0.5) -- (2.15,-1.5);
\draw[-latex, popmuted, thin] (4,-1.5) -- (4.15,-4.5);
\draw[-latex, popmuted, thin] (5,-3.5) -- (6.15,-5.5);
\draw[-latex, popmuted, thin] (6,-4.5) -- (6.3,-5.5);
\draw[-latex, popmuted, thin] (8,-5.5) -- (8.15,-6.5);
\draw[-latex, popmuted, thin] (10,-6.5) -- (10.15,-7.5);
\end{tikzpicture}
\legende{Gantt chart of the computer laboratory project (W1--W12 = weeks). Bars show task durations; the gold diamond is the milestone ``laboratory opens''; thin arrows show the main dependencies. Note that F can only start in week 7, when both D (end of week 5) and E (end of week 6) are finished.}
\end{popfigure}

\begin{methode}[Drawing a Gantt chart]
\begin{enumerate}
\item List the tasks from the WBS, in the order they will start.
\item Estimate the \textbf{duration} of each task (in days or weeks).
\item Choose the timeline unit and draw the horizontal axis.
\item For each task, draw a bar starting at its planned start week, of length equal to its duration. A task may only start when \emph{all} its predecessors are finished; tasks that can run in parallel get bars on the same weeks.
\item Mark \textbf{milestones} with diamonds at the end of important stages.
\item Optionally, draw small arrows between bars to show the main dependencies.
\end{enumerate}
With a spreadsheet (Excel, LibreOffice Calc), the same chart is made by typing tasks in one column, week numbers across a row, and filling cells with colour --- or by using a stacked bar chart.
\end{methode}

\begin{propriete}[Advantages and limits of the Gantt chart]
\textbf{Advantages:} simple to draw and to read; very visual (anyone, even a non-specialist like a school principal, understands it at a glance); shows clearly the calendar, parallel work and progress (bars can be partly shaded to show percentage completed).\\
\textbf{Limits:} dependencies between tasks are not shown rigorously (arrows quickly become unreadable on large projects); it does \textbf{not} reveal which tasks are \textbf{critical} (i.e.\ which delays will delay the whole project); it becomes clumsy for projects with hundreds of tasks.
\end{propriete}

\courssub{The PERT Chart}

\textbf{Intuition.} The Gantt chart answers ``when?'', but a project manager also needs to answer ``what blocks what?''. If the delivery of computers is late, which tasks suffer? To reason about dependencies, we draw the project as a \textbf{network}: circles for the \emph{states} of the project (events) and arrows for the \emph{activities} that move the project from one state to the next.

\begin{definition}[PERT chart]
A \cle{PERT chart} (Program Evaluation and Review Technique) is a network diagram representing a project as a directed graph:
\begin{itemize}
\item \textbf{nodes} (circles) represent \textbf{events}: instants at which certain activities are finished and others can begin. Each node is split to record the \textbf{earliest date} and the \textbf{latest date} at which the event can occur;
\item \textbf{arrows} represent \textbf{activities}, labelled with the activity name and its \textbf{duration};
\item a \textbf{dummy activity} (dashed arrow, duration 0) is added when a dependency must be shown but no real work exists.
\end{itemize}
An activity $(i,j)$ can only start when \emph{all} activities arriving at node $i$ are finished.
\end{definition}

\begin{methode}[Constructing a PERT network from a precedence table]
\begin{enumerate}
\item Read the precedence table (each task with its predecessors and duration).
\item Draw a single \textbf{start node} (event 0). From it, draw one arrow for each task that has \textbf{no predecessor}.
\item Take the tasks in order: the arrow of a task starts at the node where \emph{all} its predecessors end. If a task has several predecessors, they must all arrive at the same node --- merge nodes or add \textbf{dummy activities} (dashed arrows, duration 0) to bring the dependencies together.
\item Check the rules: \textbf{one single start node}, \textbf{one single end node}, \textbf{no loop} (you can never return to a previous node), and \textbf{at most one arrow between any pair of nodes} (otherwise insert a dummy).
\item Number the nodes so that every arrow goes from a smaller number to a larger one, and label each arrow with the task letter and its duration.
\end{enumerate}
\end{methode}

\begin{attention}[Common mistakes in PERT networks]
Two errors lose many marks at the GCE:
\begin{itemize}
\item \textbf{Creating a loop}: an arrow that returns to an earlier node. A loop means a task depends on itself --- the project could never start! Always check that arrows only go ``forward''.
\item \textbf{Two arrows between the same pair of nodes}: if tasks B and C both go from node 1 to node 2, they become indistinguishable. The correction is to give one of them its own end node and join the nodes with a \textbf{dummy activity} (dashed arrow, duration 0).
\end{itemize}
Also remember: a dummy activity is \emph{not} real work; never give it a duration other than 0.
\end{attention}

\begin{exoresolu}[PERT network of the laboratory project]
\textbf{Statement.} The precedence table of the laboratory project is: A (site preparation, 2 weeks, no predecessor); B (wiring, 2 weeks, after A); C (painting, 2 weeks, after A); D (purchase of computers, 5 weeks, no predecessor); E (network, 2 weeks, after B and C); F (hardware setup, 2 weeks, after D and E); G (software installation, 2 weeks, after F); H (training, 2 weeks, after G). Draw the PERT network.
\tcblower
\textbf{Solution.} A and D have no predecessor: they leave the start node 0. B and C both follow A: they leave node 1 (end of A). E needs B \emph{and} C, so B and C must both have arrived at the node where E starts: we end B at node 2 and C at node 3, then join node 3 to node 2 with a \textbf{dummy activity} (duration 0); E can now leave node 2, where both B and (via the dummy) C have arrived. F needs D \emph{and} E: we simply end both D and E at the same node 4, so F leaves node 4 --- no extra dummy is needed here, because D and E are allowed to share their end node. Then G and H follow in sequence to the single end node 7. The resulting network is shown in the figure below.
\end{exoresolu}

\begin{popfigure}
\begin{tikzpicture}[>=latex, font=\small]
% node style: circle split
\tikzset{evt/.style={circle, draw=popdark, thick, minimum size=1.05cm, inner sep=0pt, path picture={\draw[popdark, thin] (path picture bounding box.west) -- (path picture bounding box.east);}}}
% nodes
\node[evt, align=center] (n0) at (0,0){\shortstack{0\\[2pt]\footnotesize 0}};
\node[evt, align=center] (n1) at (2.4,1.6){\shortstack{1\\[2pt]\footnotesize 2}};
\node[evt, align=center] (n2) at (5.2,1.6){\shortstack{2\\[2pt]\footnotesize 4}};
\node[evt, align=center] (n3) at (3.6,-1.6){\shortstack{3\\[2pt]\footnotesize 4}};
\node[evt, align=center] (n4) at (7.6,0){\shortstack{4\\[2pt]\footnotesize 6}};
\node[evt, align=center] (n5) at (10,0){\shortstack{5\\[2pt]\footnotesize 8}};
\node[evt, align=center] (n6) at (12.2,0){\shortstack{6\\[2pt]\footnotesize 10}};
\node[evt, fill=popblueL, align=center] (n7) at (14.2,0){\shortstack{7\\[2pt]\footnotesize 12}};
% arrows
\draw[->, thick, popblue] (n0) -- (n1) node[midway, above, sloped, text=popink]{A (2)};
\draw[->, thick, popblue] (n1) -- (n2) node[midway, above, sloped, text=popink]{B (2)};
\draw[->, thick, popblue] (n1) -- (n3) node[midway, above, sloped, text=popink]{C (2)};
\draw[->, thick, poporange] (n0) -- (n4) node[midway, below, sloped, text=popink]{D (5)};
\draw[->, thick, popgreen] (n2) -- (n4) node[midway, above, sloped, text=popink]{E (2)};
\draw[->, thick, popgreen] (n4) -- (n5) node[midway, above, text=popink]{F (2)};
\draw[->, thick, popgreen] (n5) -- (n6) node[midway, above, text=popink]{G (2)};
\draw[->, thick, poppurple] (n6) -- (n7) node[midway, above, text=popink]{H (2)};
% dummies
\draw[->, thick, dashed, popmuted] (n3) -- (n2) node[midway, right, text=popmuted]{dummy (0)};
\end{tikzpicture}
\legende{PERT network of the laboratory project. Circles are events (number above, earliest date below, in weeks); solid arrows are activities with durations; the dashed arrow is a dummy activity forcing E to wait for both B and C.}
\end{popfigure}

Note how the dummy activity from node 3 to node 2 expresses ``E waits for C as well as for B'' without inventing any false work. The numbers inside the nodes are the \textbf{earliest dates} of the events, in weeks: node 4, for example, can only occur at week $6$, because it must wait for both D (finished at week $5$) and E (finished at week $4+2=6$), and we take the \emph{maximum}. Reading the last node, the whole project needs at least \textbf{12 weeks} --- exactly what the Gantt chart showed. These dates will be computed rigorously in the next section, when we study the critical path.

\courssub{Gantt or PERT: Which Tool for Which Job?}

\begin{aretenir}[The essential difference]
The \cle{Gantt chart} shows \textbf{WHEN} tasks happen (calendar, durations, parallel work, progress). The \cle{PERT chart} shows \textbf{HOW} tasks depend on each other (precedence, convergence, and --- as we shall see --- which tasks are critical). A good project manager uses \textbf{both}: PERT to reason about dependencies and criticality, Gantt to communicate the calendar to the team and to sponsors.
\end{aretenir}

\begin{center}
\begin{tabularx}{\linewidth}{|l|X|X|}
\hline
\rowcolor{popblueL} \textbf{Criterion} & \textbf{Gantt chart} & \textbf{PERT chart} \\
\hline
Main question answered & When does each task run? & What depends on what? \\
\hline
Visual form & Bars on a calendar & Network of nodes and arrows \\
\hline
Dependencies & Shown weakly (small arrows) & Shown rigorously (structure itself) \\
\hline
Critical tasks & Not visible & Computable (next section) \\
\hline
Ease of reading & Immediate, for anyone & Requires training \\
\hline
Best use & Communication, tracking progress & Planning, analysing delays \\
\hline
\end{tabularx}
\end{center}

\begin{exemplebox}[Choosing the tool in Yaound\'e]
The committee building a computer laboratory at a Government High School presents the project to the Parent-Teacher Association. For the meeting, the coordinator shows a \textbf{Gantt chart}: parents immediately see that the laboratory will be ready at the end of week 12. But when the supplier of computers announces a two-week delay, the coordinator returns to the \textbf{PERT network} to check whether this delay blocks the whole project or can be absorbed.
\end{exemplebox}

\begin{exercice}[From precedence table to both charts]
A school website project has tasks: A (collect content, 3 weeks, no predecessor); B (design layout, 2 weeks, no predecessor); C (write pages, 2 weeks, after A and B); D (test site, 1 week, after C); E (put online, 1 week, after D).
\begin{enumerate}
\item Draw the Gantt chart of this project on a weekly timeline, with a milestone ``site online''.
\item Draw the PERT network, using a dummy activity if necessary.
\item State, for each tool, one thing it shows better than the other.
\end{enumerate}
\end{exercice}

\begin{corrige}[Exercise 1]
\textbf{1.} Gantt: bars A from W1 to W3; B from W1 to W2; C from W4 to W5 (it waits for A, which ends in W3); D in W6; E in W7; diamond ``site online'' at the end of W7.\\
\textbf{2.} PERT: start node 0; arrows A and B leave node 0. Give A its own end node 1 and B its own end node 2, then join node 2 to node 1 with a dummy activity (duration 0) so that C starts at node 1 where both A and B have arrived. Then C (node 1 $\to$ node 3), D (node 3 $\to$ node 4), E (node 4 $\to$ node 5, the single end node).\\
\textbf{3.} The Gantt chart shows better \emph{when} each task happens on the calendar and which tasks overlap; the PERT network shows better the \emph{dependency structure} (that C needs both A and B) and will allow us to compute the critical path.
\end{corrige}

\begin{savaistu}[PERT and the Polaris submarines]
The PERT technique was developed in 1958 by the United States Navy to manage the Polaris nuclear submarine programme, which involved thousands of contractors. By modelling dependencies explicitly, the programme was completed ahead of schedule. Today the same mathematics helps Cameroonian engineers plan fibre-optic deployments for CAMTEL and MTN: before digging a single trench, the whole network of tasks is simulated on a PERT chart.
\end{savaistu}

We now know how to place tasks in time (Gantt) and how to model their dependencies (PERT). But one crucial question remains: \emph{which tasks can afford a delay, and which cannot?} That is the subject of the next section: project control, the critical path and slack analysis.

\coursec{Project Control, Critical Path and Slack Analysis}

In the previous section we built a project schedule using Gantt charts and PERT networks. We identified activities, estimated their durations and sequenced them with logical dependencies. A schedule, however, is only a plan. Once execution begins, reality intervenes: some tasks finish early, others overrun, resources become unavailable, and new requirements appear. \textbf{Project control} is the discipline of comparing actual progress against the plan, understanding the impact of variances, and taking corrective action to keep the project on track. The mathematical backbone of this control is the \cle{Critical Path Method (CPM)}, which we now develop in full rigour.

\courssub{From Scheduling to Control: The Need for Critical Path Analysis}

A PERT network shows the logical sequence of activities but does not, by itself, tell us which activities are \emph{time-critical} and which have \emph{breathing room}. Consider the computer laboratory project: if the procurement of computers (Activity D) is delayed by one week, does the whole project finish one week later? What if the electrical wiring (Activity B) is delayed by one week? The answer depends on whether the activity lies on the \cle{critical path}. The critical path is the longest path through the network; it determines the \emph{minimum possible project duration}. Any delay on a critical activity directly extends the project end date. Conversely, non-critical activities possess \cle{slack} (or \cle{float}) — a margin of delay that can be absorbed without affecting the project completion.

To quantify slack and identify the critical path we perform two systematic passes over the network:
\begin{enumerate}
    \item \textbf{Forward Pass} — computes the earliest possible start and finish times (\cle{ES}, \cle{EF}) for every activity, respecting predecessor constraints.
    \item \textbf{Backward Pass} — computes the latest allowable start and finish times (\cle{LS}, \cle{LF}) that will not delay the project, working backwards from the project deadline.
\end{enumerate}
The difference \cle{LS - ES} (or \cle{LF - EF}) is the \cle{total slack}. Activities with zero total slack form the critical path.

\courssub{Forward Pass: Earliest Start and Earliest Finish}

The forward pass answers: \emph{``If everything goes as early as possible, when can each activity start and finish?''} We start at the project start node (time 0) and move forward through the network, respecting the rule that an activity cannot start until \emph{all} its predecessors have finished.

\begin{methode}[Forward Pass Algorithm]
\begin{enumerate}
    \item Set the \cle{Earliest Start (ES)} of the first activity (or activities with no predecessors) to \textbf{0}.
    \item For each activity in topological order (i.e. after all its predecessors have been processed):
          \[ \text{ES} = \max\{\text{EF of all immediate predecessors}\} \]
          (If an activity has no predecessors, its ES = 0.)
    \item Compute \cle{Earliest Finish (EF)}:
          \[ \text{EF} = \text{ES} + \text{Duration} \]
    \item The \cle{Earliest Finish of the project} is the maximum EF among all activities that have no successors (terminal activities). This value is the \textbf{minimum project duration}.
\end{enumerate}
\end{methode}

\begin{attention}[Common Mistake: Minimum instead of Maximum]
During the forward pass, a frequent error is to take the \textbf{minimum} of predecessors' EF values. Remember: an activity must wait for \emph{the last} predecessor to finish. Hence we take the \textbf{maximum}. Using the minimum would allow an activity to start before some predecessors are done, violating the dependency logic.
\end{attention}

\courssub{Backward Pass: Latest Start and Latest Finish}

The backward pass answers: \emph{``How late can each activity start and finish without pushing the project end date beyond the minimum duration found in the forward pass?''} We start at the project end (the earliest project finish time) and move backwards.

\begin{methode}[Backward Pass Algorithm]
\begin{enumerate}
    \item Set the \cle{Latest Finish (LF)} of every terminal activity (no successors) to the \textbf{project duration} (the maximum EF from the forward pass).
    \item For each activity in reverse topological order (i.e. before all its successors have been processed):
          \[ \text{LF} = \min\{\text{LS of all immediate successors}\} \]
          (If an activity has no successors, its LF = project duration.)
    \item Compute \cle{Latest Start (LS)}:
          \[ \text{LS} = \text{LF} - \text{Duration} \]
    \item The \cle{Latest Start of the first activity} should be 0 (or the project start date). If it is negative, the project is already behind schedule.
\end{enumerate}
\end{methode}

\begin{attention}[Common Mistake: Maximum instead of Minimum]
In the backward pass, the trap is to take the \textbf{maximum} of successors' LS values. The logic is: an activity must finish early enough so that \emph{the earliest-starting successor} can begin on time. Hence we take the \textbf{minimum} of the successors' LS. Using the maximum would allow the activity to finish after a successor has already needed to start.
\end{attention}

\courssub{Slack (Float): Total and Free}

With ES, EF, LS, LF known for every activity, we compute slack.

\begin{definition}[Total Slack (Total Float)]
The \cle{Total Slack} of an activity is the amount of time it can be delayed \emph{without delaying the project completion date}, assuming all other activities follow their planned schedule.
\[ \text{Total Slack} = \text{LS} - \text{ES} = \text{LF} - \text{EF} \]
\end{definition}

\begin{definition}[Free Slack (Free Float)]
The \cle{Free Slack} of an activity is the amount of time it can be delayed \emph{without delaying the earliest start of any immediate successor}, assuming all predecessors finish at their earliest times.
\[ \text{Free Slack} = \min\{\text{ES of immediate successors}\} - \text{EF} \]
(If an activity has no successors, its free slack equals its total slack.)
\end{definition}

\begin{propriete}[Relationship between Total Slack and Free Slack]
For any activity:
\[ \text{Free Slack} \le \text{Total Slack} \]
Free slack is the portion of total slack that does not consume the slack of successor activities. Total slack is shared along a path; if an activity uses part of its total slack, the total slack available to its successors is reduced by the same amount.
\end{propriete}

\begin{aretenir}[Interpretation of Slack for the Manager]
\begin{itemize}
    \item \textbf{Zero total slack} $\Rightarrow$ the activity is \textbf{critical}. Any delay, however small, extends the project. These activities demand the closest monitoring, the most reliable resources, and contingency plans.
    \item \textbf{Positive total slack} $\Rightarrow$ the activity is \textbf{non-critical}. It has a buffer. The manager can temporarily reassign resources from non-critical to critical activities, or tolerate minor delays without panic.
    \item \textbf{Free slack} indicates how much an activity can slip \emph{independently} before it forces a successor to slip. It is useful for day-to-day scheduling flexibility.
\end{itemize}
\end{aretenir}

\courssub{The Critical Path: Definition and Properties}

\begin{definition}[Critical Path]
The \cle{Critical Path} is the path (or paths) through the project network from start to finish that has the \textbf{longest total duration}. Equivalently, it is the sequence of activities for which \cle{Total Slack = 0}. The length of the critical path equals the \textbf{minimum project duration}.
\end{definition}

\begin{propriete}[Properties of the Critical Path]
\begin{enumerate}
    \item \textbf{Uniqueness of duration:} All critical paths have the same length (the project duration). There can be multiple critical paths running in parallel.
    \item \textbf{Zero slack:} Every activity on a critical path has total slack = 0 (and free slack = 0).
    \item \textbf{Delay propagation:} Delaying any critical activity by $\Delta t$ increases the project duration by $\Delta t$ (unless another path becomes critical and limits the increase).
    \item \textbf{Non-critical paths:} Any path that is not critical has total slack $> 0$ equal to the difference between the critical path length and that path's length.
    \item \textbf{Changes in critical path:} As the project progresses and actual durations replace estimates, the critical path can shift. A non-critical path can become critical if its activities overrun their slack.
\end{enumerate}
\end{propriete}

\begin{savaistu}[Critical Path in Cameroon's ICT Projects]
In large Cameroonian e-government projects (e.g., the \emph{Projet d'Appui \`a la Modernisation de l'Administration} -- PAMA), the critical path often runs through procurement (import licences, customs clearance at Douala port) and civil works (building data centres). A one-week delay in customs can slip the whole project because those activities have zero slack. Project managers at ANTIC and MINPOSTEL routinely use CPM to negotiate priority lanes for critical equipment.
\end{savaistu}

\courssub{Worked Example: Computer Laboratory Project -- Full CPM Computation}

We return to the computer laboratory setup project introduced in the scheduling section. The activity list, durations (in weeks) and dependencies are:

\begin{center}
\begin{tabularx}{\linewidth}{|c|l|c|X|}\hline
\textbf{Code} & \textbf{Description} & \textbf{Dur.} & \textbf{Predecessors} \\ \hline
A & Site preparation \& civil works & 2 & -- \\
B & Electrical wiring \& trunking & 3 & A \\
C & Structured network cabling & 2 & A \\
D & Procurement of computers \& equipment & 4 & A \\
E & Installation of computers \& peripherals & 2 & B, C, D \\
F & OS \& software deployment & 3 & E \\
G & System testing \& commissioning & 1 & F \\
H & Staff training \& handover & 1 & G \\
I & Project documentation \& closure & 1 & G \\ \hline
\end{tabularx}
\end{center}

We will compute ES, EF, LS, LF, Total Slack, Free Slack and identify the critical path.

\begin{exoresolu}[Full CPM Computation for the Computer Laboratory Project]
\textbf{Step 1: Forward Pass (Earliest Times)}

\begin{itemize}
    \item \textbf{A:} ES = 0, EF = 0 + 2 = \textbf{2}
    \item \textbf{B:} Predecessor A (EF=2) $\Rightarrow$ ES = 2, EF = 2 + 3 = \textbf{5}
    \item \textbf{C:} Predecessor A (EF=2) $\Rightarrow$ ES = 2, EF = 2 + 2 = \textbf{4}
    \item \textbf{D:} Predecessor A (EF=2) $\Rightarrow$ ES = 2, EF = 2 + 4 = \textbf{6}
    \item \textbf{E:} Predecessors B(EF=5), C(EF=4), D(EF=6) $\Rightarrow$ ES = \textbf{max(5,4,6)} = 6, EF = 6 + 2 = \textbf{8}
    \item \textbf{F:} Predecessor E (EF=8) $\Rightarrow$ ES = 8, EF = 8 + 3 = \textbf{11}
    \item \textbf{G:} Predecessor F (EF=11) $\Rightarrow$ ES = 11, EF = 11 + 1 = \textbf{12}
    \item \textbf{H:} Predecessor G (EF=12) $\Rightarrow$ ES = 12, EF = 12 + 1 = \textbf{13}
    \item \textbf{I:} Predecessor G (EF=12) $\Rightarrow$ ES = 12, EF = 12 + 1 = \textbf{13}
\end{itemize}
\textbf{Project Duration} = max(EF of terminal activities H, I) = \textbf{13 weeks}.

\textbf{Step 2: Backward Pass (Latest Times)} -- Project end = 13.

\begin{itemize}
    \item \textbf{H:} LF = 13, LS = 13 - 1 = \textbf{12}
    \item \textbf{I:} LF = 13, LS = 13 - 1 = \textbf{12}
    \item \textbf{G:} Successors H(LS=12), I(LS=12) $\Rightarrow$ LF = \textbf{min(12,12)} = 12, LS = 12 - 1 = \textbf{11}
    \item \textbf{F:} Successor G (LS=11) $\Rightarrow$ LF = 11, LS = 11 - 3 = \textbf{8}
    \item \textbf{E:} Successor F (LS=8) $\Rightarrow$ LF = 8, LS = 8 - 2 = \textbf{6}
    \item \textbf{D:} Successor E (LS=6) $\Rightarrow$ LF = 6, LS = 6 - 4 = \textbf{2}
    \item \textbf{C:} Successor E (LS=6) $\Rightarrow$ LF = 6, LS = 6 - 2 = \textbf{4}
    \item \textbf{B:} Successor E (LS=6) $\Rightarrow$ LF = 6, LS = 6 - 3 = \textbf{3}
    \item \textbf{A:} Successors B(LS=3), C(LS=4), D(LS=2) $\Rightarrow$ LF = \textbf{min(3,4,2)} = 2, LS = 2 - 2 = \textbf{0}
\end{itemize}

\textbf{Step 3: Slack Computation and Critical Path Identification}

\begin{center}
\begin{tabular}{|c|c|c|c|c|c|c|c|c|}\hline
\rowcolor{popblueL}
\textbf{Act} & \textbf{Dur} & \textbf{ES} & \textbf{EF} & \textbf{LS} & \textbf{LF} & \textbf{Total Slack} & \textbf{Free Slack} & \textbf{Critical?} \\ \hline
A & 2 & 0 & 2 & 0 & 2 & 0 & 0 & \textbf{Yes} \\
B & 3 & 2 & 5 & 3 & 6 & 1 & 1 & No \\
C & 2 & 2 & 4 & 4 & 6 & 2 & 2 & No \\
D & 4 & 2 & 6 & 2 & 6 & 0 & 0 & \textbf{Yes} \\
E & 2 & 6 & 8 & 6 & 8 & 0 & 0 & \textbf{Yes} \\
F & 3 & 8 & 11 & 8 & 11 & 0 & 0 & \textbf{Yes} \\
G & 1 & 11 & 12 & 11 & 12 & 0 & 0 & \textbf{Yes} \\
H & 1 & 12 & 13 & 12 & 13 & 0 & 0 & \textbf{Yes} \\
I & 1 & 12 & 13 & 12 & 13 & 0 & 0 & \textbf{Yes} \\ \hline
\end{tabular}
\end{center}

\textbf{Critical Path(s):} There are two critical paths (both length 13 weeks):
\begin{enumerate}
    \item A $\to$ D $\to$ E $\to$ F $\to$ G $\to$ H
    \item A $\to$ D $\to$ E $\to$ F $\to$ G $\to$ I
\end{enumerate}
Activities B and C have positive slack (1 and 2 weeks respectively). They can be delayed within those limits without affecting the 13-week project completion.

\tcblower
\textbf{Detailed Solution Commentary:}
\begin{itemize}
    \item The forward pass correctly takes the \emph{maximum} EF at each merge (Activity E). The backward pass correctly takes the \emph{minimum} LS at each burst (Activity G and Activity A).
    \item Notice that Activities H and I are parallel terminal activities. Both have zero total slack because the project cannot finish until \emph{both} are done. If only H were required, I would have free slack = total slack = 1 week (it could finish at week 14 without delaying the project). But since both are deliverables, both are critical.
    \item The critical path runs through the \emph{longest} procurement leg (D = 4 weeks) rather than the electrical (B = 3) or cabling (C = 2) legs. This matches intuition: procurement is the pacing activity.
    \item Managerial implication: The procurement officer (Activity D) must be monitored daily. The electrician (B) and cable installer (C) have some flexibility; if the electrician is sick for a week, the project still finishes on time.
\end{itemize}
\end{exoresolu}

\begin{popfigure}
\centering
\begin{tikzpicture}[
    node style/.style={draw, rounded corners, minimum width=2.5cm, minimum height=1.2cm, align=center, font=\small},
    critical/.style={node style, fill=poppink!30, thick, draw=popink},
    normal/.style={node style, fill=popblueL!30},
    arrow/.style={->, >=stealth, thick},
    critical arrow/.style={arrow, popink, very thick},
    slack label/.style={font=\tiny, text=popmuted, anchor=north}
]
% Nodes
\node[critical, align=center] (A) at (0,0){A\\Site prep\\Dur=2\\ES=0 EF=2\\LS=0 LF=2\\Slack=0};
\node[normal, align=center] (B) at (3,2){B\\Electrical\\Dur=3\\ES=2 EF=5\\LS=3 LF=6\\Slack=1};
\node[normal, align=center] (C) at (3,0){C\\Cabling\\Dur=2\\ES=2 EF=4\\LS=4 LF=6\\Slack=2};
\node[critical, align=center] (D) at (3,-2){D\\Procurement\\Dur=4\\ES=2 EF=6\\LS=2 LF=6\\Slack=0};
\node[critical, align=center] (E) at (6,0){E\\Install PCs\\Dur=2\\ES=6 EF=8\\LS=6 LF=8\\Slack=0};
\node[critical, align=center] (F) at (9,0){F\\Software\\Dur=3\\ES=8 EF=11\\LS=8 LF=11\\Slack=0};
\node[critical, align=center] (G) at (12,0){G\\Testing\\Dur=1\\ES=11 EF=12\\LS=11 LF=12\\Slack=0};
\node[critical, align=center] (H) at (15,1){H\\Training\\Dur=1\\ES=12 EF=13\\LS=12 LF=13\\Slack=0};
\node[critical, align=center] (I) at (15,-1){I\\Doc/Closure\\Dur=1\\ES=12 EF=13\\LS=12 LF=13\\Slack=0};

% Arrows - critical path in bold red
\draw[critical arrow] (A.east) -- (D.west);
\draw[critical arrow] (D.east) -- (E.west);
\draw[critical arrow] (E.east) -- (F.west);
\draw[critical arrow] (F.east) -- (G.west);
\draw[critical arrow] (G.north east) -- (H.west);
\draw[critical arrow] (G.south east) -- (I.west);

% Non-critical arrows
\draw[arrow] (A.north east) -- (B.south west);
\draw[arrow] (A.east) -- (C.west);
\draw[arrow] (B.south east) -- (E.north west);
\draw[arrow] (C.east) -- (E.west);

% Slack labels on non-critical arrows
\node[slack label] at (1.5,1.2) {Slack=1};
\node[slack label] at (1.5,-0.2) {Slack=2};

% Start/End markers
\node[draw, circle, minimum size=0.6cm] (start) at (-1.5,0) {};
\node[draw, circle, minimum size=0.6cm] (end) at (16.5,0) {};
\draw[arrow] (start) -- (A.west);
\draw[arrow] (H.east) -- (end);
\draw[arrow] (I.east) -- (end);
\end{tikzpicture}
\legende{Activity-on-Node PERT network for the Computer Laboratory Project. Critical path activities (A, D, E, F, G, H, I) are highlighted in pink with bold red arrows. Non-critical activities B and C show their slack values on the incoming arrows.}
\end{popfigure}

\courssub{Project Monitoring and Control}

Having established the baseline schedule (the critical path and slack values), the project enters the \textbf{execution and control} phase. Control is a continuous cycle: \textbf{Measure $\to$ Compare $\to$ Decide $\to$ Act}.

\begin{aretenir}[Key Control Activities]
\begin{enumerate}
    \item \textbf{Progress Tracking:} Collect actual start/finish dates and \% complete for each activity at regular intervals (weekly for a 13-week project).
    \item \textbf{Variance Analysis:} Compare \emph{Planned} vs \emph{Actual}:
          \begin{itemize}
              \item \textbf{Schedule Variance (SV):} Is the activity ahead or behind its planned dates? (Simple version: Actual Start vs ES, Actual Finish vs EF).
              \item \textbf{Cost Variance (CV):} Actual Cost vs Budgeted Cost (if cost data is tracked).
          \end{itemize}
    \item \textbf{Milestone Reviews:} Formal checkpoints at major deliverables (e.g., ``Computers installed'', ``Software deployed''). A milestone is a zero-duration event marking the completion of a key phase.
    \item \textbf{Progress Reports:} Weekly status reports showing: activities completed this week, activities in progress, upcoming activities, issues/risks, and decisions needed.
\end{enumerate}
\end{aretenir}

\begin{definition}[Elementary Earned Value Concepts]
While full Earned Value Management (EVM) is beyond the GCE AL scope, the \emph{intuition} is examinable:
\begin{itemize}
    \item \textbf{Planned Value (PV):} Budgeted cost of work scheduled to be done by now.
    \item \textbf{Earned Value (EV):} Budgeted cost of work \emph{actually} completed (i.e., \% complete $\times$ total budget for the activity).
    \item \textbf{Actual Cost (AC):} Money actually spent.
    \item \textbf{Schedule Performance Index (SPI) = EV / PV:} $>1$ = ahead of schedule; $<1$ = behind.
    \item \textbf{Cost Performance Index (CPI) = EV / AC:} $>1$ = under budget; $<1$ = over budget.
\end{itemize}
At Upper Sixth level, you may be asked to \emph{interpret} given SPI/CPI values or compute simple \% complete vs budget spent.
\end{definition}

\begin{exemplebox}[Interpreting Simple Earned Value Data]
At week 6 of the computer lab project (planned duration 13 weeks), the project manager collects:
\begin{itemize}
    \item Activity D (Procurement, budget 2\,000\,000 FCFA): 100\% complete, Actual Cost = 2\,100\,000 FCFA.
    \item Activity E (Installation, budget 500\,000 FCFA): 50\% complete, Actual Cost = 300\,000 FCFA.
    \item All other activities: not started or on track.
\end{itemize}
\textbf{Analysis:}
\begin{itemize}
    \item D: EV = 2\,000\,000; AC = 2\,100\,000 $\Rightarrow$ CPI = 0.95 (5\% over budget). SV = 0 (finished on time).
    \item E: At week 8 (when E should be 100\% complete): EV = 250\,000; PV = 500\,000 $\Rightarrow$ SPI = 0.5 (behind schedule).
    \item \textbf{Action:} Investigate procurement overrun (maybe customs duties). For installation, add a technician to catch up (crashing).
\end{itemize}
\end{exemplebox}

\courssub{Corrective Actions, Change Management and Project Closure}

When variances are detected, the project manager must act. The response depends on the severity and the activity's criticality.

\begin{methode}[Corrective Action Decision Framework]
\begin{enumerate}
    \item \textbf{Assess Impact:} Is the delayed activity on the critical path? How much slack is consumed?
    \item \textbf{If Non-Critical with Slack Remaining:} Absorb the delay; monitor. No immediate action needed.
    \item \textbf{If Critical or Slack Exhausted:} Choose a recovery strategy:
          \begin{itemize}
              \item \textbf{Crashing:} Add resources (overtime, extra staff, express delivery) to critical activities to shorten their duration. Increases cost.
              \item \textbf{Fast-Tracking:} Perform activities in parallel that were originally sequential (e.g., start software deployment on first batch of computers while the rest are still being installed). Increases risk (rework).
              \item \textbf{Scope Reduction:} Negotiate with sponsor to drop non-essential features (e.g., defer advanced training module).
              \item \textbf{Schedule Extension:} Request a new deadline from the sponsor (last resort).
          \end{itemize}
    \item \textbf{Document the Change:} Every deviation from the baseline schedule/budget must go through a \cle{Change Control Process}: Request $\to$ Impact Analysis $\to$ Approval (by Project Board/Sponsor) $\to$ Update Baseline $\to$ Communicate.
    \item \textbf{Update the Schedule:} Recompute ES/EF/LS/LF with actual dates and remaining durations. The critical path may shift.
\end{enumerate}
\end{methode}

\begin{attention}[Crashing vs Fast-Tracking]
\begin{itemize}
    \item \textbf{Crashing} only works on critical activities and assumes the activity is \emph{resource-elastic} (adding people actually speeds it up). Brooks' Law: ``Adding manpower to a late software project makes it later.'' Crashing is effective for procurement (pay express shipping) or civil works (double shifts), less so for complex configuration.
    \item \textbf{Fast-Tracking} increases overlap risk. In our example, starting software deployment (F) before \emph{all} computers are installed (E) is fast-tracking. If the last computers have different hardware, the software image might fail $\to$ rework.
\end{itemize}
\end{attention}

\begin{aretenir}[Project Closure Checklist]
A project is not finished when the last technical task ends. Formal closure includes:
\begin{enumerate}
    \item \textbf{User Acceptance Testing (UAT) \& Sign-off:} The client (school principal, MINSEC representative) verifies the lab meets specifications and signs the acceptance certificate.
    \item \textbf{Documentation Handover:} As-built network diagrams, equipment inventories, software licences, admin passwords, maintenance contracts (MTN/Orange business support), warranty certificates.
    \item \textbf{Resource Release:} Contractors paid off, staff reassigned, rented equipment returned.
    \item \textbf{Financial Closure:} Final invoice, budget reconciliation, audit if required.
    \item \textbf{Lessons Learned Workshop:} What went well? What failed? (e.g., ``Procurement took 4 weeks instead of 3 because customs clearance at Douala was slower than expected. Next time, engage a clearing agent earlier.'') Store in organisational process assets for future projects.
    \item \textbf{Celebrate!} Recognise the team.
\end{enumerate}
\end{aretenir}

\courssub{Integration: Using CPM for Day-to-Day Decisions}

The critical path method is not a one-off calculation at the start. It is a \emph{living model}. Each week the project manager:
\begin{enumerate}
    \item Updates the network with \emph{actual} start/finish dates for completed activities.
    \item Re-estimates \emph{remaining} durations for in-progress and future activities.
    \item Re-runs the forward/backward pass (software does this instantly).
    \item Checks: Has the critical path changed? Has any slack been consumed? Are new risks appearing?
    \item Reports to stakeholders: ``We are on track to finish Week 13. Procurement (D) is critical -- customs cleared yesterday. Electrical (B) used 1 day of its 1-week slack due to cable shortage; we ordered more. No other issues.''
\end{enumerate}

\begin{savaistu}[Open Source Tools for CPM in Cameroon]
Students can practice CPM calculations with free tools:
\begin{itemize}
    \item \textbf{ProjectLibre} (open-source Microsoft Project alternative) -- runs on Windows/Linux/Mac.
    \item \textbf{GanttProject} -- lightweight, good for learning.
    \item \textbf{Python} with \texttt{networkx} or \texttt{pm4py} libraries for algorithmic implementation (excellent for Computer Science students).
    \item Spreadsheet (Excel/LibreOffice Calc): Build the forward/backward pass formulas manually -- best for understanding the mechanics.
\end{itemize}
\end{savaistu}

\begin{exercice}[Practice: Critical Path and Slack Analysis]
A mini-project to deploy a school website has the following activities (duration in days):

\begin{center}
\begin{tabular}{|c|l|c|lc|}\hline
\textbf{ID} & \textbf{Activity} & \textbf{Dur.} & \textbf{Predecessors} \\ \hline
A & Requirements gathering & 3 & -- \\
B & Design mockups & 4 & A \\
C & Set up hosting \& domain & 2 & A \\
D & Front-end development & 6 & B \\
E & Back-end development & 5 & B \\
F & Database setup & 3 & C \\
G & Integration \& testing & 4 & D, E, F \\
H & Content population & 2 & G \\
I & Go-live \& handover & 1 & H \\ \hline
\end{tabular}
\end{center}

\begin{enumerate}
    \item Draw the Activity-on-Node network diagram.
    \item Perform the Forward Pass and Backward Pass. Compute ES, EF, LS, LF, Total Slack, Free Slack for each activity.
    \item Identify the critical path(s) and the minimum project duration.
    \item If Activity C (Hosting) is delayed by 2 days due to a payment issue with the registrar, what is the impact on the project end date? Justify using slack values.
    \item The sponsor demands the project be finished 3 days earlier. Which activities would you consider crashing? Explain why.
\end{enumerate}
\end{exercice}

\begin{corrige}[Exercise: School Website Deployment]
\textbf{1. Network Diagram:} (Draw A $\to$ B, C; B $\to$ D, E; C $\to$ F; D,E,F $\to$ G; G $\to$ H; H $\to$ I)

\textbf{2. Forward/Backward Pass Table:}

\begin{center}
\begin{tabular}{|c|c|c|c|c|c|c|c|c|}\hline
\rowcolor{popblueL}
Act & Dur & ES & EF & LS & LF & Tot.Slack & Free Slack & Critical \\ \hline
A & 3 & 0 & 3 & 0 & 3 & 0 & 0 & Yes \\
B & 4 & 3 & 7 & 3 & 7 & 0 & 0 & Yes \\
C & 2 & 3 & 5 & 8 & 10 & 5 & 0 & No \\
D & 6 & 7 & 13 & 7 & 13 & 0 & 0 & Yes \\
E & 5 & 7 & 12 & 8 & 13 & 1 & 1 & No \\
F & 3 & 5 & 8 & 10 & 13 & 5 & 5 & No \\
G & 4 & 13 & 17 & 13 & 17 & 0 & 0 & Yes \\
H & 2 & 17 & 19 & 17 & 19 & 0 & 0 & Yes \\
I & 1 & 19 & 20 & 19 & 20 & 0 & 0 & Yes \\ \hline
\end{tabular}
\end{center}

\textbf{3. Critical Path:} A $\to$ B $\to$ D $\to$ G $\to$ H $\to$ I (Duration = 20 days).
(Note: Path A-B-E-G-H-I = 19 days, slack 1 on E. Path A-C-F-G-H-I = 15 days, slack 5 on C and F.)

\textbf{4. Impact of 2-day delay on C:} Activity C has Total Slack = 5 days. A 2-day delay consumes 2 days of slack, leaving 3 days slack. \textbf{Project end date unchanged} (still 20 days). The new EF of C = 7, LF remains 10, LS becomes 8. Path A-C-F-G-H-I becomes 17 days, still 3 days shorter than critical path.

\textbf{5. Crashing to save 3 days:} Must crash activities on the critical path: A, B, D, G, H, I.
\begin{itemize}
    \item A (Requirements): Hard to crash (stakeholder availability).
    \item B (Design): Could add a designer (risk of inconsistency).
    \item D (Front-end, 6 days): Best candidate -- add a developer, use component library, work weekends. Can likely save 2-3 days.
    \item G (Testing, 4 days): Can run test automation, parallel testers -- save 1-2 days.
    \item H (Content, 2 days): Add content editors -- save 1 day.
    \item I (Go-live, 1 day): Cannot crash.
\end{itemize}
Recommended: Crash D by 2 days (add developer) and G by 1 day (automated testing) $\to$ Project duration = 17 days. Cost increase: developer overtime + testing tools.
\end{corrige}

\begin{aretenir}[Summary: Project Control Essentials for GCE AL]
\begin{itemize}
    \item \textbf{Forward Pass:} ES = max(Predecessors' EF); EF = ES + Dur. Project Duration = max(EF).
    \item \textbf{Backward Pass:} LF = min(Successors' LS); LS = LF - Dur. Project End = Project Duration.
    \item \textbf{Total Slack} = LS - ES = LF - EF. \textbf{Free Slack} = min(Successors' ES) - EF.
    \item \textbf{Critical Path} = Zero Total Slack activities. Length = Project Duration.
    \item \textbf{Control:} Track actuals, compute variances, focus on critical path, manage changes formally.
    \item \textbf{Corrective Actions:} Crash (add resources to critical), Fast-track (parallelise), Reduce scope, Extend schedule.
    \item \textbf{Closure:} Acceptance, Documentation, Lessons Learned, Release, Celebrate.
\end{itemize}
\end{aretenir}

\coursec{Integration: Running a Complete Sample Project}

In the previous sections we learned how to \emph{control} a project: computing earliest and latest dates, identifying the \cle{critical path}, measuring \cle{slack} and reacting when a task slips. We now possess every tool of the project manager's toolbox. This final section puts everything together: we shall run a \textbf{complete sample project} from initiation to closure, exactly as a project team would do in a company in Douala or Yaoundé, and exactly as the GCE Advanced Level examiner expects you to do on paper. By the end, you will be able to take any small project brief and produce, in order: a charter, SMART objectives, a WBS, a precedence table, a Gantt chart, a PERT network, a critical path analysis and a control plan.

\courssub{Recap: The Complete Project Management Process}

A project is not managed by improvising; it follows a disciplined life cycle. Everything we have studied in this chapter fits into five phases.

\begin{definition}[The five phases of project management]
The \cle{project management life cycle} is the sequence of phases through which every project passes:
\begin{enumerate}
\item \textbf{Initiation}: define the need, the scope, the stakeholders; write the \cle{project charter}.
\item \textbf{Planning}: decompose the work into a \cle{WBS}, estimate durations and resources, establish the precedence table.
\item \textbf{Scheduling}: build the \cle{Gantt chart} and the \cle{PERT network}; compute the project duration.
\item \textbf{Execution and Control}: monitor progress, track the \cle{critical path} and \cle{slacks}, apply corrective actions.
\item \textbf{Closure}: deliver, evaluate, document \cle{lessons learned}, release the team.
\end{enumerate}
\end{definition}

\begin{popfigure}
\begin{center}
\begin{tikzpicture}[
font=\small,
phase/.style={draw=popblue, fill=popblueL, rounded corners=3pt, minimum height=1.1cm, minimum width=2.6cm, align=center, text=popdark},
arr/.style={-{Stealth[length=2.5mm]}, very thick, poporange},
note/.style={font=\scriptsize\itshape, text=popmuted, align=center}
]
\node[phase, align=center] (ini) at (0,0){\textbf{1. Initiation}\\ charter, scope};
\node[phase, align=center] (pla) at (3.4,0){\textbf{2. Planning}\\ WBS, durations};
\node[phase, align=center] (sch) at (6.8,0){\textbf{3. Scheduling}\\ Gantt, PERT};
\node[phase, align=center] (ctl) at (10.2,0){\textbf{4. Control}\\ critical path, slack};
\node[phase, align=center] (clo) at (13.6,0){\textbf{5. Closure}\\ lessons learned};
\draw[arr] (ini) -- (pla);
\draw[arr] (pla) -- (sch);
\draw[arr] (sch) -- (ctl);
\draw[arr] (ctl) -- (clo);
\draw[arr, dashed, poppurple] (ctl.south) .. controls (8.5,-1.6) and (5.5,-1.6) .. (pla.south)
 node[midway, below, note, text=poppurple]{corrective action: re-plan if a task slips};
\end{tikzpicture}
\end{center}
\legende{The complete project management process, from initiation to closure. Control may loop back to planning when disruptions occur.}
\end{popfigure}

\begin{methode}[The 6-step checklist to manage any small project]
Whatever the project (a school event, a software deployment, a science fair), apply these six steps in order:
\begin{enumerate}
\item \textbf{Initiate}: write the project charter (purpose, scope, budget, deadline, sponsor, team).
\item \textbf{Set objectives}: formulate SMART objectives (Specific, Measurable, Achievable, Relevant, Time-bound).
\item \textbf{Decompose}: build the WBS down to tasks of a few days each; assign a responsible person to each.
\item \textbf{Schedule}: draw the precedence table, then the Gantt chart and the PERT network.
\item \textbf{Analyse}: compute earliest/latest dates, the critical path, the total duration and all slacks.
\item \textbf{Control and close}: fix milestones, write weekly reports, react to disruptions using slack analysis, then close with a lessons-learned session.
\end{enumerate}
\end{methode}

\courssub{Step 1: Project Charter and SMART Objectives}

\begin{definition}[Project charter]
The \cle{project charter} is the short founding document that officially authorises the project. It states the purpose, the scope (what is included and excluded), the budget, the deadline, the sponsor, the project manager and the main success criteria.
\end{definition}

\begin{exemplebox}[Charter of the CS-Week project (extract)]
\begin{center}
\begin{tabularx}{\linewidth}{|l|X|}
\hline
\rowcolor{popblueL} \textbf{Element} & \textbf{Content} \\
\hline
Purpose & Organise a Computer Science Week to promote ICT in the school \\
\hline
Scope (in) & Exhibitions, coding competition, prize-giving, publicity \\
\hline
Scope (out) & Permanent computer-lab renovation (separate project) \\
\hline
Budget & \SI{450000}{FCFA} (sponsor: mobile-money operator) \\
\hline
Deadline & Opening ceremony 8 weeks after kick-off \\
\hline
Sponsor / PM & ICT teacher / elected student PM \\
\hline
Success criteria & Event on time, within budget, at least 200 visitors \\
\hline
\end{tabularx}
\end{center}
\end{exemplebox}

Objectives must be \cle{SMART}. Compare: ``have a good week'' (vague, not examinable) with a SMART version.

\begin{exemplebox}[SMART objectives]
\begin{enumerate}
\item Hold the opening ceremony of CS-Week on the Monday of week 9 (Specific, Time-bound).
\item Register at least 60 contestants for the coding competition by the end of week 6 (Measurable).
\item Keep total expenditure at or below \SI{450000}{FCFA} (Achievable, Measurable).
\item Publish a final report with photos and accounts one week after closure (Relevant, Time-bound).
\end{enumerate}
\end{exemplebox}

\courssub{Step 2: The WBS and the Precedence Table}

The planners decompose the project. Level 1 is the project; level 2 gives the work packages (Initiation, Publicity, Logistics, Competition, Closure); level 3 gives the tasks.

\begin{exemplebox}[WBS and precedence table of the CS-Week project]
Durations are in \textbf{weeks}.
\begin{center}
\begin{tabularx}{\linewidth}{|c|X|c|c|}
\hline
\rowcolor{popblueL} \textbf{Code} & \textbf{Task} & \textbf{Duration} & \textbf{Predecessors} \\
\hline
A & Write charter, get sponsor approval & 1 & --- \\
\hline
B & Design posters and announcements & 2 & A \\
\hline
C & Book hall and equipment & 2 & A \\
\hline
D & Publicity campaign (radio, social media) & 3 & B \\
\hline
E & Prepare competition problems and rules & 3 & A \\
\hline
F & Register contestants & 2 & D \\
\hline
G & Buy prizes and materials & 2 & C \\
\hline
H & Run the CS-Week (exhibitions, contest) & 1 & E, F, G \\
\hline
I & Prize-giving, final report, closure & 1 & H \\
\hline
\end{tabularx}
\end{center}
\end{exemplebox}

\begin{attention}[Common mistake: skipping the precedence table]
Many candidates jump straight to the PERT network and then mix up dependencies. \textbf{Always write the precedence table first}: it is the single source of truth from which both the Gantt chart and the PERT network are drawn. One forgotten predecessor changes the whole critical path.
\end{attention}

\courssub{Step 3: Gantt Chart and PERT Network}

The Gantt chart places each task on a time axis (weeks 1 to 10); the PERT network shows dependencies as a directed graph. Both come from the same precedence table.

\begin{popfigure}
\begin{center}
\begin{tikzpicture}[font=\scriptsize, text=popdark, scale=0.95]
% ---- Mini WBS (left) ----
\begin{scope}
\node[draw=popblue, fill=popblue, text=white, rounded corners=2pt, minimum width=2.2cm, align=center] (root) at (0,0) {\textbf{CS-Week}};
\node[draw=popblue, fill=popblueL, rounded corners=2pt, minimum width=1.55cm, align=center] (w1) at (-2.4,-1.2) {Initiation};
\node[draw=popblue, fill=popblueL, rounded corners=2pt, minimum width=1.55cm, align=center] (w2) at (-0.8,-1.2) {Publicity};
\node[draw=popblue, fill=popblueL, rounded corners=2pt, minimum width=1.55cm, align=center] (w3) at (0.8,-1.2) {Logistics};
\node[draw=popblue, fill=popblueL, rounded corners=2pt, minimum width=1.55cm, align=center] (w4) at (2.4,-1.2) {Event};
\draw[popline] (root) -- (w1); \draw[popline] (root) -- (w2);
\draw[popline] (root) -- (w3); \draw[popline] (root) -- (w4);
\node[font=\tiny, text=popmuted, align=center] at (-2.4,-1.85) {A};
\node[font=\tiny, text=popmuted, align=center] at (-0.8,-1.85) {B, D, F};
\node[font=\tiny, text=popmuted, align=center] at (0.8,-1.85) {C, E, G};
\node[font=\tiny, text=popmuted, align=center] at (2.4,-1.85) {H, I};
\node[font=\small\bfseries, text=popdark] at (0,0.8) {WBS};
\end{scope}
% ---- Gantt (right) ----
\begin{scope}[xshift=7.2cm, yshift=0.55cm]
\node[font=\small\bfseries, text=popdark] at (2.5,0.85) {Gantt chart (weeks)};
\foreach \x in {1,...,10} {
  \draw[popline, very thin] (\x*0.5,0) -- (\x*0.5,-4.0);
  \node[font=\tiny, text=popmuted] at (\x*0.5-0.25,0.15) {\x};
}
% Tasks: y positions: A=0, B=-0.5, C=-1.0, D=-1.5, E=-2.0, F=-2.5, G=-3.0, H=-3.5, I=-4.0
\foreach \y/\name/\s/\d/\col in {
  0/A/0/1/poppurple,
  -0.5/B/1/2/poporange,
  -1.0/C/1/2/poporange,
  -1.5/D/3/3/poppurple,
  -2.0/E/1/3/poporange,
  -2.5/F/6/2/poppurple,
  -3.0/G/3/2/poporange,
  -3.5/H/8/1/poppurple,
  -4.0/I/9/1/poppurple} {
  \node[font=\tiny, anchor=east, text=popdark] at (-0.05,\y-0.18) {\name};
  \fill[\col, rounded corners=1pt] (\s*0.5,\y-0.32) rectangle ({(\s+\d)*0.5},\y-0.04);
}
% Critical tasks highlighted in purple (A, D, F, H, I) - already set via \col
% Non-critical in orange (B, C, E, G)
\end{scope}
\end{tikzpicture}

\vspace{0.4cm}
\begin{tikzpicture}[font=\scriptsize, text=popdark,
evt/.style={circle, draw=popblue, fill=popblueL, minimum size=0.85cm, inner sep=1pt},
task/.style={-{Stealth[length=2mm]}, thick, popdark},
crit/.style={-{Stealth[length=2mm]}, very thick, poporange}]
\node[font=\small\bfseries, text=popdark] at (5.5,1.3) {PERT network (durations in weeks)};
% Nodes (events)
\node[evt] (s) at (0,0) {S};
\node[evt] (n1) at (2,0) {1};
\node[evt] (n2) at (4,1.5) {2};
\node[evt] (n3) at (4,-1.5) {3};
\node[evt] (n4) at (7,1.5) {4};
\node[evt] (n5) at (7,-1.5) {5};
\node[evt] (n6) at (9.5,0) {6};
\node[evt] (n7) at (11.5,0) {7};
\node[evt] (e) at (13.5,0) {E};
% Arrows
\draw[crit] (s) -- node[above, midway]{A (1)} (n1);
\draw[task] (n1) -- node[above, sloped, pos=0.4]{B (2)} (n2);
\draw[crit] (n1) -- node[below, sloped, pos=0.4]{E (3)} (n3);
\draw[task] (n1) -- node[below, sloped, pos=0.4]{C (2)} (n5);
\draw[task] (n2) -- node[above, midway]{D (3)} (n4);
\draw[crit] (n4) -- node[above, sloped, pos=0.4]{F (2)} (n6);
\draw[crit] (n3) -- node[below, sloped, pos=0.4]{dummy (0)} (n6);
\draw[task] (n5) -- node[below, sloped, pos=0.4]{G (2)} (n6);
\draw[crit] (n6) -- node[above, midway]{H (1)} (n7);
\draw[crit] (n7) -- node[above, midway]{I (1)} (e);
\end{tikzpicture}
\end{center}
\legende{Poster view of the sample project: mini WBS (top left), Gantt chart with critical tasks in purple and non-critical in orange (top right), and PERT network with the critical path in orange (bottom). The network includes a dummy activity (duration 0) to correctly model the convergence of E, F, G onto H.}
\end{popfigure}

\courssub{Step 4: Critical Path, Duration and Slacks}

The controllers now compute. Working \emph{forwards} (earliest start of a task = maximum of the earliest finishes of its predecessors) and \emph{backwards} (latest finish = minimum of the latest starts of its successors), we obtain the following table. Recall: $\text{slack} = \text{latest start} - \text{earliest start}$.

\begin{exoresolu}[Critical path and slacks of the CS-Week project]
Using the precedence table above (durations in weeks), compute the earliest and latest schedules, the slack of every task, the project duration and the critical path.
\tcblower
\textbf{Forward pass (earliest start ES, earliest finish EF):}
\begin{align*}
A &: ES=0,\ EF=1\\
B &: ES=1,\ EF=3 \qquad C: ES=1,\ EF=3 \qquad E: ES=1,\ EF=4\\
D &: ES=3,\ EF=6 \qquad G: ES=3,\ EF=5\\
F &: ES=6,\ EF=8\\
H &: ES=\max(4,8,5)=8,\ EF=9\\
I &: ES=9,\ EF=10
\end{align*}
\textbf{Backward pass (latest finish LF, latest start LS)}, with project end at week 10:
\begin{align*}
I &: LF=10,\ LS=9 \qquad H: LF=9,\ LS=8\\
F &: LF=8,\ LS=6 \qquad E: LF=8,\ LS=5 \qquad G: LF=8,\ LS=6\\
D &: LF=6,\ LS=3 \qquad B: LF=3,\ LS=1 \qquad C: LF=6,\ LS=4\\
A &: LF=\min(1,4,5)=1,\ LS=0
\end{align*}

\textbf{Slack table:}
\begin{center}
\begin{tabular}{|c|c|c|c|c|c|}
\hline
\rowcolor{popblueL} Task & ES & LS & Duration & Slack & Critical? \\
\hline
A & 0 & 0 & 1 & 0 & Yes \\
B & 1 & 1 & 2 & 0 & Yes \\
C & 1 & 4 & 2 & 3 & No \\
D & 3 & 3 & 3 & 0 & Yes \\
E & 1 & 5 & 3 & 4 & No \\
F & 6 & 6 & 2 & 0 & Yes \\
G & 3 & 6 & 2 & 3 & No \\
H & 8 & 8 & 1 & 0 & Yes \\
I & 9 & 9 & 1 & 0 & Yes \\
\hline
\end{tabular}
\end{center}
\textbf{Conclusion:} the \textbf{project duration is 10 weeks} and the \textbf{critical path is}
\[ A \rightarrow B \rightarrow D \rightarrow F \rightarrow H \rightarrow I. \]
\end{exoresolu}

\textbf{A management problem!} Our schedule needs 10 weeks, but the charter allows only 8. This is exactly the kind of finding a good team must report. Options: reduce publicity $D$ from 3 to 2 weeks (paid radio spots instead of free ones), allow registration $F$ to overlap with the last week of $D$ (fast-tracking), and run $H$ and the report part of $I$ in parallel. After compression, the critical path becomes $A(1)+B(2)+D(2)+F(1.5)+H(1)+I(0.5)=8$ weeks: the deadline is met, and this negotiation is documented in the control plan.

\begin{attention}[Common exam mistake]
After computing slacks, many candidates stop. The examiner expects you to \textbf{state explicitly}: (1) the \textbf{project duration} (``the project lasts 10 weeks'') and (2) the \textbf{critical path written as a sequence of tasks}. A correct table without these two sentences loses marks.
\end{attention}

\courssub{Step 5: Control Plan and Disruption Simulation}

\begin{definition}[Control plan]
The \cle{control plan} defines how progress is monitored: the \cle{milestones} (key checkpoints with dates), the frequency and format of progress reports, the person responsible for tracking, and the escalation procedure when a deviation is detected.
\end{definition}

\begin{exemplebox}[Control plan of the CS-Week project]
\begin{itemize}
\item \textbf{Milestones}: M1 = charter approved (end of week 1); M2 = publicity launched (end of week 3); M3 = registration closed (end of week 6); M4 = event delivered (week 8).
\item \textbf{Weekly report} (every Friday, one page): tasks done, tasks late, budget spent, risks, decisions needed.
\item \textbf{Rule}: any delay on a \emph{critical} task is escalated to the sponsor within 24 hours.
\end{itemize}
\end{exemplebox}

\begin{experience}[Simulating a disruption: the supplier delivers late]
\textbf{Scenario.} At the end of week 4, the prize supplier announces that task $G$ (buy prizes and materials, planned weeks 4--5) will finish \textbf{2 weeks late}, i.e.\ at the end of week 7 instead of week 5.

\textbf{Analysis using slack.} Task $G$ has a slack of 3 weeks (from our table). A 2-week delay consumes only part of the slack: $G$ now finishes at week 7, still before its latest allowed finish (week 8). The critical path $A \rightarrow B \rightarrow D \rightarrow F \rightarrow H \rightarrow I$ is \emph{unchanged}, and the event date is \textbf{safe}.

\textbf{Corrective actions anyway (good practice):}
\begin{enumerate}
\item record the deviation in the weekly report and recompute $G$'s remaining slack (now 1 week);
\item ask the supplier for a written delivery date and identify a backup supplier in Bafoussam market;
\item monitor $G$ weekly, since it is now \emph{near-critical}.
\end{enumerate}
\textbf{What if the delay were 4 weeks?} Then $G$ would finish at week 9, beyond $H$'s latest start: $G$ would become critical and the project would slip by 1 week. The PM would then compress $H$ or $I$ (e.g.\ prepare report templates in advance) to recover.
\end{experience}

\courssub{Step 6: Closure, Presentation and Peer Evaluation}

\begin{definition}[Project closure and lessons learned]
\cle{Project closure} is the final phase in which the deliverables are formally accepted by the sponsor, resources are released, accounts are settled and a \cle{lessons-learned} document is written. \emph{Lessons learned} record what went well, what went wrong and what the team would do differently, so that future projects benefit from the experience.
\end{definition}

For the CS-Week, closure means: the prize-giving ceremony held, the final report (with accounts and photos) handed to the Principal, the sponsor thanked, and a 30-minute team debrief. Each team then \textbf{presents} its charter, WBS, Gantt, PERT and critical path analysis to the class, and the other teams evaluate it using a peer grid: correctness of dependencies, correctness of the critical path, quality of the control plan, clarity of the presentation.

\begin{aretenir}[Synthesis: WBS vs Gantt vs PERT]
\begin{center}
\begin{tabularx}{\linewidth}{|l|X|X|X|}
\hline
\rowcolor{popblueL} & \textbf{WBS} & \textbf{Gantt chart} & \textbf{PERT network} \\
\hline
\textbf{Purpose} & Decompose the project into manageable tasks & Show tasks on a calendar time axis & Show dependencies and compute dates \\
\hline
\textbf{Input} & Project scope and charter & Task list + durations & Precedence table + durations \\
\hline
\textbf{Output} & Hierarchical task tree (work packages) & Bar chart; who does what and when & Critical path, duration, slacks \\
\hline
\textbf{Strength} & Nothing is forgotten & Easy to read and communicate & Rigorous delay analysis \\
\hline
\textbf{Limit} & No time information & Dependencies hard to see & Less intuitive to read \\
\hline
\end{tabularx}
\end{center}
The three tools are \emph{complementary}: a complete project file contains all three.
\end{aretenir}

\courssub{GCE Exam Technique}

Typical Advanced Level questions on this chapter give you a small project described in words and ask you to:
\begin{enumerate}
\item extract the task list with durations and predecessors (build the precedence table);
\item draw the WBS or the Gantt chart;
\item construct the PERT network;
\item compute earliest/latest dates, slacks, the critical path and the minimum project duration;
\item comment on a scenario (``task X is delayed by 2 days: what is the impact?'').
\end{enumerate}

\begin{methode}[Winning strategy in the exam]
\begin{enumerate}
\item Read the brief twice; underline every duration and every ``after'', ``once'', ``before'' (these give the predecessors).
\item Write the precedence table \emph{first} and check it against the text.
\item Draw the PERT network from the table, adding dummy tasks where needed.
\item Do the forward pass, then the backward pass, then the slack column.
\item Finish with two explicit sentences: the \textbf{project duration} and the \textbf{critical path}.
\item For a delay scenario: compare the delay with the task's slack; delay $\le$ slack means no impact on the end date.
\end{enumerate}
\end{methode}

\begin{exercice}[Your turn: a mini-project]
A school wants to deploy a management software. Tasks: A, needs analysis (2 weeks, none); B, buy server (3 weeks, after A); C, install software (2 weeks, after B); D, train teachers (1 week, after C); E, migrate data (2 weeks, after C); F, go live (1 week, after D and E).
\begin{enumerate}
\item Draw the precedence table and the PERT network.
\item Compute all slacks, the project duration and the critical path.
\item Training (D) is delayed by 2 weeks. What is the impact on go-live?
\end{enumerate}
\end{exercice}

\begin{corrige}[Exercise 1]
Forward pass: $A: 0\!-\!2$; $B: 2\!-\!5$; $C: 5\!-\!7$; $D: 7\!-\!8$; $E: 7\!-\!9$; $F: ES=\max(8,9)=9$, $EF=10$. Duration = \textbf{10 weeks}. Backward pass: $LF(F)=10$; $LF(D)=LF(E)=9$, so $LS(D)=8$ (slack 1), $LS(E)=7$ (slack 0); $C, B, A$ have slack 0. \textbf{Critical path:} $A \rightarrow B \rightarrow C \rightarrow E \rightarrow F$. Task $D$ has 1 week of slack; a 2-week delay exceeds it by 1 week, so go-live slips by \textbf{1 week} (to week 11) unless $D$ or $F$ is compressed.
\end{corrige}

\begin{savaistu}[Project management in Cameroon]
Large Cameroonian projects --- the deployment of fibre-optic backbone by CAMTEL, the rollout of mobile-money services by MTN and Orange, or the e-government platforms supervised by ANTIC --- all use the very tools you have just practised: charters, WBS, Gantt and PERT charts, milestone reviews and lessons-learned workshops. Mastering them on an 8-week school project is the first step towards managing national-scale ICT projects.
\end{savaistu}

\newpage

\coursec{Integration activity}

\courssub{Aims of the activity}

This integration activity closes the chapter by asking you to \emph{run} a complete project, not merely describe one. Working in teams of 4--5, you will mobilise every competence developed in Lessons 52 to 56:

\begin{itemize}
\item define a project and write a \cle{project charter} with \cle{SMART objectives};
\item identify and classify \cle{stakeholders};
\item decompose the work into a \cle{Work Breakdown Structure} (WBS) down to \cle{work packages};
\item build a \cle{precedence table}, draw a \cle{Gantt chart} and a \cle{PERT network};
\item compute \cle{ES}, \cle{EF}, \cle{LS}, \cle{LF}, the \cle{critical path} and \cle{slacks} (total and free);
\item monitor the project, react to a \cle{disruption} and propose \cle{corrective actions};
\item present and defend a plan orally before a decision-making audience.
\end{itemize}

\begin{aretenir}[Success criteria]
Your team succeeds if the plan is \textbf{feasible within 8 weeks and \SI{150000}{FCFA}}, if the critical path is correctly identified, and if your reaction to the disruption keeps the opening ceremony on schedule.
\end{aretenir}

\courssub{Situation}

\textbf{Government High School Mbankomo} (a town near Yaoundé) has just received a new computer room of 25 machines donated by a local NGO. To celebrate and to promote digital literacy, the \textbf{Principal} asks your Upper Sixth ICT class to organise a \textbf{Computer Science Week} (CS-Week). The event must take place \textbf{exactly 8 weeks from today} and the total budget granted by the Parents' Association (PTA) is \textbf{\SI{150000}{FCFA}}.

The week must include: an \textbf{opening ceremony} with invited guests (the Sub-Divisional Officer, PTA officials), a three-day \textbf{exhibition} of student projects, a \textbf{typing competition} (Forms 3--5), a \textbf{coding contest} (Lower and Upper Sixth), and a closing \textbf{award ceremony} with prizes.

Your team of 4--5 students is the \emph{project team}; the class teacher acts as \emph{project sponsor's representative}. After discussions, the team fixes the following tasks, durations (in weeks) and costs (in FCFA):

\begin{center}
\begin{tabularx}{\linewidth}{|c|X|c|c|c|}
\hline
\rowcolor{popblueL} \textbf{Code} & \textbf{Task} & \textbf{Dur.} & \textbf{Pred.} & \textbf{Cost} \\
\hline
A & Write charter, get Principal's approval & 1 & --- & 0 \\
\hline
B & Book hall, chairs, sound system & 1 & A & 20000 \\
\hline
C & Buy prizes and certificates & 2 & A & 60000 \\
\hline
D & Design and print posters, invitations & 2 & A & 15000 \\
\hline
E & Prepare lab: install software, test 25 PCs & 2 & A & 10000 \\
\hline
F & Publicity campaign (announcements, radio) & 2 & D & 10000 \\
\hline
G & Register participants, select contestants & 1 & F & 5000 \\
\hline
H & Rent projector and generator & 1 & B & 25000 \\
\hline
I & Rehearsal and final checks & 1 & C, E, G, H & 5000 \\
\hline
J & CS-Week itself (opening to award ceremony) & 1 & I & 0 \\
\hline
\end{tabularx}
\end{center}

Total planned cost: $20000+60000+15000+10000+10000+5000+25000+5000 = \SI{150000}{FCFA}$.

\textbf{The disruption.} At the end of week 4, the projector supplier in Yaoundé phones: delivery of the projector (task H, planned for week 5 in the team's schedule) will be delayed by \textbf{one week}. The opening ceremony cannot happen without a projector.

\courssub{Tasks}

\begin{enumerate}
\item \textbf{Charter.} Write the project charter: project title, purpose, two SMART objectives, scope (what is included and excluded), budget, deadline, and success criteria.
\item \textbf{Stakeholders.} List at least six stakeholders and classify each as \emph{sponsor}, \emph{team}, \emph{customer/beneficiary} or \emph{supplier}; state one expectation of each.
\item \textbf{WBS.} Draw the WBS with three levels: project $\rightarrow$ five deliverables $\rightarrow$ work packages (use tasks A--J at the lowest level).
\item \textbf{Scheduling.} From the precedence table above, draw a Gantt chart (weeks 1--8 on the horizontal axis).
\item \textbf{PERT and critical path.} Draw the PERT network, compute ES/EF by the forward pass and LS/LF by the backward pass, then give the critical path, the project duration, and the total slack of every non-critical task.
\item \textbf{Budget control.} At the end of week 4, tasks A, B, D, E are finished and C is half done; actual spending is \SI{85000}{FCFA}. Compute the planned value of work done and comment on the budget situation.
\item \textbf{Disruption management.} Analyse the impact of the one-week delay of task H on the schedule (use its slack). Propose \textbf{two} corrective actions and justify your choice.
\item \textbf{Defence.} Prepare a 5-minute presentation (charter, WBS, Gantt, critical path, disruption response) and defend it before the class acting as the school administration.
\end{enumerate}

\courssub{Model answer}

\textbf{1. Project charter (extract).}

\begin{definition}[Charter --- CS-Week]
\textbf{Title:} Computer Science Week, GHS Mbankomo. \\
\textbf{Purpose:} Inaugurate the new computer room and promote digital skills among students. \\
\textbf{Objectives (SMART):} 
\begin{enumerate}
\item Organise four events (exhibition, typing contest, coding contest, award ceremony) during week 8, with at least 200 participating students.
\item Complete the project within 8 weeks and \SI{150000}{FCFA}.
\end{enumerate}
\textbf{Scope:} Included --- the four events, publicity, logistics (hall, sound, projector, generator), prizes, lab preparation. Excluded --- permanent lab maintenance, internet subscription, teacher training. \\
\textbf{Budget:} \SI{150000}{FCFA} (PTA grant). \\
\textbf{Deadline:} End of week 8 (opening ceremony on Monday of week 8). \\
\textbf{Success criteria:} Opening ceremony held on the Monday of week 8; no budget overrun; at least 200 students participate.
\end{definition}

\textbf{2. Stakeholders.}

\begin{center}
\begin{tabularx}{\linewidth}{|l|X|X|}
\hline
\rowcolor{popblueL} \textbf{Stakeholder} & \textbf{Role} & \textbf{Expectation} \\
\hline
Principal & Sponsor & Event on time, school image enhanced \\
\hline
PTA & Sponsor (funds) & Budget respected, transparent accounting \\
\hline
ICT teacher & Supervisor / Sponsor's rep. & Pedagogical quality, student learning \\
\hline
Student team & Project team & Clear tasks, fair workload, recognition \\
\hline
Students of the school & Beneficiaries & Interesting, fair competitions, prizes \\
\hline
Projector supplier & Supplier & Clear order, prompt payment \\
\hline
Sub-Divisional Officer & Guest of honour / Authority & Smooth ceremony, protocol respected \\
\hline
\end{tabularx}
\end{center}

\textbf{3. Work Breakdown Structure (WBS).}

\begin{center}
\begin{tikzpicture}[
  level 1/.style={sibling distance=4.5cm, level distance=1.5cm},
  level 2/.style={sibling distance=2.2cm, level distance=1.5cm},
  level 3/.style={sibling distance=1.5cm, level distance=1.2cm},
  every node/.style={draw, rounded corners, fill=popblueL, text width=2.2cm, align=center, font=\small}
]
\node {CS-Week Project}
  child {node {Management} 
    child {node {A: Charter \& approval}}}
  child {node {Logistics} 
    child {node {B: Book hall, chairs, sound}}
    child {node {H: Rent projector \& generator}}}
  child {node {Communication} 
    child {node {D: Design/print posters}}
    child {node {F: Publicity campaign}}}
  child {node {Competitions} 
    child {node {C: Buy prizes \& certificates}}
    child {node {G: Register participants}}}
  child {node {Technical readiness} 
    child {node {E: Prepare lab}}
    child {node {I: Rehearsal \& checks}}
    child {node {J: CS-Week execution}}};
\end{tikzpicture}
\end{center}

\textbf{4. Gantt chart.}

\begin{popfigure}
\begin{tikzpicture}[x=0.7cm, y=0.55cm, font=\small]
% Time axis
\draw[->] (0,0) -- (9.5,0) node[right]{Weeks};
\foreach \w in {1,...,8} {
  \draw (\w,0.1) -- (\w,-0.1) node[below]{\w};
  \draw[popline, dashed] (\w,0) -- (\w,-9.5);
}
% Task bars
% A: week 1
\fill[popblue] (0,-0.5) rectangle (1,-1) node[midway, white]{A};
% B: week 2
\fill[popgreen] (1,-1.5) rectangle (2,-2) node[midway, white]{B};
% C: weeks 2-3
\fill[popgreen] (1,-2.5) rectangle (3,-3) node[midway, white]{C};
% D: weeks 2-3
\fill[popgreen] (1,-3.5) rectangle (3,-4) node[midway, white]{D};
% E: weeks 2-3
\fill[popgreen] (1,-4.5) rectangle (3,-5) node[midway, white]{E};
% F: weeks 4-5
\fill[poporange] (3,-5.5) rectangle (5,-6) node[midway, white]{F};
% G: week 6
\fill[poporange] (5,-6.5) rectangle (6,-7) node[midway, white]{G};
% H: week 3 (early) but scheduled at week 5 (late) in plan
\fill[popgreen] (4,-7.5) rectangle (5,-8) node[midway, white]{H};
% I: week 7
\fill[poppurple] (6,-8.5) rectangle (7,-9) node[midway, white]{I};
% J: week 8
\fill[poppurple] (7,-9.5) rectangle (8,-10) node[midway, white]{J};
% Critical path highlight: A, D, F, G, I, J
\draw[poppink, thick] (0,-0.75) -- (1,-0.75) -- (1,-3.75) -- (3,-3.75) -- (3,-5.75) -- (5,-5.75) -- (5,-6.75) -- (6,-6.75) -- (6,-8.75) -- (7,-8.75) -- (7,-9.75) -- (8,-9.75);
% Labels on left
\foreach \t/\y in {A/-0.75, B/-1.75, C/-2.75, D/-3.75, E/-4.75, F/-5.75, G/-6.75, H/-7.75, I/-8.75, J/-9.75}
  \node[left, align=right, text width=2.2cm] at (-0.3,\y) {\t};
\end{tikzpicture}
\legende{Gantt chart for CS-Week (critical path highlighted in pink). Task H is shown at its late start (week 5) as planned by the team.}
\end{popfigure}

\begin{center}
\begin{tabularx}{\linewidth}{|c|X|c|c|}
\hline
\rowcolor{popblueL} \textbf{Task} & \textbf{Bar (weeks)} & \textbf{ES} & \textbf{EF} \\
\hline
A & 1 & 0 & 1 \\
B & 2 & 1 & 2 \\
C & 2--3 & 1 & 3 \\
D & 2--3 & 1 & 3 \\
E & 2--3 & 1 & 3 \\
F & 4--5 & 3 & 5 \\
G & 6 & 5 & 6 \\
H & 5 & 4 & 5 \\
I & 7 & 6 & 7 \\
J & 8 & 7 & 8 \\
\hline
\end{tabularx}
\end{center}

\textbf{5. PERT network, critical path and slacks.}

\begin{methode}[Forward and backward pass]
\begin{enumerate}
\item \textbf{Forward pass:} For each task in topological order, $ES = \max(EF_{\text{pred}})$ (0 if no predecessor), $EF = ES + \text{duration}$.
\item \textbf{Backward pass:} Set project duration $T = \max(EF)$. For the last task(s), $LF = T$, $LS = LF - \text{duration}$. Then in reverse topological order, $LF = \min(LS_{\text{succ}})$, $LS = LF - \text{duration}$.
\item \textbf{Total slack:} $S = LS - ES = LF - EF$. Tasks with $S=0$ are critical.
\end{enumerate}
\end{methode}

Applying the method to our precedence table:

\begin{center}
\begin{tabular}{|c|c|c|c|c|c|c|}
\hline
\rowcolor{popblueL} \textbf{Task} & \textbf{Dur.} & \textbf{ES} & \textbf{EF} & \textbf{LS} & \textbf{LF} & \textbf{Total Slack} \\
\hline
A & 1 & 0 & 1 & 0 & 1 & \textbf{0} \\
B & 1 & 1 & 2 & 4 & 5 & 3 \\
C & 2 & 1 & 3 & 4 & 6 & 3 \\
D & 2 & 1 & 3 & 1 & 3 & \textbf{0} \\
E & 2 & 1 & 3 & 4 & 6 & 3 \\
F & 2 & 3 & 5 & 3 & 5 & \textbf{0} \\
G & 1 & 5 & 6 & 5 & 6 & \textbf{0} \\
H & 1 & 2 & 3 & 5 & 6 & 3 \\
I & 1 & 6 & 7 & 6 & 7 & \textbf{0} \\
J & 1 & 7 & 8 & 7 & 8 & \textbf{0} \\
\hline
\end{tabular}
\end{center}

\begin{propriete}[Critical path]
The critical path is \textbf{A $\rightarrow$ D $\rightarrow$ F $\rightarrow$ G $\rightarrow$ I $\rightarrow$ J}, of duration $1+2+2+1+1+1 = 8$ weeks. Tasks B, C, E and H each have a total slack of 3 weeks. Any delay on a critical task directly delays the whole project.
\end{propriete}

\begin{savaistu}[Free slack vs total slack]
Total slack is the delay a task can suffer without delaying the project finish. Free slack is the delay a task can suffer without delaying the early start of any successor. In this network, free slack equals total slack for all non-critical tasks because none of them share a successor with a critical task in a way that would reduce free slack. (Examinable: know the difference and how to compute both.)
\end{savaistu}

\textbf{6. Budget control (end of week 4).}

At the end of week 4, the work completed is:
\begin{itemize}
\item A (0 FCFA), B (20000), D (15000), E (10000) fully done.
\item C half done (planned cost 60000, so 30000 earned).
\end{itemize}
Planned Value (PV) = $0 + 20000 + 15000 + 10000 + 30000 = \SI{75000}{FCFA}$.
Actual Cost (AC) = \SI{85000}{FCFA}.
Cost Variance (CV) = PV $-$ AC = \SI{-10000}{FCFA} (overrun).

\begin{attention}[Budget overrun]
The project is \SI{10000}{FCFA} over budget on the work performed so far. The team must investigate immediately: perhaps prizes (task C) cost more than estimated, or there were hidden costs in logistics. Corrective action: negotiate discounts on remaining purchases, reduce publicity costs (task F), or seek a small top-up from the PTA. The overrun must be recovered before week 8.
\end{attention}

\textbf{7. Disruption management.}

\textbf{Analysis.} Task H (Rent projector and generator) has a total slack of 3 weeks (LS = 5, LF = 6). In the team's schedule, H was planned for week 5 (start at week 5, finish at week 6). The supplier announces a one-week delay, meaning the earliest H can now start is week 6 (finish week 7). Since LF = 6, a finish at week 7 would exceed the late finish by 1 week, delaying task I (which depends on H) and thus the critical path. However, note that H also has a predecessor B (finished week 2). The team can \emph{re-plan} H to start earlier if the projector becomes available earlier, but the supplier says delivery is delayed by one week from the planned week 5. If the planned week 5 was the late start (LS=5), then the new earliest start is week 6, which pushes LF to 7. But wait: the critical path analysis shows LF=6 for H. If H finishes at week 7, I (which has ES=6) would be delayed to start at week 7, pushing I to week 8 and J to week 9, extending the project to 9 weeks. \textbf{Therefore, the one-week delay on H, if H was scheduled at its late start, would delay the project by one week.}

But the team has a buffer: H has 3 weeks of total slack. The team can \textbf{reschedule H to its early start (week 3)} or any week between 3 and 5. However, at the end of week 4, week 3 has passed. The earliest H can now start is week 5 (if the projector arrives on time) or week 6 (with the delay). Since the disruption is announced at end of week 4, the team can still schedule H for week 5 (if the supplier delivers at the beginning of week 5) or week 6 (if delayed). The supplier says delivery delayed by one week from planned week 5, so delivery in week 6. That means H can only start in week 6, finish week 7. This would delay the project.

\textbf{Corrective actions:}
\begin{enumerate}
\item \textbf{Action 1 (Exploit slack by early procurement):} Contact the supplier immediately to see if a partial delivery (projector only) can be made by week 5, allowing H to start week 5 and finish week 6 (within LF). If not, rent a projector from a local cybercafé or another school for week 5 only (estimated extra cost \SI{5000}{FCFA}), and use the supplier's generator later. This keeps the project on schedule.
\item \textbf{Action 2 (Fast-track rehearsal):} If the projector absolutely cannot be available before week 6, reorganise the rehearsal (task I) to split: do the parts not needing the projector in week 6, and the projector-dependent part in week 7, but then compress the CS-Week (task J) by starting the opening ceremony later in week 8? Not acceptable because the opening ceremony is fixed on Monday week 8. So Action 1 is the only viable one.
\end{enumerate}

\textbf{Decision:} The team chooses \textbf{Action 1} (rent a projector locally for week 5) because it uses the existing slack (H can be done in week 5) at minimal extra cost (\SI{5000}{FCFA}), which can be covered by reducing the poster budget (task D) by \SI{5000}{FCFA} (print fewer posters, use digital displays). The supplier's generator can still be used in week 6 if needed. This keeps the critical path intact and the opening ceremony on Monday week 8.

\begin{attention}[Common mistakes]
\begin{itemize}
\item Forgetting that slack belongs to a \emph{path}, not only to a task; starting the backward pass from a wrong project duration.
\item Confusing free slack with total slack.
\item Panicking at a disruption before checking the slack of the affected task.
\item Assuming that a delay on a non-critical task never matters; it matters if the delay exceeds the task's total slack.
\item In budget control, confusing Planned Value (PV) with Earned Value (EV) or Actual Cost (AC).
\end{itemize}
\end{attention}

\textbf{8. Defence.}

The presentation follows the structure: Charter $\rightarrow$ WBS $\rightarrow$ Gantt $\rightarrow$ Critical path $\rightarrow$ Disruption response. Each team member speaks for about one minute. The Gantt chart and PERT table are projected. The team emphasises that the critical path is A-D-F-G-I-J, that the budget overrun is being addressed, and that the projector delay is absorbed by renting a local projector for one week, keeping the opening ceremony on schedule.

\begin{aretenir}[Key takeaways for the exam]
\begin{itemize}
\item A project charter must contain SMART objectives, scope, budget, deadline, success criteria.
\item Stakeholders: sponsor, team, customer/beneficiary, supplier.
\item WBS: hierarchical decomposition (project $\rightarrow$ deliverables $\rightarrow$ work packages).
\item Gantt chart: time-scaled bar chart; PERT: network diagram with dependencies.
\item Critical path: longest path (zero total slack); determines project duration.
\item Slack: total slack = LS $-$ ES; free slack = min(ES$_{\text{succ}}$) $-$ EF.
\item Budget control: PV (planned value of work done) vs AC (actual cost); CV = PV $-$ AC.
\item Disruption management: check slack of affected task; if delay $\le$ slack, absorb; else fast-track, crash, or re-scope.
\end{itemize}
\end{aretenir}

\newpage

\coursec{Methods and worked examples}

\courssub{Fundamentals of Project Management}

\begin{definition}[Project]
A \cle{project} is a \emph{temporary} endeavour undertaken to create a \emph{unique} product, service, or result. It has a defined start and end date, a specific objective, and consumes resources (people, money, equipment). Routine, repetitive work (e.g.\ daily backups, monthly payroll) is an \emph{operation}, not a project.
\end{definition}

\begin{definition}[Triple Constraint]
Every project balances three interdependent constraints:
\begin{itemize}
\item \cle{Scope} — what must be delivered (features, functions, quality standards);
\item \cle{Time} — the schedule, deadlines, and duration;
\item \cle{Cost} — the budget, human effort, and material resources.
\end{itemize}
\cle{Quality} sits at the centre: changing one constraint forces a change in at least one other.
\end{definition}

\begin{aretenir}[Key distinctions for the GCE]
\begin{itemize}
\item \textbf{Project vs Operation:} Projects are temporary and unique; operations are ongoing and repetitive.
\item \textbf{SMART objective:} Specific, Measurable, Achievable, Realistic, Time-bound.
\item \textbf{Stakeholders:} Sponsor (funds, high influence), Project Manager (leads), Team (executes), Users (operate the result), Suppliers (provide goods/services).
\item \textbf{Phases:} Initiation $\to$ Planning $\to$ Execution $\to$ Monitoring \& Control $\to$ Closure.
\end{itemize}
\end{aretenir}

\begin{methode}[Characterise a project and identify its constraints]
To analyse any project situation in the GCE examination:
\begin{enumerate}
\item \textbf{Check the definition.} Confirm the situation is a \cle{project}: temporary (start and end) and unique result. Routine work is an \emph{operation}.
\item \textbf{State the objective in SMART form}: Specific, Measurable, Achievable, Realistic, Time-bound.
\item \textbf{Identify the triple constraint}: \cle{scope} (deliverables), \cle{time} (deadline), \cle{cost} (budget). Quality is affected by all three.
\item \textbf{List the stakeholders}: sponsor, project manager, team, users, suppliers; note each one's interest and influence.
\item \textbf{Identify the project phases}: initiation, planning, execution, monitoring and control, closure.
\end{enumerate}
\textbf{Reflex:} if the examiner changes one constraint (e.g.\ shorter deadline), state explicitly which other constraint(s) must move (more cost, or reduced scope).
\end{methode}

\begin{savaistu}[Project management in Cameroon]
The \cle{ANTIC} (Agence Nationale des Technologies de l'Information et de la Communication) oversees major e-government projects (e.g.\ \emph{e-SISTEM} for civil service, \emph{e-Procurement}). Telecom operators (MTN, Orange) run network rollout projects using PERT/Gantt tools. Mobile Money platform upgrades are classic ICT projects with tight scope-time-cost trade-offs.
\end{savaistu}

\begin{exoresolu}[Analysing the ``E-Santé Yaoundé'' project]
The Yaoundé Central Hospital wants a computerised patient-registration system. The Director gives the ICT team 4 months and a budget of 8\,000\,000 FCFA. The system must register patients, print consultation tickets and produce daily statistics. Halfway through the project, the Director asks that the system be ready in 3 months instead of 4, without extra money.
\begin{enumerate}
\item Show that this is a project and state its objective in SMART form.
\item Identify the triple constraint and the main stakeholders.
\item Explain the consequences of the Director's new demand.
\end{enumerate}
\tcblower
\textbf{1. A project and its SMART objective.}\\
The situation is a \cle{project} because it is \emph{temporary} (4 months, defined start and end) and produces a \emph{unique} result (a registration system tailored to this hospital). It is not routine maintenance work, which would be an operation.\\
A SMART objective: \emph{``Deliver, within 4 months and for at most 8\,000\,000 FCFA, a working patient-registration system for Yaoundé Central Hospital that registers patients, prints consultation tickets and produces daily statistical reports.''} It is Specific (three named functions), Measurable (system works, reports produced), Achievable and Realistic (a small, well-defined system), Time-bound (4 months).

\textbf{2. Triple constraint and stakeholders.}
\begin{center}
\begin{tabularx}{\linewidth}{|l|X|}\hline
\rowcolor{popblueL} \textbf{Constraint} & \textbf{Value in this project} \\ \hline
Scope & Registration, ticket printing, daily statistics \\ \hline
Time & 4 months \\ \hline
Cost & 8\,000\,000 FCFA (team salaries, hardware, software licences) \\ \hline
\end{tabularx}
\end{center}
Stakeholders: the \emph{sponsor} (hospital Director, high influence), the \emph{project manager} and \emph{ICT team} (do the work), the \emph{users} (receptionists, doctors, nurses — high interest), the \emph{patients} (affected by the result), and possible \emph{suppliers} (computers, network equipment).

\textbf{3. Consequences of the new demand.}\\
The \emph{time} constraint is tightened (4 $\to$ 3 months) while \emph{cost} is frozen. By the triple-constraint principle, at least one other element must give way:
\begin{itemize}
\item either the \textbf{scope is reduced} (e.g.\ drop the statistics module for version 1);
\item or \textbf{quality is endangered} (less testing, more risk of bugs at deployment);
\item or the team must \textbf{work overtime}, which usually increases hidden cost and fatigue.
\end{itemize}
The correct project-management reflex is to \emph{negotiate}: present the Director with options (reduce scope, or add budget for temporary staff, or accept higher risk) rather than silently accepting an unrealistic constraint.
\end{exoresolu}

\begin{attention}[Common GCE trap]
Do not confuse \emph{scope} with \emph{requirements}. Scope is the \emph{sum of deliverables}; requirements are the detailed conditions each deliverable must meet. In the WBS you break down scope, not requirements.
\end{attention}

\courssub{Project Planning and the Work Breakdown Structure (WBS)}

\begin{definition}[Work Breakdown Structure (WBS)]
A \cle{WBS} is a hierarchical decomposition of the \emph{total scope of work} into manageable components called \cle{work packages}. It organises and defines the project's deliverables. The WBS is \emph{deliverable-oriented}, not task-oriented — it shows \textbf{what} must be produced, not \textbf{when} or \textbf{who}.
\end{definition}

\begin{propriete}[The 100\% Rule]
The children of any WBS node must represent \emph{exactly 100\%} of the work of the parent — nothing missing, nothing extra. This ensures no work falls through the cracks and no gold-plating occurs.
\end{propriete}

\begin{methode}[Construct a WBS]
\begin{enumerate}
\item Write the \textbf{final deliverable} (project name) at the top — level 1.
\item Decompose into \textbf{major deliverables or phases} — level 2 (e.g.\ analysis, design, implementation, testing, deployment).
\item Decompose each branch into \textbf{work packages} — level 3: small, assignable tasks (typically 8 to 80 hours of work) that one person or a small team can own.
\item Apply the \textbf{100\% rule}: the children of a node must represent \emph{all} the work of the parent — nothing missing, nothing extra.
\item \textbf{Number} the elements hierarchically: 1, 1.1, 1.1.1, \dots
\item Check: a WBS shows \emph{deliverables and work}, \textbf{not} durations, dates or dependencies (those belong to scheduling).
\end{enumerate}
\textbf{Reflex:} use nouns or verb+noun for work packages (``Design database'', not ``Database'' alone when it is a task).
\end{methode}

\begin{experience}[Building a WBS in a spreadsheet]
Open LibreOffice Calc or Excel. Column A: WBS code (1, 1.1, 1.1.1). Column B: Work package name. Column C: Owner. Column D: Estimated hours. Use indentation (Data $\to$ Group) to visualise hierarchy. This WBS dictionary becomes the input for scheduling.
\end{experience}

\begin{exoresolu}[WBS for a school website project]
Government High School Nkolbisson asks its Computer Club to develop the school's website. The work covers: requirements gathering, graphic design, content writing, coding of the pages, testing, and putting the site online.
\begin{enumerate}
\item Draw the WBS of this project on three levels, with hierarchical numbering.
\item State two checks you make to validate your WBS.
\end{enumerate}
\tcblower
\textbf{1. The WBS.} Level 1 is the final deliverable; level 2 the phases; level 3 the work packages.
\begin{center}
\begin{tikzpicture}[
 grow=right,
 level 1/.style={sibling distance=52mm, level distance=30mm},
 level 2/.style={sibling distance=13mm, level distance=34mm},
 every node/.style={draw=popblue, fill=popblueL, rounded corners=2pt, font=\scriptsize, align=center, text width=24mm},
 edge from parent/.style={draw=popmuted}]
\node[fill=poporange, text=white, text width=22mm]{\textbf{1. School Website}}
 child { node {1.5 Deployment} 
   child { node {1.5.2 Put site online} }
   child { node {1.5.1 Configure hosting} } }
 child { node {1.4 Testing}
   child { node {1.4.2 Fix defects} }
   child { node {1.4.1 Test pages \& links} } }
 child { node {1.3 Implementation}
   child { node {1.3.2 Write content} }
   child { node {1.3.1 Code pages} } }
 child { node {1.2 Design}
   child { node {1.2.2 Design page layout} }
   child { node {1.2.1 Choose colours \& logo} } }
 child { node {1.1 Analysis}
   child { node {1.1.2 Write specification} }
   child { node {1.1.1 Collect requirements} } };
\end{tikzpicture}
\legende{Three-level WBS of the school website project (read rightwards from the root).}
\end{center}

\textbf{2. Validation checks.}
\begin{itemize}
\item \textbf{100\% rule:} the union of 1.1 to 1.5 covers \emph{all} the work mentioned (requirements, design, content, coding, testing, going online) and nothing outside the project (e.g.\ ``maintain the site for one year'' is excluded — it is an operation).
\item \textbf{Work packages are assignable and estimable:} each level-3 item (e.g.\ 1.3.1 ``Code pages'') can be given to one club member and estimated in hours. If a package were too big (e.g.\ ``Build the whole site''), it would be decomposed further.
\end{itemize}
\end{exoresolu}

\begin{attention}[WBS vs Schedule]
A WBS has \textbf{no arrows, no dates, no durations}. It is a static scope baseline. The schedule (Gantt, PERT) is built \emph{from} the WBS work packages by adding durations, dependencies, and resources.
\end{attention}

\courssub{Project Scheduling with Gantt and PERT Charts}

\begin{definition}[Gantt Chart]
A \cle{Gantt chart} is a bar chart that illustrates a project schedule. The horizontal axis represents time; the vertical axis lists tasks. Each task is a horizontal bar whose length is proportional to its duration. Milestones (zero-duration events) are shown as diamonds. A Gantt chart shows \emph{when} tasks happen and their overlap, but \textbf{does not by itself reveal the critical path}.
\end{definition}

\begin{definition}[PERT Network]
A \cle{PERT} (Program Evaluation and Review Technique) network is a directed graph where nodes represent tasks (Activity-on-Node) or events (Activity-on-Arrow) and arrows represent precedence constraints. It is used to compute earliest/latest times, slack, and the critical path.
\end{definition}

\begin{methode}[Build a Gantt chart]
\begin{enumerate}
\item List the \textbf{tasks} (from the WBS work packages) with their \textbf{durations} and \textbf{predecessors}.
\item Draw a \textbf{time axis} (days, weeks or months) on the horizontal; list tasks on the vertical, in logical order.
\item For each task, draw a \textbf{horizontal bar} starting after its predecessors finish (for simple Gantt: at its planned start date), with length $=$ duration.
\item Mark \textbf{milestones} (key events, zero duration) with a diamond.
\item Read off: the \textbf{project duration} (end of the last bar), \textbf{parallel tasks} (overlapping bars), and \textbf{idle resources}.
\end{enumerate}
\textbf{Reflex:} a Gantt chart shows \emph{when} tasks happen; it does not by itself prove which tasks are critical — that needs the network (PERT) analysis of Method 4.
\end{methode}

\begin{exoresolu}[Gantt chart for a cybercafé network installation]
Mr.\ Ekane opens a cybercafé in Buea. The tasks are:
\begin{center}
\begin{tabular}{|c|l|c|c|}\hline
\rowcolor{popblueL} \textbf{Task} & \textbf{Description} & \textbf{Duration (days)} & \textbf{Predecessor} \\ \hline
A & Paint and wire the room & 4 & -- \\ \hline
B & Buy computers & 3 & -- \\ \hline
C & Install network cabling & 2 & A \\ \hline
D & Set up computers & 2 & B, C \\ \hline
E & Install software and test & 2 & D \\ \hline
\end{tabular}
\end{center}
\begin{enumerate}
\item Draw the Gantt chart, starting each task as early as possible.
\item Give the minimum project duration and identify which tasks run in parallel.
\end{enumerate}
\tcblower
\textbf{1. Gantt chart.} A and B have no predecessor: they start at day 0 in parallel. C waits for A (day 4). D waits for B \emph{and} C, so it starts at $\max(3, 4+2) = 6$. E follows D at day 8 and ends at day 10.
\begin{center}
\begin{tikzpicture}[x=0.85cm, y=0.62cm]
\draw[->, popmuted] (0,0) -- (10.8,0) node[right, font=\scriptsize]{days};
\foreach \x in {0,1,...,10} { \draw[popline] (\x,0) -- (\x,5.6); \node[below, font=\scriptsize, popmuted] at (\x,0) {\x}; }
\foreach \y/\name in {1/A, 2/B, 3/C, 4/D, 5/E} \node[left, font=\scriptsize\bfseries, popdark] at (0,\y+0.4) {\name};
\fill[popblue]   (0,1.1) rectangle (4,1.7);
\fill[popgreen]  (0,2.1) rectangle (3,2.7);
\fill[poppurple] (4,3.1) rectangle (6,3.7);
\fill[poporange] (6,4.1) rectangle (8,4.7);
\fill[poppink]   (8,5.1) rectangle (10,5.7);
\end{tikzpicture}
\legende{Gantt chart of the cybercafé project (earliest starts).}
\end{center}

\textbf{2. Duration and parallelism.} The last bar (E) ends at day 10, so the \textbf{minimum duration is 10 days}. Tasks \textbf{A and B run in parallel} from day 0 to day 3. Note that B finishes at day 3 but D cannot start until day 6: the person responsible for B has 3 idle days — a resource-planning observation the Gantt chart makes visible.
\end{exoresolu}

\begin{methode}[Draw the PERT network and compute the critical path]
\begin{enumerate}
\item From the task table (task, duration, predecessors), draw \textbf{nodes} (tasks) and \textbf{arrows} showing dependencies: an arrow $X \to Y$ means $X$ must finish before $Y$ starts. (Activity-on-Node convention used in GCE.)
\item \textbf{Forward pass} — earliest start/finish:
\[ ES_{\text{first}} = 0, \qquad EF = ES + d, \qquad ES_j = \max_{\text{predecessors } i}(EF_i). \]
\item \textbf{Backward pass} — latest start/finish, starting from the project end $T = \max(EF)$:
\[ LF_{\text{last}} = T, \qquad LS = LF - d, \qquad LF_i = \min_{\text{successors } j}(LS_j). \]
\item \textbf{Total slack} of a task: $S = LS - ES = LF - EF$.
\item The \textbf{critical path} is the chain of tasks with \emph{zero slack}; its length is the minimum project duration.
\end{enumerate}
\textbf{Reflex:} at a merge (several arrows in), take the \emph{maximum} going forward; at a split (several arrows out), take the \emph{minimum} going backward.
\end{methode}

\begin{exoresolu}[Critical path of the cybercafé project]
Using the task table of the previous exercise (A: 4 days; B: 3 days; C: 2 days after A; D: 2 days after B and C; E: 2 days after D):
\begin{enumerate}
\item Draw the PERT network.
\item Compute $ES$, $EF$, $LS$, $LF$ and the total slack of every task.
\item Deduce the critical path and the minimum project duration.
\end{enumerate}
\tcblower
\textbf{1. PERT network.}
\begin{center}
\begin{tikzpicture}[x=2.2cm, y=1.4cm,
 task/.style={draw=popblue, fill=popblueL, rounded corners=2pt, minimum width=14mm, minimum height=8mm, font=\scriptsize\bfseries},
 arr/.style={->, thick, popdark}]
\node[task] (A) at (0,1) {A (4)};
\node[task] (B) at (0,0) {B (3)};
\node[task] (C) at (2,1) {C (2)};
\node[task] (D) at (4,0.5) {D (2)};
\node[task] (E) at (6,0.5) {E (2)};
\draw[arr] (A) -- (C);
\draw[arr] (C) -- (D);
\draw[arr] (B) -- (D);
\draw[arr] (D) -- (E);
\end{tikzpicture}
\legende{PERT network: arrows show precedence, brackets show durations.}
\end{center}

\textbf{2. Forward and backward passes.}\\
\emph{Forward pass} ($EF = ES + d$; at a merge, take the max):
\begin{align*}
ES_A &= 0, & EF_A &= 0+4 = 4;\\
ES_B &= 0, & EF_B &= 0+3 = 3;\\
ES_C &= EF_A = 4, & EF_C &= 4+2 = 6;\\
ES_D &= \max(EF_B, EF_C) = \max(3,6) = 6, & EF_D &= 6+2 = 8;\\
ES_E &= EF_D = 8, & EF_E &= 8+2 = 10.
\end{align*}
Project duration $T = 10$ days.\\
\emph{Backward pass} ($LS = LF - d$; at a split, take the min):
\begin{align*}
LF_E &= 10, & LS_E &= 10-2 = 8;\\
LF_D &= LS_E = 8, & LS_D &= 8-2 = 6;\\
LF_C &= LS_D = 6, & LS_C &= 6-2 = 4;\\
LF_B &= LS_D = 6, & LS_B &= 6-3 = 3;\\
LF_A &= LS_C = 4, & LS_A &= 4-4 = 0.
\end{align*}
\textbf{Summary table} ($S = LS - ES$):
\begin{center}
\begin{tabular}{|c|c|c|c|c|c|c|}\hline
\rowcolor{popblueL} \textbf{Task} & $d$ & $ES$ & $EF$ & $LS$ & $LF$ & \textbf{Slack} \\ \hline
A & 4 & 0 & 4 & 0 & 4 & \textbf{0} \\ \hline
B & 3 & 0 & 3 & 3 & 6 & \textbf{3} \\ \hline
C & 2 & 4 & 6 & 4 & 6 & \textbf{0} \\ \hline
D & 2 & 6 & 8 & 6 & 8 & \textbf{0} \\ \hline
E & 2 & 8 & 10 & 8 & 10 & \textbf{0} \\ \hline
\end{tabular}
\end{center}

\textbf{3. Critical path.} The zero-slack tasks are A, C, D, E: the \textbf{critical path is A $\to$ C $\to$ D $\to$ E}, of length $4+2+2+2 = 10$ \textbf{days}, which equals the minimum project duration. Any delay on A, C, D or E delays the whole project by the same amount. Task B has 3 days of slack: it can start any time between day 0 and day 3 without affecting the opening date.
\end{exoresolu}

\begin{savaistu}[GCE vocabulary: Activity-on-Node vs Activity-on-Arrow]
Cameroon GCE uses \textbf{Activity-on-Node} (tasks are boxes, arrows show precedence). Some textbooks use Activity-on-Arrow (events are nodes, tasks are arrows). Both give the same critical path; stick to Activity-on-Node for the exam.
\end{savaistu}

\courssub{Project Control, Critical Path and Slack Analysis}

\begin{definition}[Total Slack (Float)]
The \cle{total slack} of a task is the amount of time it can be delayed without delaying the \emph{project finish date}. $S_{\text{total}} = LS - ES = LF - EF$. On the critical path, total slack is zero.
\end{definition}

\begin{definition}[Free Slack]
The \cle{free slack} of a task is the amount of time it can be delayed without delaying the \emph{earliest start of any immediate successor}. $S_{\text{free}}(i) = \min_{\text{successors } j}(ES_j) - EF_i$. Free slack $\le$ total slack always.
\end{definition}

\begin{methode}[Use slack to control a project]
\begin{enumerate}
\item Compute \textbf{total slack} $S = LS - ES$ for every task (Method 4).
\item Compute \textbf{free slack} for each task: $S_{\text{free}}(i) = \min_{\text{successors } j}(ES_j) - EF_i$.
\item During \textbf{monitoring}, compare planned vs actual progress (Gantt tracking, milestones) and compute the \textbf{variance}.
\item If a task slips: if the slip $\le$ total slack, no impact on the deadline; otherwise the project end shifts by (slip $-$ slack) and you must act: \textbf{reallocate resources} from slack tasks to critical tasks, \textbf{fast-track} (parallelise tasks that were sequential), or \textbf{crash} (add resources to shorten a critical task, at extra cost).
\end{enumerate}
\textbf{Reflex:} free slack $\le$ total slack always; on the critical path both are zero.
\end{methode}

\begin{exoresolu}[Controlling the cybercafé project]
Same project (A:4; B:3; C:2 after A; D:2 after B,C; E:2 after D; $T = 10$ days).
\begin{enumerate}
\item Compute the free slack of task B and interpret it.
\item At the end of day 5, the supervisor finds that A finished on time but C (cabling) has only done 1 day of work out of 2. What is the impact on the opening date if nothing is done? Propose two corrective actions.
\item Instead, suppose the purchase of computers (B) is delayed by 2 days. Does the opening date change? Justify.
\end{enumerate}
\tcblower
\textbf{1. Free slack of B.} The only successor of B is D, with $ES_D = 6$:
\[ S_{\text{free}}(B) = ES_D - EF_B = 6 - 3 = 3 \text{ days}. \]
Interpretation: B can finish up to 3 days late without delaying \emph{any} other task. Here free slack equals total slack (3) because D's start is fixed by C, not by B.

\textbf{2. C is one day behind.} C is \textbf{critical} (slack $= 0$). Losing 1 day on C pushes D to day 7--9 and E to day 9--11: \textbf{the opening is delayed by 1 day}, to day 11, if nothing is done. Corrective actions:
\begin{itemize}
\item \textbf{Crash:} assign a second technician to cabling so C finishes on day 6 as planned (extra labour cost, but deadline saved);
\item \textbf{Fast-track / reallocate:} the worker who finished B at day 3 is idle; move him to help on C, or let the software installation of E be prepared in parallel (pre-downloading software images) so E shortens by 1 day.
\end{itemize}

\textbf{3. B delayed by 2 days.} B has a \textbf{total slack of 3 days}. A delay of 2 days $\le 3$ days: B then finishes at day 5, still before day 6 when D must start. \textbf{The opening date does not change} (still day 10); only B's remaining slack falls from 3 to 1 day. This is exactly why slack analysis matters: it tells the manager \emph{which delays are dangerous} (critical tasks) and \emph{which can be absorbed}.
\end{exoresolu}

\begin{attention}[Monitoring traps]
\begin{itemize}
\item Do not confuse \emph{slack} with \emph{buffer}. Slack is a mathematical property of the network; a buffer is a deliberate time reserve added by the manager.
\item Reporting ``task 50\% complete'' is meaningless without knowing if it is on the critical path. Always report \emph{earned value} or \emph{slack consumption}.
\end{itemize}
\end{attention}

\courssub{Crashing: Shortening a Project at Minimum Cost}

\begin{definition}[Crashing]
\cle{Crashing} is a schedule compression technique that adds resources to critical-path tasks to shorten their duration at the least incremental cost. Only critical tasks can be crashed to reduce project duration; crashing non-critical tasks wastes money.
\end{definition}

\begin{methode}[Crash a project]
\begin{enumerate}
\item For each task that can be shortened, compute the \textbf{crash cost per day} (cost slope):
\[ \text{cost slope} = \frac{\text{crash cost} - \text{normal cost}}{\text{normal duration} - \text{crash duration}}. \]
\item Identify the \textbf{critical path}; only shortening critical tasks reduces the project duration.
\item Crash the critical task with the \textbf{smallest cost slope} first, one day at a time.
\item After each day crashed, \textbf{re-check the critical path}: a new path may become critical, and then both paths must be shortened together.
\item Stop when the target duration is reached or no further crashing is possible; add up the \textbf{total extra cost}.
\end{enumerate}
\textbf{Reflex:} never crash a non-critical task — it costs money and saves no time.
\end{methode}

\begin{exoresolu}[Crashing the cybercafé project]
Mr.\ Ekane wants to open in \textbf{9 days instead of 10}. Crashing data:
\begin{center}
\begin{tabular}{|c|c|c|c|c|}\hline
\rowcolor{popblueL} \textbf{Task} & \textbf{Normal (days)} & \textbf{Crash (days)} & \textbf{Normal cost (FCFA)} & \textbf{Crash cost (FCFA)} \\ \hline
A & 4 & 3 & 80\,000 & 110\,000 \\ \hline
C & 2 & 1 & 40\,000 & 90\,000 \\ \hline
D & 2 & 1 & 50\,000 & 75\,000 \\ \hline
E & 2 & 2 & 60\,000 & 60\,000 \\ \hline
\end{tabular}
\end{center}
Which task should be crashed, and at what extra cost?
\tcblower
\textbf{Step 1 — Cost slopes} (cost slope $= \dfrac{\Delta \text{cost}}{\Delta \text{days}}$):
\begin{center}
\begin{tabular}{|c|c|c|c|}\hline
\rowcolor{popblueL} \textbf{Task} & \textbf{Max reduction} & \textbf{Extra cost} & \textbf{Cost slope (FCFA/day)} \\ \hline
A & $4-3 = 1$ & $110\,000 - 80\,000 = 30\,000$ & 30\,000 \\ \hline
C & $2-1 = 1$ & $90\,000 - 40\,000 = 50\,000$ & 50\,000 \\ \hline
D & $2-1 = 1$ & $75\,000 - 50\,000 = 25\,000$ & 25\,000 \\ \hline
E & $2-2 = 0$ & 0 & cannot be crashed \\ \hline
\end{tabular}
\end{center}
\textbf{Step 2 — Critical path.} From Method 4, the critical path is A $\to$ C $\to$ D $\to$ E (10 days). Only A, C, D are candidates (E cannot be shortened).

\textbf{Step 3 — Cheapest first.} The smallest cost slope on the critical path is \textbf{D at 25\,000 FCFA/day}. Crash D from 2 to 1 day: the path becomes $4+2+1+2 = 9$ days.

\textbf{Step 4 — Re-check the critical path.} With D $=1$: path A--C--D--E $= 9$ days; path B--D--E $= 3+1+2 = 6$ days. The original path is still the unique critical path, and the target of 9 days is reached.

\textbf{Conclusion.} Crash \textbf{task D by 1 day}; the project lasts \textbf{9 days} for a total extra cost of \textbf{25\,000 FCFA} (total cost rises from 230\,000 to 255\,000 FCFA). Note the trap: A has a bigger possible reduction but costs 30\,000 FCFA/day — crashing D first is cheaper.
\end{exoresolu}

\begin{savaistu}[Crashing in Cameroonian practice]
When CAMTEL rolls out fibre-to-the-home in a new neighbourhood, crashing might mean hiring extra splicing teams (overtime pay) to meet a ministerial deadline. The cost slope is the overtime premium per day saved. Project managers at MTN Cameroon use Primavera P6 or MS Project to automate this analysis.
\end{savaistu}

\courssub{Integration: Running a Complete Sample Project}

\begin{methode}[Run a sample project end-to-end]
\begin{enumerate}
\item \textbf{Initiate:} write a short project charter (objective, sponsor, budget, deadline).
\item \textbf{Plan:} build the WBS $\to$ derive the task list with durations and predecessors $\to$ draw the Gantt chart and the PERT network $\to$ find the critical path and slacks.
\item \textbf{Execute and control:} track progress against the Gantt baseline; when a deviation appears, quantify its effect using slack, then decide (reallocate, fast-track, crash).
\item \textbf{Close:} deliver, obtain acceptance, and write the \emph{lessons learned} report.
\end{enumerate}
\textbf{Reflex (GCE):} integration questions chain all skills — answer in this order and justify each number you use.
\end{methode}

\begin{exoresolu}[The ``Vote Night'' project — full integration]
A students' association in Bamenda organises an election-results night. Tasks: A — book hall (2 d); B — print ballots (3 d, after A); C — configure computers and network (2 d, after A); D — install results software (1 d, after C); E — train operators (2 d, after B and C); F — full rehearsal (1 d, after D and E).
\begin{enumerate}
\item Draw the PERT network and compute the minimum duration.
\item Give the slack of every task and the critical path.
\item On day 4, printing (B) is reported 1 day late. Consequence and decision?
\end{enumerate}
\tcblower
\textbf{1. Network and duration.}\\
Forward pass:
\begin{align*}
ES_A&=0,\ EF_A=2; & ES_B&=2,\ EF_B=5; & ES_C&=2,\ EF_C=4;\\
ES_D&=4,\ EF_D=5; & ES_E&=\max(5,4)=5,\ EF_E=7; & ES_F&=\max(5,7)=7,\ EF_F=8.
\end{align*}
Minimum duration: $T = 8$ \textbf{days}.

\textbf{2. Backward pass and slacks.} $LF_F = 8$, $LS_F = 7$; $LF_E = 7$, $LS_E = 5$; $LF_D = 7$, $LS_D = 6$; $LF_C = \min(LS_D, LS_E) = \min(6,5) = 5$, $LS_C = 3$; $LF_B = LS_E = 5$, $LS_B = 2$; $LF_A = \min(LS_B, LS_C) = \min(2,3) = 2$, $LS_A = 0$.
\begin{center}
\begin{tabular}{|c|c|c|c|c|c|c|}\hline
\rowcolor{popblueL} \textbf{Task} & $d$ & $ES$ & $EF$ & $LS$ & $LF$ & \textbf{Slack} \\ \hline
A & 2 & 0 & 2 & 0 & 2 & \textbf{0} \\ \hline
B & 3 & 2 & 5 & 2 & 5 & \textbf{0} \\ \hline
C & 2 & 2 & 4 & 3 & 5 & \textbf{1} \\ \hline
D & 1 & 4 & 5 & 6 & 7 & \textbf{2} \\ \hline
E & 2 & 5 & 7 & 5 & 7 & \textbf{0} \\ \hline
F & 1 & 7 & 8 & 7 & 8 & \textbf{0} \\ \hline
\end{tabular}
\end{center}
\textbf{Critical path: A $\to$ B $\to$ E $\to$ F} ($2+3+2+1 = 8$ days). Note that C and D (the IT branch) carry slack: the technical team has margin, the printing does not.

\textbf{3. B one day late.} B is critical (slack 0): the rehearsal F moves to day 8--9 and the event is \textbf{delayed by 1 day} unless action is taken. Decision options: reallocate the idle IT technician (who has 1--2 days of slack on C/D) to help finish the printing preparation, or shorten E by training operators in two parallel groups. \textbf{Lesson for closure:} printing was the hidden bottleneck — in future, order ballots earlier or add a buffer day on the critical path. This closes the full cycle: charter $\to$ WBS $\to$ schedule $\to$ critical path $\to$ control $\to$ lessons learned.
\end{exoresolu}

\begin{exercice}[Practice: ``Mobile Money Kiosk'' deployment]
A fintech startup deploys a Mobile Money kiosk in Douala. Tasks:
\begin{center}
\begin{tabular}{|c|l|c|cc|}\hline
\rowcolor{popblueL} \textbf{ID} & \textbf{Task} & \textbf{Dur. (days)} & \textbf{Pred.} \\ \hline
A & Site survey \& contract & 3 & -- \\ \hline
B & Order hardware (POS, printer, router) & 5 & -- \\ \hline
C & Electrical wiring \& security cage & 4 & A \\ \hline
D & Install hardware & 2 & B, C \\ \hline
E & Configure software \& test transactions & 3 & D \\ \hline
F & Staff training \& pilot & 2 & E \\ \hline
G & Go-live \& handover & 1 & F \\ \hline
\end{tabular}
\end{center}
\begin{enumerate}
\item Draw the WBS (3 levels) and the PERT network.
\item Compute ES, EF, LS, LF, total slack, free slack for each task.
\item Identify the critical path and the project duration.
\item Task B (hardware delivery) is delayed by 3 days. Impact on go-live? Propose a crashing option if task E can be reduced from 3 to 2 days at an extra cost of 150\,000 FCFA.
\end{enumerate}
\end{exercice}

\begin{corrige}[Exercise: Mobile Money Kiosk]
\textbf{1. WBS (3 levels).}
\begin{center}
\begin{tikzpicture}[
 grow=right,
 level 1/.style={sibling distance=50mm, level distance=28mm},
 level 2/.style={sibling distance=12mm, level distance=32mm},
 every node/.style={draw=popblue, fill=popblueL, rounded corners=2pt, font=\scriptsize, align=center, text width=22mm},
 edge from parent/.style={draw=popmuted}]
\node[fill=poporange, text=white, text width=20mm]{\textbf{1. MM Kiosk Deployment}}
 child { node {1.7 Go-live} child { node {1.7.1 Handover} } }
 child { node {1.6 Training} child { node {1.6.2 Pilot run} } child { node {1.6.1 Train staff} } }
 child { node {1.5 Config \& Test} child { node {1.5.2 Test transactions} } child { node {1.5.1 Config software} } }
 child { node {1.4 Install HW} child { node {1.4.2 Connect peripherals} } child { node {1.4.1 Mount POS \& router} } }
 child { node {1.3 Wiring \& Cage} child { node {1.3.2 Build security cage} } child { node {1.3.1 Electrical wiring} } }
 child { node {1.2 Procurement} child { node {1.2.2 Receive hardware} } child { node {1.2.1 Place orders} } }
 child { node {1.1 Initiation} child { node {1.1.2 Sign contract} } child { node {1.1.1 Site survey} } };
\end{tikzpicture}
\end{center}

\textbf{2. PERT network and forward/backward pass.}
\begin{center}
\begin{tikzpicture}[x=2.0cm, y=1.2cm,
 task/.style={draw=popblue, fill=popblueL, rounded corners=2pt, minimum width=13mm, minimum height=7mm, font=\scriptsize\bfseries},
 arr/.style={->, thick, popdark}]
\node[task] (A) at (0,2) {A (3)};
\node[task] (B) at (0,1) {B (5)};
\node[task] (C) at (2,2) {C (4)};
\node[task] (D) at (4,1.5) {D (2)};
\node[task] (E) at (6,1.5) {E (3)};
\node[task] (F) at (8,1.5) {F (2)};
\node[task] (G) at (10,1.5) {G (1)};
\draw[arr] (A) -- (C);
\draw[arr] (C) -- (D);
\draw[arr] (B) -- (D);
\draw[arr] (D) -- (E);
\draw[arr] (E) -- (F);
\draw[arr] (F) -- (G);
\end{tikzpicture}
\end{center}

Forward pass:
$ES_A=0, EF_A=3$; $ES_B=0, EF_B=5$; $ES_C=3, EF_C=7$; $ES_D=\max(5,7)=7, EF_D=9$; $ES_E=9, EF_E=12$; $ES_F=12, EF_F=14$; $ES_G=14, EF_G=15$.
Project duration $T = 15$ days.

Backward pass:
$LF_G=15, LS_G=14$; $LF_F=14, LS_F=12$; $LF_E=12, LS_E=9$; $LF_D=9, LS_D=7$; $LF_C=7, LS_C=3$; $LF_B=7, LS_B=2$; $LF_A=3, LS_A=0$.

\textbf{3. Slack table.}
\begin{center}
\begin{tabular}{|c|c|c|c|c|c|c|c|}\hline
\rowcolor{popblueL} \textbf{Task} & $d$ & $ES$ & $EF$ & $LS$ & $LF$ & \textbf{Total Slack} & \textbf{Free Slack} \\ \hline
A & 3 & 0 & 3 & 0 & 3 & 0 & 0 \\ \hline
B & 5 & 0 & 5 & 2 & 7 & 2 & 2 \\ \hline
C & 4 & 3 & 7 & 3 & 7 & 0 & 0 \\ \hline
D & 2 & 7 & 9 & 7 & 9 & 0 & 0 \\ \hline
E & 3 & 9 & 12 & 9 & 12 & 0 & 0 \\ \hline
F & 2 & 12 & 14 & 12 & 14 & 0 & 0 \\ \hline
G & 1 & 14 & 15 & 14 & 15 & 0 & 0 \\ \hline
\end{tabular}
\end{center}

\textbf{4. Critical path:} A $\to$ C $\to$ D $\to$ E $\to$ F $\to$ G (15 days). Tasks B has 2 days total (and free) slack.

\textbf{5. B delayed by 3 days.} B's total slack is 2 days. A 3-day delay exceeds slack by 1 day $\to$ project end shifts to day 16. To recover: crash E from 3 to 2 days (cost slope 150\,000 FCFA/day). New path A-C-D-E-F-G = 14 days, but B-D-E-F-G becomes $5+3+2+2+1=13$ days (with B delayed to 8, path = 8+2+2+1=13). Critical path remains A-C-D-E-F-G at 14 days. Go-live moves from day 15 to day 14 (saved 1 day) at cost 150\,000 FCFA. If the 3-day delay on B is accepted without crashing, go-live is day 16.
\end{corrige}

\begin{aretenir}[Chapter summary for revision]
\begin{itemize}
\item \textbf{Project} = temporary + unique. Triple constraint: scope, time, cost (quality central).
\item \textbf{WBS} = hierarchical scope decomposition (100\% rule), deliverable-oriented, no dates.
\item \textbf{Gantt} = time-scaled bars; shows overlap, not critical path.
\item \textbf{PERT} = network diagram; forward pass (ES/EF), backward pass (LS/LF), slack = LS-ES.
\item \textbf{Critical path} = zero-slack tasks; length = minimum project duration.
\item \textbf{Total slack} = delay without moving project end; \textbf{free slack} = delay without moving any successor.
\item \textbf{Control:} monitor variance, use slack to absorb slips, else reallocate/fast-track/crash.
\item \textbf{Crashing:} only critical tasks; lowest cost slope first; re-check critical path after each day.
\item \textbf{Integration:} charter $\to$ WBS $\to$ schedule $\to$ critical path $\to$ control $\to$ lessons learned.
\end{itemize}
\end{aretenir}

\newpage

\coursec{Exercises}

\courssub{I check my knowledge}

\begin{exercice}[MCQ on Project Management Basics]
Choose the correct answer for each question.
\begin{enumerate}
    \item Which of the following best defines a project?
    \begin{itemize}
        \item[A.] A repetitive process that produces the same output continuously.
        \item[B.] A temporary endeavour undertaken to create a unique product, service, or result.
        \item[C.] An ongoing operational activity with no defined end date.
        \item[D.] A routine task performed by a functional department.
    \end{itemize}
    \item The Triple Constraint in project management consists of:
    \begin{itemize}
        \item[A.] Time, Cost, and Quality.
        \item[B.] Scope, Time, and Cost.
        \item[C.] Scope, Quality, and Risk.
        \item[D.] Time, Risk, and Communication.
    \end{itemize}
    \item A Work Breakdown Structure (WBS) is primarily used to:
    \begin{itemize}
        \item[A.] Schedule activities on a calendar.
        \item[B.] Identify project stakeholders.
        \item[C.] Decompose the project scope into manageable work packages.
        \item[D.] Calculate the critical path.
    \end{itemize}
    \item In a PERT chart, the critical path is the path that:
    \begin{itemize}
        \item[A.] Has the most activities.
        \item[B.] Has the longest total duration and determines the minimum project duration.
        \item[C.] Has the largest amount of slack.
        \item[D.] Contains only milestone activities.
    \end{itemize}
\end{enumerate}
\end{exercice}

\begin{exercice}[True/False with Justification]
State whether each statement is True or False. Justify your answer briefly.
\begin{enumerate}
    \item A project milestone has a duration of zero and marks a significant point in the project timeline.
    \item Slack (or float) is the amount of time an activity can be delayed without delaying the project finish date.
    \item The Work Breakdown Structure (WBS) should be organised by project phases only, not by deliverables.
    \item In a Gantt chart, the critical path activities are always shown in red and have no slack.
    \item A stakeholder is anyone who is actively working on the project tasks.
\end{enumerate}
\end{exercice}

\begin{exercice}[Definitions and Distinctions]
Provide clear, concise definitions for the following terms and distinguish between the pairs where indicated.
\begin{enumerate}
    \item Project
    \item Stakeholder
    \item Milestone
    \item Slack (Float)
    \item Distinguish between a \emph{project} and \emph{operations}.
    \item State the three elements of the Triple Constraint and explain how they interact.
\end{enumerate}
\end{exercice}

\begin{exercice}[Direct Calculations: Slack and Critical Path]
Consider the following activity list for a small project. All durations are in days.
\begin{center}
\begin{tabular}{|c|c|c|}
\hline
Activity & Duration & Immediate Predecessors \\
\hline
A & 3 & -- \\
B & 4 & -- \\
C & 2 & A \\
D & 5 & A \\
E & 3 & B \\
F & 4 & C, D \\
G & 2 & E \\
H & 3 & F, G \\
\hline
\end{tabular}
\end{center}
\begin{enumerate}
    \item Draw the Activity-on-Node (AON) network diagram.
    \item Compute the Earliest Start (ES), Earliest Finish (EF), Latest Start (LS), and Latest Finish (LF) for each activity.
    \item Calculate the Total Slack for each activity.
    \item Identify the critical path and the minimum project duration.
\end{enumerate}
\end{exercice}

\courssub{I practise}

\begin{exercice}[WBS for VSAT Installation in Bamenda]
A cybercafé in Bamenda wants to install a VSAT internet connection. Build a three-level Work Breakdown Structure (WBS) for this project. Use a standard coding scheme (e.g., 1.0, 1.1, 1.1.1). Include at least 15 work packages at the lowest level. Present your WBS in a hierarchical table or indented list.
\end{exercice}

\begin{exercice}[PERT Network and Critical Path Analysis]
The table below shows activities for a software development project. Durations are in weeks.
\begin{center}
\begin{tabular}{|c|c|c|}
\hline
Activity & Duration (weeks) & Predecessors \\
\hline
A & 2 & -- \\
B & 3 & -- \\
C & 4 & A \\
D & 3 & A \\
E & 5 & B \\
F & 2 & C \\
G & 4 & D, E \\
H & 3 & F \\
I & 2 & G, H \\
\hline
\end{tabular}
\end{center}
\begin{enumerate}
    \item Construct the PERT network diagram (Activity-on-Node).
    \item Perform the forward pass to determine ES and EF for all activities.
    \item Perform the backward pass to determine LS and LF for all activities.
    \item Compute the Total Float for each activity.
    \item Identify the critical path(s) and the project duration.
    \item If activity D is delayed by 2 weeks, what is the impact on the project completion date? Justify.
\end{enumerate}
\end{exercice}

\begin{exercice}[Gantt Chart Interpretation]
The Gantt chart below represents the planned vs. actual progress of a project at the end of Week 5. The project started on Monday of Week 1. Each bar represents one activity. The planned bars are shown in blue (top), actual progress in green (bottom, solid for completed, hatched for in-progress). Today is Friday of Week 5.

\begin{center}
\begin{tikzpicture}[scale=0.7]
% Time axis
\draw[->] (0,0) -- (14,0) node[right] {Weeks};
\foreach \x in {1,2,...,13} \draw (\x,0.1) -- (\x,-0.1) node[below] {\x};
% Activities
% A
\draw[popblue, thick] (0,4) -- (3,4) node[left] {A};
\draw[popgreen, thick] (0,3.8) -- (3,3.8);
% B
\draw[popblue, thick] (0,3) -- (4,3) node[left] {B};
\draw[popgreen, thick] (0,2.8) -- (4,2.8);
% C
\draw[popblue, thick] (3,2) -- (6,2) node[left] {C};
\draw[popgreen, thick, pattern=north east lines] (3,1.8) -- (5,1.8);
% D
\draw[popblue, thick] (3,1) -- (7,1) node[left] {D};
\draw[popgreen, thick] (3,0.8) -- (5,0.8);
% E
\draw[popblue, thick] (4,0) -- (8,0) node[left] {E};
\draw[popgreen, thick, pattern=north east lines] (4,-0.2) -- (6,-0.2);
\end{tikzpicture}
\end{center}
\begin{enumerate}
    \item Which activities are behind schedule at Week 5? Explain.
    \item Which activities have slack (float) based on the planned chart? How can you tell?
    \item What is the project status at Week 5 in terms of overall progress?
    \item If the project must finish by the end of Week 10, which activities are most critical to monitor?
\end{enumerate}
\end{exercice}

\begin{exercice}[Project Control Measures for School Website]
A secondary school in Yaoundé wants to launch a new website before the GCE registration period (which starts in 12 weeks). The project team has defined the scope and created a WBS and schedule.
\begin{enumerate}
    \item Propose three SMART objectives for this project.
    \item Design a milestone plan with at least five milestones, including dates (week numbers).
    \item Suggest three specific project control measures (tools/techniques) the project manager should implement to ensure the website is ready on time. Explain how each measure helps.
\end{enumerate}
\end{exercice}

\courssub{I prepare for the GCE}

\begin{exercice}[Structured Question: PERT/CPM Full Analysis] \hfill (25 marks)
A project to set up a computer network for a new campus has the following activities, durations (in days), and dependencies:

\begin{center}
\begin{tabular}{|c|c|c|}
\hline
Activity & Duration (days) & Predecessors \\
\hline
A & 5 & -- \\
B & 3 & -- \\
C & 4 & A \\
D & 6 & A \\
E & 2 & B \\
F & 7 & C \\
G & 4 & D, E \\
H & 3 & F \\
I & 5 & G, H \\
\hline
\end{tabular}
\end{center}

\begin{enumerate}
    \item Draw the Activity-on-Node (AON) network diagram for the project. \hfill (4 marks)
    \item Calculate the Earliest Start (ES) and Earliest Finish (EF) times for each activity using the forward pass. \hfill (5 marks)
    \item Calculate the Latest Start (LS) and Latest Finish (LF) times for each activity using the backward pass. \hfill (5 marks)
    \item Determine the Total Float (Slack) for each activity. \hfill (3 marks)
    \item Identify the critical path and state the minimum project duration. \hfill (3 marks)
    \item Activity D is delayed by 4 days due to late delivery of cables. Activity E is completed 2 days early. Determine the new project duration and explain whether the critical path changes. \hfill (5 marks)
\end{enumerate}
\end{exercice}

\begin{exercice}[Scenario: School Website Project] \hfill (20 marks)
A school in Douala decides to develop a website to publish GCE results and allow online registration. The project must be completed in 10 weeks. The project manager has identified the following high-level deliverables: Requirements Specification, Design, Development, Testing, Deployment, and Training.

\begin{enumerate}
    \item Write a SMART project objective statement for this project. \hfill (3 marks)
    \item Create a Work Breakdown Structure (WBS) down to work package level (at least 3 levels) for the project. Use a coding system. \hfill (6 marks)
    \item From the WBS, list 8 activities with estimated durations (in weeks) and logical dependencies to create a schedule. Present them in a table. \hfill (5 marks)
    \item Draw a Gantt chart (sketch) for the project based on your activities. Indicate the critical path. \hfill (3 marks)
    \item Describe two risk management strategies the project manager should apply to ensure the website is ready before the registration period. \hfill (3 marks)
\end{enumerate}
\end{exercice}

\begin{exercice}[Essay: Impact of Delays on Project Schedule] \hfill (15 marks)
In a project to install a Local Area Network (LAN) for a regional hospital, the critical path has a duration of 30 days. Activity X is on the critical path with a duration of 5 days. Activity Y is not on the critical path; it has a duration of 4 days and a Total Float of 6 days.

\begin{enumerate}
    \item Explain the concepts of \emph{critical path} and \emph{total float} in project scheduling. \hfill (4 marks)
    \item Activity X is delayed by 3 days due to equipment failure. Calculate the new project duration and explain the impact. \hfill (3 marks)
    \item Activity Y is delayed by 5 days due to staff absence. Calculate the new project duration and explain the impact. \hfill (3 marks)
    \item If both delays occur simultaneously (X delayed by 3 days, Y delayed by 5 days), what is the overall project delay? Justify your answer with reference to the critical path and float. \hfill (5 marks)
\end{enumerate}
\end{exercice}

\newpage

\coursec{Solutions to the exercises}

\begin{corrige}[Exercise 1 — Multiple choice questions]
\begin{enumerate}
\item \textbf{Answer: (b).} A project is a \emph{temporary} endeavour (it has a defined start and end) undertaken to create a \emph{unique} product, service or result. Option (a) describes routine operations, (c) is far too broad, and (d) describes a permanent organisational unit, the opposite of a project.
\item \textbf{Answer: (b).} The triple constraint (or ``iron triangle'') links \cle{scope}, \cle{time} and \cle{cost}: changing one of them necessarily affects at least one of the other two.
\item \textbf{Answer: (c).} The critical path is the \emph{longest} path through the network; its length fixes the \emph{minimum} possible duration of the project, because no path can be completed faster than its own total duration.
\item \textbf{Answer: (b).} Slack (float) is the maximum delay an activity can accumulate \emph{without} delaying the end of the project: $S = LS - ES = LF - EF$.
\end{enumerate}
\end{corrige}

\begin{corrige}[Exercise 2 — True or false]
\begin{enumerate}
\item \textbf{False.} A project is temporary and unique, so it is managed with specific tools (WBS, Gantt, PERT) and a dedicated team; routine operations are repetitive and permanent, managed through standard procedures and hierarchical line management.
\item \textbf{True.} By definition, an activity with zero total slack has no margin: any delay on it delays the whole project, which is exactly the property of activities on the critical path.
\item \textbf{False.} The Gantt chart shows \emph{time positioning} (start, duration, progress) clearly, but dependencies are shown only implicitly (or with thin arrows). The PERT network is precisely designed to display logical precedence relations between activities.
\item \textbf{False.} A milestone is a significant \emph{event} or checkpoint (for example ``contract signed''); it has \emph{zero} duration, not a long one.
\item \textbf{True.} The WBS is a hierarchical, deliverable-oriented decomposition of the project into levels, down to work packages small enough to be estimated, assigned and controlled.
\end{enumerate}
\end{corrige}

\begin{corrige}[Exercise 3 — Key definitions]
\begin{enumerate}
\item \textbf{Project:} a temporary endeavour, with a defined start and end, undertaken to create a unique product, service or result within given constraints of scope, time and cost. \emph{Example:} installing a computer laboratory of 20 machines at GBHS Bamenda before the start of the school year.
\item \textbf{Stakeholder:} any person, group or organisation that can affect or is affected by the project, or that has an interest in its outcome. \emph{Example:} for a mobile-money kiosk network project in Douala, stakeholders include the owner, the agents, MTN or Orange, and the customers.
\item \textbf{Milestone:} a significant event or checkpoint of zero duration marking the completion of an important stage. \emph{Example:} ``website validated by the school administration'' on 15 March.
\item \textbf{Deliverable:} a tangible, verifiable product or result that must be produced to complete the project or one of its phases. \emph{Example:} the signed feasibility report, or the installed and configured VSAT antenna.
\item \textbf{Slack (float):} the maximum delay an activity can accumulate without delaying the end of the project, computed as $S = LS - ES = LF - EF$. \emph{Example:} if printing T-shirts for the school open day can slip by 3 days without affecting the ceremony, that activity has 3 days of slack.
\end{enumerate}
\end{corrige}

\begin{corrige}[Exercise 4 — Direct calculations]
\begin{enumerate}
\item Earliest finish: $EF = ES + d = 6 + 4 = 10$ days.
\item Latest start: $LS = LF - d = 15 - 4 = 11$ days.
\item Total slack: $S = LS - ES = 11 - 6 = 5$ days (check: $S = LF - EF = 15 - 10 = 5$). Since $S = 5 \neq 0$, activity $T$ is \textbf{not} on the critical path: it can be delayed by up to 5 days without delaying the project.
\end{enumerate}
\end{corrige}

\begin{corrige}[Exercise 5 — Building a WBS]
\textbf{1. and 2.} A possible three-level WBS with coding:

\begin{center}
\begin{tabularx}{\linewidth}{|l|X|}\hline
\rowcolor{popblueL} \textbf{Code} & \textbf{Element} \\ \hline
1 & VSAT internet connection for the cybercafé \\ \hline
1.1 & Studies and procurement \\ \hline
1.1.1 & Choose the VSAT provider and subscription \\ \hline
1.1.2 & Obtain quotations and sign the contract \\ \hline
1.2 & Site preparation \\ \hline
1.2.1 & Prepare the roof support for the antenna \\ \hline
1.2.2 & Bring power supply and earthing to the equipment room \\ \hline
1.3 & Installation and configuration \\ \hline
1.3.1 & Mount and align the VSAT antenna \\ \hline
1.3.2 & Install modem, router and Wi-Fi distribution \\ \hline
1.4 & Testing and commissioning \\ \hline
1.4.1 & Test bandwidth and stability \\ \hline
1.4.2 & Train staff and hand over documentation \\ \hline
\end{tabularx}
\end{center}

\textbf{3.} Two work packages at the lowest level:
\begin{itemize}
\item \textbf{1.3.1 Mount and align the VSAT antenna} --- deliverable: the antenna physically installed and aligned with the satellite; milestone: ``antenna operational (signal locked)''.
\item \textbf{1.4.1 Test bandwidth and stability} --- deliverable: a signed test report showing measured bandwidth; milestone: ``connection accepted by the owner''.
\end{itemize}
\emph{Examiner's note:} any coherent decomposition is accepted provided it is hierarchical, deliverable-oriented, coded consistently, and the lowest-level packages are small enough to be assigned to one responsible person.
\end{corrige}

\begin{corrige}[Exercise 6 — Reading a Gantt chart]
\textbf{1.} Planned Gantt chart (weeks 1 to 8):

\begin{popfigure}
\begin{tikzpicture}[x=0.95cm,y=0.55cm]
\draw[->] (0,0) -- (9.3,0);
\foreach \w in {1,...,8}{
  \draw[popline] (\w-1,0) -- (\w-1,5.6);
  \node[below, font=\small] at (\w-0.5,0) {\w};
}
\draw[popline] (8,0) -- (8,5.6);
\node[below, font=\small] at (0,0) {0};
\fill[popblue] (0,4.6) rectangle (2,5.2); \node[left, font=\small] at (0,4.9) {A};
\fill[popgreen] (2,3.6) rectangle (4,4.2); \node[left, font=\small] at (0,3.9) {B};
\fill[poporange] (2,2.6) rectangle (5,3.2); \node[left, font=\small] at (0,2.9) {C};
\fill[poppurple] (4,1.6) rectangle (6,2.2); \node[left, font=\small] at (0,1.9) {D};
\fill[poppink] (6,0.6) rectangle (7,1.2); \node[left, font=\small] at (0,0.9) {E};
\draw[popdark, dashed, thick] (5,0) -- (5,5.6) node[above, font=\small] {end of week 5};
\end{tikzpicture}
\legende{Planned Gantt chart; the dashed line marks the end of week 5.}
\end{popfigure}

\textbf{2.} Situation at the end of week 5:
\begin{itemize}
\item \textbf{A} (planned weeks 1--2): 100\,\% finished $\Rightarrow$ \emph{on schedule} (finished).
\item \textbf{B} (planned weeks 3--4): 100\,\% finished $\Rightarrow$ \emph{on schedule} (finished).
\item \textbf{C} (planned weeks 3--5): only 50\,\% done instead of 100\,\% $\Rightarrow$ \emph{late} (about 1.5 weeks of work remain).
\item \textbf{D} (planned weeks 5--6): should have started at week 5 but has not started $\Rightarrow$ \emph{late}.
\item \textbf{E} (planned week 7): not due yet $\Rightarrow$ \emph{on schedule} (not started, as planned).
\end{itemize}

\textbf{3.} Project status: C is late but non-critical with 2 weeks of slack, so its delay (about 1--1.5 weeks) can be absorbed without moving the end date. D, however, is \emph{critical} and has not started on time: every week of delay on D directly postpones the end of the project. Conclusion: the project is \textbf{at risk of finishing late}; the manager must immediately reallocate resources to start D (and accelerate C, since C's remaining work may block D if D depends on the room being ready).
\end{corrige}

\begin{corrige}[Exercise 7 — Effect of delays]
\begin{enumerate}
\item $X$ is on the critical path, so its slack is zero: any delay is fully transmitted to the project. A delay of 3 days on $X$ \textbf{postpones the project end date by exactly 3 days}.
\item $Y$ has a total slack of 4 days. A delay of 5 days exceeds the slack by $5 - 4 = 1$ day. The first 4 days are absorbed by the float; the remaining day shifts the project. \textbf{Effect: the project end date is postponed by 1 day} (and $Y$ becomes critical).
\item Original minimum duration: 30 days. The delay on $X$ adds 3 days; the delay on $Y$ adds 1 day. Assuming the two paths are independent, the new end date is fixed by the worst shift: $30 + \max(3, 1) = 33$ days. \textbf{New project end date: day 33} (a 3-day overall delay).
\end{enumerate}
\end{corrige}

\begin{corrige}[Exercise 8 — SMART objectives and milestones]
\textbf{1.} Two SMART objectives (Specific, Measurable, Achievable, Relevant, Time-bound):
\begin{itemize}
\item ``Publish the school's new website, containing the presentation pages, the results portal and the online pre-registration form, no later than 30 April, i.e.\ at least two weeks before GCE registration opens.''
\item ``Achieve a page loading time below 3 seconds on a 3G connection and train 2 administrative staff members to update the site before 15 May.''
\end{itemize}

\textbf{2.} Milestone plan (dates relative to project start, week 0):
\begin{center}
\begin{tabularx}{\linewidth}{|l|X|c|}\hline
\rowcolor{popblueL} \textbf{Milestone} & \textbf{Meaning} & \textbf{Target} \\ \hline
M1 & Requirements validated by the administration & end of week 2 \\ \hline
M2 & Graphic design and site map approved & end of week 4 \\ \hline
M3 & Development finished, internal tests passed & end of week 8 \\ \hline
M4 & Site online and staff trained & end of week 10 \\ \hline
\end{tabularx}
\end{center}

\textbf{3.} Three control measures:
\begin{itemize}
\item weekly progress meetings comparing actual progress with the Gantt baseline (schedule control);
\item validation of each deliverable against the requirements before starting the next phase (quality control);
\item tracking expenditure against the budget and managing changes through a formal change-request procedure (cost and scope control).
\end{itemize}
\end{corrige}

\begin{corrige}[Exercise 9 — Structured question: PERT network and critical path]
\textbf{Indicative mark scheme (20 marks):} network 6, earliest dates 4, latest dates 4, slack and critical path 4, duration and delay analysis 2.

\textbf{1. PERT network (6 marks).} Predecessors: A none; B after A; C and D after B; E after D; F after C \emph{and} E; G after F; H after G. Since F has two predecessors, a dummy activity (or a converging node structure) is used.

\begin{popfigure}
\begin{tikzpicture}[->, >=stealth, thick, node distance=1.9cm,
  evt/.style={circle, draw=popdark, fill=popblueL, minimum size=7mm, font=\small}]
\node[evt] (n0) {0};
\node[evt, right of=n0] (n1) {1};
\node[evt, right of=n1, xshift=6mm] (n2) {2};
\node[evt, above right=4mm and 2.2cm of n2] (n3) {3};
\node[evt, below right=4mm and 2.2cm of n2] (n4) {4};
\node[evt, right=2.6cm of n4] (n5) {5};
\node[evt, right=2.2cm of n5] (n6) {6};
\node[evt, right of=n6] (n7) {7};
\node[evt, right of=n7] (n8) {8};
\draw (n0) -- node[above, font=\small]{A (2)} (n1);
\draw (n1) -- node[above, font=\small]{B (3)} (n2);
\draw (n2) -- node[above, font=\small]{C (5)} (n3);
\draw (n2) -- node[below, font=\small]{D (4)} (n4);
\draw (n4) -- node[below, font=\small]{E (3)} (n5);
\draw (n3) -- node[above, font=\small]{dummy (0)} (n5);
\draw (n5) -- node[above, font=\small]{F (2)} (n6);
\draw (n6) -- node[above, font=\small]{G (3)} (n7);
\draw (n7) -- node[above, font=\small]{H (1)} (n8);
\end{tikzpicture}
\legende{PERT network; the dummy activity carries the dependency C $\to$ F.}
\end{popfigure}

\emph{Marking:} 1 mark per correct precedence (A$\to$B, B$\to$C, B$\to$D, D$\to$E, C+E$\to$F with dummy, F$\to$G$\to$H), up to 6.

\textbf{2. Earliest dates (forward pass, 4 marks).} $ES = \max(EF \text{ of predecessors})$, $EF = ES + d$:
\begin{center}
\begin{tabular}{|l|c|c|c|c|}\hline
\rowcolor{popblueL} \textbf{Act.} & $d$ & $ES$ & $EF$ & \\ \hline
A & 2 & 0 & 2 & \\ \hline
B & 3 & 2 & 5 & \\ \hline
C & 5 & 5 & 10 & \\ \hline
D & 4 & 5 & 9 & \\ \hline
E & 3 & 9 & 12 & \\ \hline
F & 2 & $\max(10,12)=12$ & 14 & \\ \hline
G & 3 & 14 & 17 & \\ \hline
H & 1 & 17 & 18 & \\ \hline
\end{tabular}
\end{center}

\textbf{3. Latest dates (backward pass, 4 marks).} $LF(\text{H}) = 18$; $LF = \min(LS \text{ of successors})$, $LS = LF - d$:
\begin{center}
\begin{tabular}{|l|c|c|c|}\hline
\rowcolor{popblueL} \textbf{Act.} & $LF$ & $LS$ & \\ \hline
H & 18 & 17 & \\ \hline
G & 17 & 14 & \\ \hline
F & 14 & 12 & \\ \hline
E & 12 & 9 & \\ \hline
C & 12 & 7 & \\ \hline
D & 9 & 5 & \\ \hline
B & 5 & 2 & \\ \hline
A & 2 & 0 & \\ \hline
\end{tabular}
\end{center}

\textbf{4. Slack and critical path (4 marks).} $S = LS - ES$:
\begin{center}
\begin{tabular}{|l|c|c|c|c|c|c|c|c|}\hline
\rowcolor{popblueL} \textbf{Act.} & A & B & C & D & E & F & G & H \\ \hline
$S$ & 0 & 0 & 2 & 0 & 0 & 0 & 0 & 0 \\ \hline
\end{tabular}
\end{center}
Only C has slack (2 days). The \textbf{critical path is A -- B -- D -- E -- F -- G -- H}.

\textbf{5. Minimum duration and delay on D (2 marks).} Minimum duration $= EF(H) = \textbf{18 days}$. D is critical (slack 0), so a delay of 2 days on D shifts E, F, G, H and \textbf{postpones the project end by exactly 2 days: new duration 20 days}. (Note: a delay on C of up to 2 days would have had no effect, since C has 2 days of slack.)
\end{corrige}

\begin{corrige}[Exercise 10 — Structured question: planning a school management system]
\textbf{Indicative mark scheme (20 marks):} (a) 5, (b) 4, (c) 5, (d) 4, (e) 2.

\textbf{1. Definitions and stakeholders (5 marks).}
\begin{itemize}
\item \emph{Project} (1): a temporary endeavour undertaken to create a unique product, service or result.
\item \emph{Stakeholder} (1): any person or group who can affect or is affected by the project.
\item \emph{Milestone} (1): a zero-duration event marking the completion of an important stage.
\item Four stakeholders (2, i.e.\ 0.5 each): the college proprietor/administration, the teachers, the students and their parents, the development team (also acceptable: the bursar, the ministry inspectorate).
\end{itemize}

\textbf{2. Triple constraint (4 marks).} The triple constraint links \cle{scope}, \cle{time} and \cle{cost}: the three are interdependent, so modifying one forces an adjustment on at least one other (2 marks). Illustration (2 marks): if the budget is reduced, the college may have to \emph{reduce the scope} (drop the online payment module, keep only fees and marks) \emph{or extend the schedule} (hire fewer developers, so development takes longer) --- or accept lower quality, which is risky.

\textbf{3. Three-level WBS (5 marks).} Example of a coherent, coded decomposition:
\begin{center}
\begin{tabularx}{\linewidth}{|l|X|}\hline
\rowcolor{popblueL} \textbf{Code} & \textbf{Element} \\ \hline
1 & School management system \\ \hline
1.1 & Analysis and design \\ \hline
1.1.1 & Collect requirements (fees, marks, report cards) \\ \hline
1.1.2 & Design database and interfaces \\ \hline
1.2 & Development \\ \hline
1.2.1 & Fees module \\ \hline
1.2.2 & Marks and report cards module \\ \hline
1.3 & Testing and deployment \\ \hline
1.3.1 & Test modules and fix defects \\ \hline
1.3.2 & Install, migrate data and train users \\ \hline
\end{tabularx}
\end{center}
\emph{Marking:} hierarchical structure 2, three levels 1, coherent coding 1, relevance to the case 1.

\textbf{4. Gantt chart (4 marks).} Dependencies: analysis (w.\,1--2) $\to$ design (w.\,3--5) $\to$ development (w.\,6--11) $\to$ testing (w.\,12--13); training (1 week) can start after design, e.g.\ week 6 (in parallel with development).

\begin{popfigure}
\begin{tikzpicture}[x=0.62cm,y=0.55cm]
\draw[->] (0,0) -- (14.4,0);
\foreach \w in {1,...,13}{
  \draw[popline] (\w-1,0) -- (\w-1,5.6);
  \node[below, font=\scriptsize] at (\w-0.5,0) {\w};
}
\draw[popline] (13,0) -- (13,5.6);
\fill[popblue] (0,4.6) rectangle (2,5.2); \node[left, font=\scriptsize] at (0,4.9) {Analysis};
\fill[popgreen] (2,3.6) rectangle (5,4.2); \node[left, font=\scriptsize] at (0,3.9) {Design};
\fill[poporange] (5,2.6) rectangle (11,3.2); \node[left, font=\scriptsize] at (0,2.9) {Development};
\fill[poppurple] (11,1.6) rectangle (13,2.2); \node[left, font=\scriptsize] at (0,1.9) {Testing};
\fill[poppink] (5,0.6) rectangle (6,1.2); \node[left, font=\scriptsize] at (0,0.9) {Training};
\end{tikzpicture}
\legende{Gantt chart of the school management system project.}
\end{popfigure}

The \textbf{planned finish is the end of week 13} (critical chain: analysis $\to$ design $\to$ development $\to$ testing $= 2+3+6+2 = 13$ weeks).

\textbf{5. Limits of the Gantt chart (2 marks).} (i) It does not show logical dependencies clearly $\Rightarrow$ use a \textbf{PERT network}. (ii) It does not show which delays are dangerous (slack/criticality) $\Rightarrow$ use \textbf{critical path and slack analysis} (also accepted: cost control $\to$ budget tracking tools).
\end{corrige}

\begin{corrige}[Exercise 11 — Essay question: managing an e-government style project]
\textbf{Indicative mark scheme (20 marks):} (a) 8, (b) 4, (c) 5, (d) 3. The examiner expects a structured essay (introduction, developed paragraphs, conclusion), correct vocabulary, and systematic application to the council's platform.

\textbf{Introduction.} The project ``online tax and documents platform'' is a temporary, unique endeavour; it must be managed with the recognised project management framework.

\textbf{1. The five process groups applied to the project (8 marks; about 1.5 per group).}
\begin{itemize}
\item \emph{Initiating:} define the need (citizens queue for hours to pay taxes), authorise the project with a project charter, identify stakeholders (council, taxpayers, ANTIC for security, the payment operators).
\item \emph{Planning:} define scope (online payment, document download), build the WBS, schedule (Gantt/PERT), budget, risk plan (power cuts, cyber-attacks) and quality plan.
\item \emph{Executing:} develop the platform, buy and configure servers, integrate mobile money payment (MTN MoMo, Orange Money), train council staff, communicate with citizens.
\item \emph{Monitoring and controlling:} compare actual progress and spending with the baseline, track the critical path, manage change requests, verify security requirements continuously.
\item \emph{Closing:} run the final acceptance tests, obtain formal sign-off by the council, hand over documentation, write the lessons-learned report, release the team.
\end{itemize}

\textbf{2. Role and qualities of the project manager (4 marks).} The project manager plans, organises, leads and controls: he/she coordinates the team, manages the budget and schedule, communicates with stakeholders and manages risks (2 marks). Required qualities (2 marks): leadership, communication and negotiation skills, rigour and organisation, technical understanding of ICT and security, ability to decide under pressure.

\textbf{3. Critical path and slack analysis (5 marks).} The security audit is \emph{critical} (slack 0): a one-week delay \textbf{postpones the go-live date by exactly one week} (1.5 marks). The manager must react: reallocate resources, fast-track subsequent tasks, renegotiate the launch date (1 mark). The logo design has 10 days of slack; a delay of 12 days exceeds the slack by $12 - 10 = 2$ days, so it \textbf{postpones the project by 2 days} and the logo activity becomes critical (2 marks). Conclusion of the analysis: the two delays are on different paths; the overall delay is fixed by the worst shift, i.e.\ $\max(7, 2) = 7$ days if the paths are independent (0.5 mark).

\textbf{4. Three success factors in the Cameroonian context (3 marks).} Any three, well justified: strong commitment of the hierarchy (the council) and clear sponsorship; involvement and training of end users (staff and citizens); reliable infrastructure planning (power supply, internet connectivity, e.g.\ via CAMTEL fibre); partnership with local operators (mobile money integration); rigorous security and data protection in line with ANTIC guidelines; realistic budgeting and progressive deployment.

\textbf{Conclusion.} A disciplined application of the process groups, combined with critical path and slack analysis, allows the council to deliver the platform on time, within budget, and with the expected quality.
\end{corrige}

\newpage

\coursec{Summary sheet}

\begin{popfigure}
\begin{tikzpicture}[mindmap, grow cyclic, every node/.style={concept, text=white, align=center, font=\small}, concept color=popblue, level 1/.append style={level distance=3.6cm, sibling angle=60}, level 2/.append style={level distance=2.4cm, sibling angle=45}]
\node[concept, font=\small\bfseries, align=center]{Project\\Management}
  child[concept color=poporange] { node {Fundamentals} 
      child { node[align=center]{Project, PM,\\triple constraint} }
      child { node[align=center]{Phases,\\stakeholders} } }
  child[concept color=popgreen] { node[align=center]{Planning \&\\WBS}
      child { node[align=center]{Scope,\\deliverables} }
      child { node[align=center]{Decomposition\\rules} } }
  child[concept color=poppurple] { node {Scheduling}
      child { node[align=center]{Gantt\\chart} }
      child { node[align=center]{PERT\\network} } }
  child[concept color=popdark] { node[align=center]{Critical Path\\\& Slack}
      child { node[align=center]{Forward /\\backward pass} }
      child { node[align=center]{Total \&\\free slack} } }
  child[concept color=popgold] { node {Control}
      child { node[align=center]{Monitoring,\\KPIs} }
      child { node[align=center]{Change \&\\risk mgmt} } }
  child[concept color=poppink] { node {Integration}
      child { node[align=center]{Full sample\\project} } };
\end{tikzpicture}
\legende{Mind map of Chapter ICT20: Applying project management fundamentals.}
\end{popfigure}

\begin{bilan}
\textbf{Key definitions}
\begin{itemize}
\item A \cle{project} is a \textbf{temporary} endeavour undertaken to create a \textbf{unique} product, service or result, with a defined start and end, limited resources and a specific objective (e.g.\ deploying a school management system in a college in Yaound\'e).
\item \cle{Project management} is the application of knowledge, skills, tools and techniques to project activities to meet the project requirements.
\item The \cle{triple constraint} (iron triangle): \textbf{scope -- time -- cost}; quality sits at the centre. Changing one side affects the others.
\item A \cle{stakeholder} is any person or organisation affected by the project (sponsor, project manager, team, users, suppliers, e.g.\ a ministry or a company such as MTN Cameroon).
\item A \cle{deliverable} is a measurable, verifiable output of a phase or task; a \cle{milestone} is a significant event with zero duration.
\item The \cle{Work Breakdown Structure} (WBS) is a deliverable-oriented hierarchical decomposition of all the work in the project; the lowest level is the \cle{work package}.
\item A \cle{Gantt chart} is a bar chart showing tasks against a calendar; a \cle{PERT chart} is a network diagram (nodes and arrows) showing dependencies between tasks.
\item The \cle{critical path} is the longest path through the network; it gives the minimum project duration. Its tasks have \textbf{zero total slack}.
\item \cle{Total slack} (float): delay a task can absorb without delaying project end. \cle{Free slack}: delay without delaying the earliest start of any successor.
\end{itemize}

\textbf{Formulas and results to know}
\begin{itemize}
\item Project life cycle: \textbf{Initiation $\rightarrow$ Planning $\rightarrow$ Execution $\rightarrow$ Monitoring/Control $\rightarrow$ Closure}.
\item WBS ``100\% rule'': the WBS must capture \emph{all} the work; a work package is typically decomposable until it can be estimated, assigned and controlled (often 8--80 hours of effort).
\item Forward pass: $ES = \max(EF \text{ of predecessors})$, \quad $EF = ES + \text{duration}$.
\item Backward pass: $LF = \min(LS \text{ of successors})$, \quad $LS = LF - \text{duration}$.
\item Total slack: $TF = LS - ES = LF - EF$; \quad Free slack: $FF = \min(ES_{\text{successors}}) - EF$.
\item Project duration $=$ length of the critical path $=$ $\max(EF)$ over terminal tasks.
\item PERT expected duration (if three estimates given): $t_e = \dfrac{t_o + 4t_m + t_p}{6}$.
\end{itemize}

\textbf{Methods to master}
\begin{itemize}
\item \textbf{Building a WBS}: (1) identify the final deliverable; (2) split into major deliverables/phases; (3) decompose into work packages; (4) check the 100\% rule; (5) code each element (1, 1.1, 1.1.1, \ldots).
\item \textbf{Drawing a Gantt chart}: list tasks with durations and dependencies, place a time scale, draw bars, mark milestones, show the critical path in a different colour.
\item \textbf{Critical path analysis}: (1) list tasks, durations, predecessors; (2) draw the network; (3) forward pass for $ES/EF$; (4) backward pass for $LS/LF$; (5) compute slack; (6) tasks with $TF=0$ form the critical path.
\item \textbf{Project control}: compare planned vs actual (schedule, cost, quality), hold review meetings, manage risks (identify, assess probability/impact, plan responses) and process changes through formal \emph{change control}.
\end{itemize}

\textbf{Common mistakes to avoid}
\begin{itemize}
\item Confusing a project with routine \emph{operations} (repetitive, ongoing work is not a project).
\item Putting activities (verbs) instead of deliverables (nouns) in the WBS, or forgetting work so the 100\% rule is violated.
\item Believing the critical path is the \emph{shortest} path --- it is the \textbf{longest}.
\item Computing slack with $EF-ES$ (that is the duration!) instead of $LS-ES$.
\item Forgetting that delaying a critical task delays the whole project, while a non-critical task may slip within its slack.
\item Ignoring stakeholders and risks during planning; no baseline means no control is possible.
\end{itemize}

\textbf{GCE tips}
\begin{itemize}
\item In scheduling questions, always present a neat table: task, duration, predecessor, $ES$, $EF$, $LS$, $LF$, slack --- examiners reward the method, not only the answer.
\item State the critical path \emph{and} the minimum project duration explicitly, e.g.\ ``Critical path: A--C--E--G; duration $= 18$ days''.
\item For the integration lesson, be ready to take a mini-scenario (e.g.\ a school website or a mobile-money awareness campaign), produce its WBS, schedule it, find the critical path and propose control measures --- justify each step.
\item Define every term you use (project, milestone, slack\ldots): precise definitions earn easy marks.
\end{itemize}
\end{bilan}

\findecours

\end{document}
